\documentclass[JEP,XML,SOM,Unicode, NoFloatCountersInSection]{cedram}
\datereceived{2018-09-18}
\dateaccepted{2019-10-15}
\dateepreuves{2019-10-30}

\usepackage{mathrsfs}
\let\mathcal\mathscr
\multlinegap0pt

\newenvironment{enumeratei}
{\bgroup\def\theenumi{\roman{enumi}}\def\theenumii{\Alph{enumii}}\begin{enumerate}}
{\end{enumerate}\egroup}
\newcommand{\RedefinitSymbole}[1]{%
\expandafter\let\csname old\string#1\endcsname=#1
\let#1=\relax
\newcommand{#1}{\csname old\string#1\endcsname\,}%
}
\RedefinitSymbole{\forall} \RedefinitSymbole{\exists}
\newcommand{\sep}{\,:\,}
\newcommand{\spfrac}[2]{\sfrac{#1}{(#2)}}
\newcommand{\psfrac}[2]{\sfrac{(#1)}{#2}}
\newcommand{\Psfrac}[2]{(\sfrac{#1}{#2})}
\newcommand{\sbullet}{{\scriptscriptstyle\bullet}}
\newcommand{\csbullet}{{\raisebox{1pt}{$\sbullet$}}}

\usepackage{stmaryrd}

\numberwithin{equation}{section}

\theoremstyle{plain}
\newtheorem{theorema}{Theorem}
\renewcommand*{\thetheorema}{\Alph{theorema}}

\newtheorem{theorem}{Theorem}[section]
\newtheorem{lemma}[theorem]{Lemma}
\newtheorem{proposition}[theorem]{Proposition}
\newtheorem{cor}[theorem]{Corollary}

\theoremstyle{definition}
\newtheorem{rem}[theorem]{Remark}
\newtheorem{definition}[theorem]{Definition}

\newcommand{\ext}{\mathrm{ext}}
\newcommand{\sm}{\mathrm{small}}
\newcommand{\lint}{\llbracket}
\newcommand{\rint}{\rrbracket}

\newcommand{\bg}{\mathbf{g}}

\newcommand{\Diam}{\mathrm{Diam}}
\newcommand{\ind}{\mathbf{1}}
\newcommand{\restrict}{\!\!\upharpoonright}
\DeclareMathOperator{\Supp}{Supp}
\newcommand{\larg}{\mathrm{large}}
\newcommand{\cc}{\complement}

\newcommand{\cQ}{{\ensuremath{\mathcal Q}} }
\newcommand{\cA}{{\ensuremath{\mathcal A}} }
\newcommand{\cB}{{\ensuremath{\mathcal B}} }
\newcommand{\cP}{{\ensuremath{\mathcal P}} }
\newcommand{\cH}{{\ensuremath{\mathcal H}} }
\newcommand{\cC}{{\ensuremath{\mathcal C}} }
\newcommand{\cL}{{\ensuremath{\mathcal L}} }
\newcommand{\cR}{{\ensuremath{\mathcal R}} }
\newcommand{\cZ}{{\ensuremath{\mathcal Z}} }
\newcommand{\cK}{{\ensuremath{\mathcal K}} }
\newcommand{\bP}{{\ensuremath{\mathbf P}} }
\newcommand{\bQ}{{\ensuremath{\mathbf Q}} }
\newcommand{\bE}{{\ensuremath{\mathbf E}} }
\newcommand{\bC}{{\ensuremath{\mathbf C}} }
\newcommand{\bL}{{\ensuremath{\mathbf L}} }
\newcommand{\bB}{{\ensuremath{\mathbf B}} }

\newcommand{\sumtwo}[2]{\sum_{\substack{#1 \\ #2}}}
\newcommand{\limtwo}[2]{\lim_{\substack{#1 \\ #2}}}

\newcommand{\bbC}{{\ensuremath{\mathbb C}} }
\newcommand{\bbN}{{\ensuremath{\mathbb N}} }
\newcommand{\bbP}{{\ensuremath{\mathbb P}} }
\newcommand{\bbR}{{\ensuremath{\mathbb R}} }
\newcommand{\bbZ}{{\ensuremath{\mathbb Z}} }

\newcommand{\gb}{\beta}
\newcommand{\gga}{\gamma}
\newcommand{\gep}{\varepsilon}
\newcommand{\gG}{\Gamma}
\newcommand{\go}{\omega}
\newcommand{\gO}{\Omega}
\newcommand{\gl}{\lambda}
\newcommand{\gL}{\Lambda}
\newcommand{\tr}{\mathrm{tr}}

\newcommand{\ldom}{\preccurlyeq}
\newcommand{\tf}{\textsc{f}}

\begin{document}


\title[Wetting and layering for SOS II]{Wetting and layering for Solid-on-Solid II: Layering transitions, Gibbs states, and regularity
of the free energy}

\author[\initial{H.} \lastname{Lacoin}]{\firstname{Hubert} \lastname{Lacoin}}
\address{IMPA, Institudo de Matemática Pura e Aplicada\\
Estrada Dona Castorina 110, Rio de Janeiro, CEP-22460-320, Brasil}
\email{lacoin@impa.br}
\urladdr{http://w3.impa.br/~lacoin/}

\thanks{The author acknowledges the support of a productivity grand from CNPq as well as a JCNE grant from FAPERJ}

\begin{abstract}
We consider the Solid-on-Solid model interacting with a wall, which is the statistical mechanics model associated with the integer-valued field $(\phi(x))_{x\in \mathbb{Z}^2}$, and the energy functional
\[
V(\phi)=\beta \sum_{x\sim y}|\phi(x)-\phi(y)|-\sum_{x}\left(h\mathbf{1}_{\{\phi(x)=0\}}-\infty\mathbf{1}_{\{\phi(x)<0\}} \right).
\]
We prove that for $\beta$ sufficiently large, there exists a decreasing sequence $(h^*_n(\beta))_{n\ge 0}$, satisfying $\lim_{n\to\infty}h^*_n(\beta)=h_w(\beta)$, and such that: $(A)$ The free energy associated with the system is infinitely differentiable on $\mathbb{R} \setminus \left(\{h^*_n\}_{n\ge 1}\cup h_w(\beta)\right)$, and not differentiable on $\{h^*_n\}_{n\ge 1}$. $(B)$~For each $n\ge 0$ within the interval $(h^*_{n+1},h^*_n)$ (with the convention $h^*_0=\infty$), there exists a unique translation invariant Gibbs state which is localized around height $n$, while at a point of non-differentiability, at least two ergodic Gibbs states coexist. The respective typical heights of these two Gibbs states are $n-1$ and $n$. The value $h^*_n$ corresponds thus to a first order layering transition from level $n$ to level $n-1$. These results combined with those obtained in \cite{cf:part1} provide a complete description of the wetting and layering transition for SOS.
\end{abstract}

\subjclass{60K35, 60K37, 82B27, 82B44}

\keywords{Solid-on-Solid, wetting, layering transitions, Gibbs states}

\alttitle{Mouillage et stratification pour le modèle SOS II: transitions de niveau, états de Gibbs et régularité de l'énergie libre}

\begin{altabstract}
Nous considérons le modèle \og Solid-On-Solid\fg (SOS) incluant une interaction avec une paroi. Il s'agit du modèle de mécanique statistique associé au champ à valeurs entières $(\phi(x))_{x\in \mathbb{Z}^2}$ et à~la fonctionnelle d'énergie
\[
V(\phi)=\beta \sum_{x\sim y}|\phi(x)-\phi(y)|-\sum_{x}\left(h\mathbf{1}_{\{\phi(x)=0\}}-\infty\mathbf{1}_{\{\phi(x)<0\}} \right).
\]
Nous démontrons que pour des valeurs de $\beta$ suffisamment grandes, il existe une suite décroissante $(h^*_n(\beta))_{n\ge 0}$, satisfaisant $\lim_{n\to\infty}h^*_n(\beta)=h_w(\beta)$, et telle que: $(A)$ l'énergie libre associée au système est infiniment dérivable sur $\mathbb{R} \setminus \left(\{h^*_n\}_{n\ge 1}\cup h_w(\beta)\right)$, et n'admet pas de dérivée aux points $\{h^*_n\}_{n\ge 1}$; $(B)$ pour tout entier $n\ge 0$, pour les valeurs de $h$ dans l'intervalle $(h^*_{n+1},h^*_n)$ (avec la convention $h^*_0=\infty$), il existe une unique mesure de Gibbs correspondant à une hauteur de localisation $n$, alors qu'aux points de non-dérivabilité il y a multiplicité des états de Gibbs, en particulier il en existe deux correspondant aux hauteurs de localisation $n-1$ et~$n$ respectivement. La valeur $h^*_n$ marque donc une transition de niveau entre la hauteur $n$ et la hauteur $n-1$. Ces résultats et ceux prouvés dans {\selectlanguage{english}\cite{cf:part1}} fournissent une description complète des transitions de niveau et de la transition de mouillage pour le modèle SOS.
\end{altabstract}

\altkeywords{Modèle SOS, mouillage, transitions de niveau, état de Gibbs}

\maketitle
\tableofcontents


\section{Introduction}

The Solid-On-Solid model (SOS) introduced in \cite{cf:BCF, cf:T} provides a simplified framework
to study the behavior of two dimensional interfaces in three dimensional systems which display phase coexistence,
such as the Ising model with mixed boundary condition \cite{cf:Dob}.
The SOS interfaces have a simpler description than the ones that appear in most three dimensional lattice models: they are graphs of functions from a subset of~$\bbZ^2$ to~$\bbZ$ and thus have the simplest possible topological structure. The Gibbs weight associated with each possible interface realization also
have a simple expression.
This makes the SOS model considerably easier to analyze than, say, Ising interfaces. On~the other hand, as the simplification performed to obtain
the SOS description starting from a lattice model with phase coexistence, such as the low temperature Ising or Potts model, are not too drastic,
it is believed that results obtained for Solid-On-Solid model may have a predictive value for a large class of interfaces \cite{cf:BCF,cf:T,cf:WGL}.
For this reason, a particular attention has been given to results obtained for SOS concerning the transition from
to rigid interfaces at low temperature to rough ones at high temperature \cite{cf:BM, cf:FS, cf:Swend}, and to the study of layering and wetting transitions
in presence of an interaction with a wall (in a wetting or pre-wetting setup) \cite{cf:ADM, cf:AY, cf:BMF,cf:CM, cf:CM2, cf:chal,cf:DM} (see also \cite{cf:Bas} for a non-rigorous derivation of the layering transitions for the Ising model).
We~refer to the recent review \cite{cf:IV} for a richer introduction to effective interface models as well to the introduction of \cite{cf:part1}
for additional motivation and references.

Our objective is to give a full description of the transitions occurring for the wetting problem,
when an interface interacts with a solid wall which occupies a full half-space.
The problem has been investigated in \cite{cf:chal} where it was shown that the \textit{wetting transition}
occurs for a positive value $h_w(\gb)$
of the intensity of the interaction with the wall: When $h>h_w(\gb)$ the interface is typically localized in a neighborhood of the wall, while for $h<h_w(\gb)$
is repulsed away from it.
In \cite{cf:AY}, a heuristic analysis of the interface stability yielded the prediction that besides this wetting transition, the system should undergo
countably many \textit{layering transitions} which correspond to discrete change of the typical height of the interface.
This analysis also provided a low temperature expansion for the value of first layering critical points.
The first rigorous results concerning this conjectured layering phenomenon were obtained in \cite{cf:ADM}
(results were obtained earlier for the related and more tractable pre-wetting problem, see for instance \cite{cf:CM, cf:DM}): For any given $n\ge 0$, the existence of a
regime where interfaces are localized at height $n$ was evidenced, analyticity of the free energy and results concerning uniqueness of Gibbs states in that regime
were also proved.
The results in~\cite{cf:ADM} nonetheless leave some challenging questions open:
\begin{itemize}
\item [(A)] The existence of a regime with localization at height $n$ is only proved under the assumption that $\gb\ge c \log n$ (with our notation) for some constant $c>0$
(\cf \cite[Rem.\,(4) p.\,528]{cf:ADM}). This limitation obstructs the understanding of how the layering transitions accumulate on the right of $h_w(\gb)$ when $n$ tends to infinity.
\item [(B)] The layering transitions corresponding to the changes of typical height, say from $n$ to $n-1$ cannot be analyzed,
as the intervals on which localization is
shown to occur are not adjacent. This is because the perturbative approach used in \cite{cf:ADM} does not enable to come close to the layering critical points.
\end{itemize}
The present paper overcomes these limitations and proves that for $\gb$ sufficiently large
the free energy is infinitely differentiable everywhere except on a countable set which
corresponds to the layering critical points. On this set, the first derivative of the free energy is shown to be discontinuous.
The existence of Gibbs states is also proved when the asymptotic contact fraction with the substrate is positive (non-existence is proved in the other case), together with uniqueness on intervals
where the free energy is differentiable, and non-uniqueness at points of non differentiability.
Combined with the results obtained in \cite{cf:part1} concerning the value critical point $h_w(\gb)$
and the sharp asymptotics for the free energy, this yields a complete picture of the system's behavior.

\subsubsection*{Acknowledgements}
We are grateful to Kenneth Alexander for enlightening discussions.
This work has been performed in part
during a visit at the Institut Henri Poincaré (2017, spring-summer trimester) supported by the Fondation de Sciences
Mathématiques de Paris. The author thanks the Institute for its hospitality.

\section{Model and results}

\subsection{The Solid on Solid Model on $\bbZ^2$}

Consider $\gL$ a finite subset of $\bbZ^2$ (equipped with its usual lattice structure) and let $\partial \gL$ denote its external boundary
\[
\partial \gL:= \{ x \in \bbZ^2 \setminus \gL \sep \exists y \in \gL, \ x \sim y \}.
\]
Given $\psi\in \bbZ^{\bbZ^2}$, we define the Hamiltonian for SOS in the domain $\gL$ with boundary condition $\psi$ on the set $\gO_{\gL}:= \bbZ^\gL$ by
\begin{equation}\label{defhamil}
\cH^{\psi}_\gL(\phi):= \frac{1}{2}\sumtwo{x, y \in \gL}{ x \sim y} |\phi(x)-\phi(y)|
+\sumtwo{x\in \gL, y\in \partial\gL}{x\sim y} |\phi(x)-\psi(y)|, \quad \forall \phi \in \gO_{\gL}.
\end{equation}
Given $\gb>0$, we define
the SOS measure with boundary condition $\psi$, $\bP^{\psi}_{\gL,\gb}$ on $\gO_{\gL}$ by
\begin{equation*}%\label{defSOS}
\bP^{\psi}_{\gL,\gb}(\phi):= \frac{1}{\mathcal Z^{\psi}_{\gL,\gb}}e^{-\gb\cH^{\psi}_\gL(\phi)} \quad \text{ where }
\mathcal Z^{\psi}_{\gL,\gb}:=\sum_{\phi\in \gO_{\gL}} e^{-\gb\cH^{\psi}_\gL(\phi)}.
\end{equation*}
For most purposes, we only have to consider the constant boundary conditions $\psi \equiv n$ for $n\ge 0$.
In that case we simply write $\bP^{n}_{\gL,\gb}$ and $\cZ^{n}_{\gL,\gb}$.
We drop the superscript~$n$ in the notation in the special case $n=0$.
Note that by translation invariance $\mathcal Z^n_{\gL,\gb}=\mathcal Z_{\gL,\gb}$ does not depend on $n$.
We also define the \textit{free energy} (sometimes also referred to as \textit{pressure}) for the SOS model by
\begin{equation*}
\tf(\gb):=\limtwo{|\gL|\to \infty}{|\partial\gL|/|\gL|\to 0} \frac{1}{|\gL|}\log \cZ_{\gL,\gb},
\end{equation*}
where the limit can be taken over any sequence of finite sets $(\gL_N)_{N\ge 0}$ such such ratio between the cardinality of
$\gL_N$ and that of its boundary vanishes. A justification of the existence of the limit is given in the introduction of \cite{cf:part1}.
We used $|\cdot|$ to denote the cardinality of a set. In the remainder of the paper we also use the same notation to denote the $\ell_1$ distance on the lattice but this should not yield confusion.

When $\gb$ is sufficiently large, it is known \cite[Th.\,2]{cf:BM} that $\bP_{\gL,\gb}$ converges (in the sense of finite dimensional marginal)
to an infinite volume measure $\bP_{\gb}$ or \textit{Gibbs state} (see Definition \ref{defgs}). We introduce a quantitative version of the statement which requires
the introduction of some classic terminology.

We say that a function $f: \ \gO_{\infty}:=(\bbZ)^{\bbZ^2}\to \bbR$ is local if there exists $(x_1,\dots,x_k)$ and $\tilde f: (\bbZ)^{k}\to \bbR$
such that $f(\phi)=\tilde f(\phi(x_1),\dots,\phi(x_k))$. The minimal choice (for the inclusion)
for the set of indices $\{x_1,\dots, x_k\}$ is called the support of $f$ ($\Supp(f)$).
With some abuse of notation, whenever $\gL$ contains the support of $f$, we extend $f$ to $\gO_{\gL}$ in the obvious way.
An event is called local if its indicator function is a local function.

Given $\bP_{\gL}$ a sequence of measures on $\gO_{\gL}$ and $\bP$ a measure on $\gO_{\infty}$.
We say that~$\bP_{\gL}$ \textit{converges locally} to $\bP$ when $\gL$ exhausts $\bbZ^2$ if for any sequence $\gL_N$ exhausting $\bbZ^2$,
and any local function $f$ we have
\[
\lim_{N\to \infty} \bP_{\gL_N}[f(\phi)]=\bP[f(\phi)].
\]

For $A$ and $B$ two finite subsets of $\bbZ^2$ we set
\begin{equation*}%\label{disset}
d(A,B):=\min_{x\in A, y\in B} |x-y|,
\end{equation*}
where $|\cdot|$ denote the $\ell_1$ distance. The proofs in \cite{cf:BM} imply that $\bP_{\gL,\gb}$
converges exponentially fast in some sense to some measure on $\gO_{\infty}$. The statement below can also be proved using the
techniques introduced in Section \ref{clusexp} (see Remark \ref{teorema}).

\begin{theorema}\label{infinitevol}
There exists $\gb_0>1$ and $c$ such that
for any $\gb>\gb_0$,
there exists a measure $\bP_{\gb}$ defined on $\gO_{\infty}$ such that for every local function
$f: \gO_{\infty}\to [0,1]$ with $\Supp(f)=A$, and every $\gL$ which contains $A$,
\begin{equation*}%\label{decayz}
\bigl| \bE_{\gL,\gb}[f(\phi)]-\bE_{\gb}[f(\phi)] \bigr|\le C_{\gb} |A| e^{-c \gb d(\partial \gL, A)}.
\end{equation*}
\end{theorema}

\subsection{The wetting problem for the SOS model}

For $\phi \in \gO_{\gL}$ and $A\subset \bbZ$ (or $\bbR$), we set
\[
\phi^{-1}(A):=\{ x\in \gL \sep \phi(x)\in A\}.
\]
We write also $\phi^{-1}A$ when more convenient.
Using the notation $\bbZ_+:=\bbZ\cap[0,\infty)$, we define
\begin{equation}\label{positivity}
\gO^+_{\gL}:=(\bbZ_+)^{\gL}=\{\phi \in \gO_{\gL} \sep x\in \gL,\ \phi(x)\ge 0\}.
\end{equation}
This convention of adding the superscript $+$ to indicate a restriction to the set of positive functions
is used in other contexts throughout the paper.

Given $h\in \bbR$ we consider $\bP^{\psi,h}_{\gL,\gb}$ which is a modification of $\bP^{\psi}_{\gL,\gb}$
where the interface $\phi$ is constrained to remain positive
and gets an energetic reward $h$ for each contact with $0$.
It is defined as follows
\begin{equation*}%\label{def}
\bP^{\psi,h}_{\gL,\gb}(\phi):= \frac{1}{\mathcal Z^{\psi,h}_{\gL,\gb}}e^{-\gb\cH^\psi_\gL(\phi)+ h|\phi^{-1}\{0\}|},
\quad \text{where } \mathcal Z^{\psi,h}_{\gL,\gb}:=\sum_{\phi\in \gO^+_{\gL}} e^{-\gb\cH^\psi_\gL(\phi)+h|\phi^{-1}\{0\}|}.
\end{equation*}
In this case also, we replace $\psi$ by $n$ in the notation for the special case $\psi\equiv n$.
The aim of our study is to investigate the localization transition in $h$ for $\bP^{n,h}_{\gL,\gb}$ which appears in the limit when $\gL$ exhausts $\bbZ^2$.
A key quantity to study the phenomenon is the corresponding free energy
\begin{equation*}%\label{nodepends}
\tf(\gb,h):=\limtwo{|\gL|\to \infty}{|\partial \gL|/|\gL|\to 0} \frac{1}{|\gL|}\log \mathcal Z^{n,h}_{\gL,\gb}.
\end{equation*}
The reader can check that as a consequence of the inequality
\[
|\cH^n_\gL(\phi)- \cH^m_\gL(\phi)|\le 4 |m-n| |\partial \gL|
\]
the quantity
$\tf(\gb,h)$ indeed does not depend on $n$.

To clarify notation, in the remainder of the paper, we often consider the limit along the sequence
$\gL_N:=\lint 1, N\rint^2$ (using the notation $\lint a,b \rint=[a,b]\cap \bbZ$).
We write $\cZ_{N,\gb}$ for $\cZ_{\gL_N,\gb}$ and adopt a similar convention for other quantities.

The function $h\mapsto \tf(\gb,h)$ is non-decreasing and convex in $h$ (as a limit of non-decreasing convex function).
At points where $\tf(\gb,h)$ is differentiable, convexity makes it possible to exchange the positions of limit and derivative
thus $\partial_h \tf(\gb,h)$ corresponds to the
asymptotic contact fraction. Thus for every $n\in \bbN$ we have
\begin{equation}\label{difencial}
\partial_h \tf(\gb,h)=\lim_{N\to \infty} \frac{1}{N^2} \bE^{n,h}_{N,\gb}[|\phi^{-1}(0)|],
\end{equation}
wherever $\partial_h\tf$ is defined (by convexity this is everywhere except possibly on a countable set).

We define
$h_w(\gb)$ to be the value of $h$ which marks the wetting transition between a localized phase (where the asymptotic contact fraction is positive)
and a delocalized phase
\begin{equation*}
\begin{split}
h_w(\gb)&:= \inf\{ h\in \bbR \sep \partial_h \tf(\gb,h) \text{ exists and is positive}\}\\
&\phantom{:}= \sup\{ h\in \bbR: \tf(\gb,h)= \tf(\gb)\}.
\end{split}
\end{equation*}

\subsection{The asymptotic behavior for the free energy}

In previous work \cite{cf:part1}, we established the value of $h_w(\gb)$ answering a question left open since the pioneering work of Chalker \cite{cf:chal}, and
we were able to describe the asymptotic behavior of $\tf(\gb,h)$ close to the critical point. To state this result we need to introduce a few quantities.
Letting ${\bf 0}$ and ${\bf 1}$ denote the vertices $(0,0)$ and $(1,0)$ respectively, we define, for $\gb>\gb_0$
\begin{equation}\label{singledouble}
\begin{split}
\alpha_1(\gb)&:= \lim_{n\to \infty} e^{4\gb} \bP_{\gb} \left[ \phi({\bf 0})\ge n\right],\\
\alpha_2(\gb)&:= \lim_{n\to \infty} e^{6\gb} \bP_{\gb} \left[ \min(\phi({\bf 0}), \phi({\bf 1})) \ge n\right].
\end{split}
\end{equation}
For a proof of the existence of these quantity, we refer to \cite[Prop.\,4.6]{cf:part1}.
We also set $J:=e^{-2\gb}$ and for $u\in \bbR$
\begin{equation*}
\bar \tf(\gb,u)=\tf\Bigl(\gb,\log \Bigl(\frac {e^{4\gb}}{e^{4\gb}-1}\Bigr)+u \Bigr)-\tf(\gb).
\end{equation*}

\begin{theorema}\label{oldmain}
If $\gb>\gb_0$ (of Theorem \ref{infinitevol}), we have
\[
h_w(\gb)=\log\Bigl(\frac{e^{4\gb}}{e^{4\gb}-1}\Bigr).
\]
Furthermore
\begin{equation}\label{lequivdelamor}
\bar \tf(\gb,u)\stackrel{u\to 0+}{\sim} F(\gb,u),
\end{equation}
where
\begin{equation}\label{uds}
F(\gb,u):=\max_{n\in \bbZ_{+}}\Bigl(\alpha_1 J^{2n}u- \frac{2 \alpha_2(J^3-J^4)}{1-J^3} J^{3n} \Bigr).
\end{equation}
\end{theorema}

A detailed heuristic justification for the expression \eqref{uds} is given in \cite[\S 3]{cf:part1}. It can be roughly summarized as follows: one can rewrite the partition function of the model in a way that suppresses the positivity constraint $\{ \forall x \in \gL_N \sep \phi(x)\ge 0\}$ at the cost of adding energetic penalties
in sites where $\phi(x)\le 0$, which
exactly account for the entailed entropic loss.

With this rewriting when $h=\log\left(\spfrac{e^{4\gb}}{e^{4\gb}-1}\right)+u$, each spike in the lower half space of width $1$ gets a positive retribution $u$ in the exponential Boltzmann weight, while spikes of width $2$ and larger are penalized (the size of spikes is measured at level~$0$, meaning that we are looking at connected components of $\phi^{-1}(0)$).
In particular, spikes of size two yield a penalty equal to $\spfrac{ (J^3-J^4)}{1-J^3}$.
For this reason for small $u$, the best localization strategy is choose a high boundary condition $n$ so that the density of spikes of size one (which is of order $e^{-4\beta n}$ by a Peierls type argument) is much larger than that of larger spikes (which is of order $e^{-6\beta n}$).
The expression~\eqref{uds} comes from the fact that in the small $u$ limit, only the contribution of spikes of size one and two become relevant and their relative density is well approximated by \eqref{singledouble}. The free-energy is obtained by choosing $n$ that optimizes the balance between the gain of single spikes and the cost of double spikes.
Note that the function $F(\gb,u)$ is piecewise affine on $\bbR_+$, and presents angular points at
\begin{equation}\label{defun}
u=u_n:=\frac{2\alpha_2}{\alpha_1(1+J)}J^{n+2},
\end{equation}
for $n\ge 1$.
While Theorem \ref{oldmain} does not imply the convergence of $\partial_u \bar \tf(\gb,u)$ and thus, the presence
of angular points on the free energy curves, convexity implies that for large values of $n$ the contact fraction changes abruptly
around $u_n$.

An explanation for this phenomenon which is corroborated by the above heuristic (again we refer to \cite{cf:part1} for a deeper insight) is that the typical behavior of $\phi$ changes radically
around $u_n$: when $u\le u_n+o(J^{n})$ the surface $\phi$ tends to localize at height $n$, meaning that $\phi(x)=n$ for a majority of points and that
connected components of the set $\{ x \sep \phi(x)\ne n\}$ are all of small diameters,
while when $u\ge u_n+o(J^{n})$ the typical height should be $n-1$.

This indicates that there should exist a value $u^*_n$ which delimits a phase transition between these two kinds of behavior. Moreover it should satisfy, asymptotically
for large values of $n$,
$u^*_n=u_n+o(J^{n})$.
The change of behavior around $u^*_n$ should provoke a discontinuity in the contact fraction,
so that these sequences of phase transition should be manifested by discontinuities
for $\partial_u \bar \tf(\gb,u)$. This prediction can be interpreted as a refined version
of the conjecture presented in \cite[Stat.\,p.\,228]{cf:ADM}. Earlier versions of the same conjecture is found in \cite{cf:AY} and \cite[\S 4.3]{cf:BMF}.

\subsection{Main results}

In the present paper, we bring the above stated conjecture on a rigorous ground by showing the existence of an infinite sequence of point of discontinuity for
$\partial_u \bar \tf(\gb,u)$. We complement this result by relevant information concerning the regularity of $\tf(\gb,u)$ between these transition points,
and a statement concerning uniqueness of Gibbs states.
These results are summarized in the caption of Figure \ref{fig:freen}.

\begin{theorem}\label{main}
For $\gb\ge \gb_1$ sufficiently large,
there exists a decreasing sequence $(u^*_n(\gb))_{n\ge 1}$, which satisfies
\[
\frac{1}{200} J^{n+2} \le u^*_n \le 200 J^{n+2}
\]
and
\begin{equation}\label{asimptic}
\lim_{n\to \infty} e^{2\gb n} u^*_n=\frac{2\alpha_2J^{2}}{\alpha_1(1+J)}.
\end{equation}
which is such that
\begin{enumeratei}
\item
the function $u\mapsto \bar \tf(\gb,u)$ is infinitely differentiable on $(u^*_{n+1}(\gb),u^*_n(\gb))$ for any $n\ge 1$ and also on $(u^*_{1}(\gb),\infty)$,
\item
for any $n\ge 1$, $\bar \tf(\gb,u)$ is not differentiable at $u^*_n(\gb)$, meaning that that the left and right derivative at $u^*_n$ do not coincide
\begin{equation*}
\partial^-_u \bar \tf(\gb,u^*_n)<\partial^+_u \bar \tf(\gb,u^*_n).
\end{equation*}
\end{enumeratei}
\end{theorem}

\begin{rem}
We believe that the free energy is in fact analytic in $u$ on the domain where it is differentiable. While such a statement could in principle be directly deduced from the convergence of the cluster expansion,
it would require to be able to obtain a convergence result for complex values of $u$, more precisely for each $n$ one should prove convergence of the expansion
on an open subset of $\bbC$ which contains the real interval $(u^*_{n+1},u^*_{n})$.
\end{rem}

\begin{rem}
Note that the main result in \cite{cf:ADM} includes a statement about analyticity. While this is not explicitly stated in the proofs, it appears that
the cluster expansion considered in \cite{cf:ADM} also converges when the parameter $u$ considered in~\hbox{\cite[Eq.\,(1.15)]{cf:ADM}} is allowed to have a small imaginary part \cite{cf:Al}.
It seems plausible that with some (significant) efforts, our proof of Proposition \ref{truta} could be adapted to handle also small imaginary perturbation
of $u$, which would yield analyticity of $\bar \tf(\gb,u)$ in the intervals of the type $[200 J^{n+2}, \frac{1}{200} J^{n+1}]$, $n\ge 1$ and on $[200 J^2,\infty)$.
However
monotonicity plays a too central role in Proposition \ref{mejor} to extend this kind of argument.
For this reason the proof of analyticity in the neighborhood of $u^*_n$ appears like a more challenging task.
\end{rem}

To state our second result about convergence for the measure $\bP^{n,h}_{\gL,\gb}$ we need to recall some terminology.

\begin{definition}\label{defgs}
An infinite volume measure or \textit{Gibbs state} for parameter $(\gb,h)$ is a measure $\nu_{\gb,h}$ on $(\bbZ_+)^{\bbZ^2}$
such that for any finite $\gL\subset \bbZ^2$, we have for $\nu_{\gb,h}$-almost all~$\psi$,
\begin{equation}\label{DLR}
\nu_{\gb,h}\left[ \phi\restrict_{\gL}\in \csbullet \ | \ \phi\restrict_{\gL^{\cc}}=\psi\restrict_{\gL^{\cc}} \right]= \bP^{\psi,h}_{\gL,\gb}\left[ \phi \in \csbullet \right].
\end{equation}
\end{definition}

It is not difficult to check that
the relation \eqref{DLR}, often referred to as the Dobrushin-Lanford-Ruelle (DLR) Equation, is valid if one replaces $\nu_{\gb,h}$
by a measure $\bP^{\Psi,h}_{\gb,\gL'}$ defined on a domain $\gL'$ which includes $\gL$ and with arbitrary boundary condition~$\Psi$.
As a consequence, the measures obtained as local limits of $\bP^{\psi,h}_{N,\gb}$ for $\psi\in \bbZ_+^{\bbZ^2}$ are Gibbs state.

For technical purpose we define $\star$-connectivity, as the connectivity associated with the network $\bbZ^2$ were diagonal edges of the type
$\{x,x+(1,1)\}$ have been added (but~not the other diagonals).

\begin{definition}
A Gibbs state $\nu_{\gb,h}$ is said to
\begin{enumeratei}
\item
be \textit{translation invariant} if under $\nu_{\gb,h}$, the distribution of $\phi$ and $\theta_z(\phi):= (\phi(z+\nobreak x))_{x\in \bbZ^2}$
are the same;
\item
have \textit{finite mean} if for all $x\in \bbZ^2$
\begin{equation*}%\label{finitemean}
\nu_{\gb,h}\left[ \phi(x) \right]<\infty;
\end{equation*}
\item
\textit{percolate at level} $n$ if $\nu_{\gb,h}$ almost surely, $\phi^{-1}(n)$ has unique infinite connected component in $\bbZ^2$
and all connected component of $\phi^{-1}[n+1,\infty)$ and $\phi^{-1}(-\infty,n-1]$ are finite.
\end{enumeratei}
\end{definition}

Setting $h^*_n=h_w(\gb)+u^*_n$,
our second result essentially claims that translation invariant Gibbs states are unique for $h\in (h_w(\gb),\infty)\setminus \{h^*_n\}_{n\ge 1}$
and that multiple translation invariant Gibbs states coexist at the layering points $(h^*_n)_{n\ge 1}$.

\begin{theorem}\label{Gibbs}
For $\gb$ sufficiently large, the following holds true.
\begin{enumeratei}
\item
For $h\le h_w(\gb)$, there exists no Gibbs state for $(\gb,h)$.
\item
When $n\ge 0$, $h\in (h^*_{n-1},h^*_{n})$ then there exists a unique finite mean translation invariant Gibbs state
which we call $\bP^{n,h}_{\gb}$.
Moreover $\bP^{n,h}_{\gb}$ percolates at level $n$.

When $h=h^*_n$, $n\ge 1$, then there exist several finite mean translation invariant Gibbs states.
In particular, we can identify two extremal states $\bP^{n-1,h^*_n}_{\gb}$ and $\bP^{n,h^*_n}_{\gb}$ which satisfy:
\begin{enumeratei}
\item\mbox{}\vspace*{-10pt}
\begin{equation}\label{variousgibbs}
\begin{split}
\bP^{n,h^*_n}_{\gb}[\phi(x)=0]&= \partial^-_h\tf(\gb,h^*_n), \\
\bP^{n-1,h^*_n}_{\gb}[\phi(x)=0]&= \partial^+_h\tf(\gb,h^*_n).
\end{split}
\end{equation}
\item
$\bP^{n-1,h^*_n}_{\gb}$ and $\bP^{n,h^*_n}_{\gb}$ respectively percolate at level $n-1$ and $n$.
\item
We have $\bP^{n-1,h^*_n}_{\gb}\preccurlyeq \bP^{n,h^*_n}_{\gb}$. Any other finite mean translation invariant Gibbs state $\nu$
for parameters $(\gb,h^*_n)$ satisfies
\begin{equation}\label{ssanduba}
\bP^{n-1,h^*_n}_{\gb}\preccurlyeq \nu \preccurlyeq \bP^{n,h^*_n}_{\gb}.
\end{equation}
\end{enumeratei}
\end{enumeratei}
\end{theorem}

\begin{rem}
We believe that there are no infinite mean translation invariant Gibbs state for $h>h_w(\gb)$ and that the finite mean assumption is present only for technical reasons.
We would also tend to believe that in analogy with low temperature two dimensional Ising model (see \cite{cf:A, cf:CoV, cf:H} for results and proofs)
$\bP^{n-1,h^*_n}_{\gb}$ and $\bP^{n,h^*_n}_{\gb}$ are in fact the only
ergodic Gibbs states when $h=h^*_n$, but proving such a statement is out of the scope of this paper. While the arguments used for the Ising case in the above references remain valid at the heuristic level, adapting them to to the SOS model represents a significant technical challenge.
\end{rem}

\begin{figure}[ht]
\begin{center}
\hspace*{-5mm}\includegraphics[scale=.85]{figures/freeen}
\end{center}
\vspace*{-.5\baselineskip}
\caption{\label{fig:freen}\smaller
A schematic representation of the free energy curve. The curve is $C^{\infty}$ except at the points of abscissa $h^*_i$ (represented as square dots) where the derivative on the left and on the right differ. The $h^*_i$ approximately forms a geometric sequence which accumulates to the right of $h_w(\beta)$.
The numbers below the $h$-axis corresponds to the height at which the interfaces localizes
for the corresponding value of~$h$, that the height at which interface percolate almost surely under the unique Gibbs state. At the critical values~$h^*_i$ there is no unicity of the Gibbs state and there are two possibility for the localization height. Below $h_w$ there is no Gibbs state and the interface is repulsed to infinite height. }
\end{figure}

Finally we conclude the exposition with a result showing that our Gibbs states exhibit exponential decay of correlation.

\begin{proposition}\label{Gibbset}
For $\gb$ sufficiently large, there exists constants $c$ and $C$ such that
for every $n\ge 0$ and any $h\in [h^*_{n+1},h^*_{n}]$, we have for any pair of local functions $f:\gO_{\infty}\to [0,1]$ and $g:\gO_{\infty}\to [0,1]$ with respective supports
$A$ and $B$ we have
\begin{equation*}
\left|\bE^{n,h}_{\gb}\left[ f(\phi)g(\phi) \right]- \bE^{n,h}_{\gb}\left[ f(\phi)\right]\bE^{n,h}_{\gb}\left[g(\phi) \right] \right|\le C
|A| e^{-c \gb d(A,B)},
\end{equation*}
and thus in particular using the notation $\delta_x:=\ind_{\{\phi(x)=0\}}$.
\begin{equation*}
| \bE^{n,h}_{\gb}\left[ \delta_x\delta_y \right]-\bE^{n,h}_{\gb}\left[ \delta_x\right]\bE^{n,h}_{\gb}\left[ \delta_y \right]|\le C e^{-c \gb |x-y|}.
\end{equation*}

\end{proposition}

\subsection{Possible extension to other surface models}

The proof of our result relies a lot on the contour representation for
two dimensional interfaces which is quite specific to two dimensional Solid on Solid. However there are some other surface models which should exhibit similar phenomenology at low temperature. Let us discuss them shortly.

\subsubsection{Other planar regular lattice}

Solid on solid can also be defined on the triangular and hexagonal lattice. While our approach in this paper relies a lot on planarity, it does not rely on the specifics of the lattice. Thus modulo minor modification (\hbox{\cite[\S 2.4]{cf:part1}}) the proof of Theorem \ref{main} should transpose to these cases.

\subsubsection{Higher dimensional Solid on Solid}
The behavior of Low temperature Solid on Solid does not depend at all on the dimension, and thus we expect that a result similar to Theorem \ref{main} to hold (indeed Theorem B is for $d\ge 3$ with appropriate change see \cite[\S 2.4]{cf:part1}).
However our proof does not adapt to this case. One possibility to extend the result would be to use another kind of cluster expansion (similar to the one presented in \cite{cf:BM}, however, much of our argument rely on the specifics of the contour decomposition and adapting it seems to present significant challenges.\enlargethispage{\baselineskip}%

\subsubsection{ Discrete Gaussian models and $|\nabla \phi|^p$ surfaces}

Remaining in dimension two, one can also look at minor modifications of the Hamiltonian. On possibility is to replace $|\phi(x)-\phi(y)|$ in \eqref{defhamil} by $|\phi(x)-\phi(y)|^p$, $p\ge 1$ (see \eg \cite{cf:PFF} and references therein for more on these types of model).
The case $p=2$ corresponds to the \textit{discrete Gaussian model}.

These models should exhibit a simpler wetting transition as soon as $p>1$
(heuristics strongly suggests for instance that $h_c(\beta)=0$) but the layering phenomenon should nonetheless be present.
When $p\in (1,\infty)$, there is no contour decomposition of the model (in the sense that there is no possibility to rewrite the Hamiltonian as a local function of the contours of $\phi$ even in the homogeneous case).
It the case $p=\infty$ however (which corresponds to imposing the constraint that the gradient of the field should take value in $\{-1,0,1\}$) there is hope that the tools developed in the present paper can be applied to prove the presence of countably many first order layering transitions accumulating at the right of $0$.

\subsection{Organization of the proof}

All the results exposed in the previous section are going to be derived as consequences of the
convergence of a cluster expansion associated with a certain contour representation of the partition functions (see the introduction of \cite{cf:BFP} and mentioned references
for a review of cluster expansion techniques).
To prove the convergence, we need to obtain very fine estimates on finite size partition functions: these are obtained by combining various ingredients such as
asymptotic properties of the SOS model (Proposition \ref{rouxrou}),

Therefore our first task is to introduce the necessary framework for the exposition of this result.
This is the purpose of Section \ref{tools}, in which we introduce various technical tools, including contour representations,
cluster expansion methods, and FKG inequalities.

In Section \ref{zorga}, we introduce the main technical result of the paper, Theorem \ref{converteo}, which implies the convergence of a cluster expansion associated with the measures
$\bP^{n,h}_{\gL,\gb}$
and give the main steps of its proof.
We also explain how Theorem \ref{main} can be deduced from this convergence result.

In Section \ref{ponverteo}, which is the technical core of the paper, we perform the proof of Theorem \ref{converteo} in full details.
For better readability, the proof of the more technical estimates presented in Section \ref{ponverteo} are performed separately, in Section \ref{secrestrict}.

In Section \ref{dalast}, we explore the consequences of Theorem \ref{converteo} on the measure $\bP^h_{\gL,\gb}$ and prove in particular Proposition \ref{Gibbset}.
Finally in Section \ref{daverylast} we prove the remaining statements of Theorem \ref{Gibbs}.

\section{Technical preliminaries}\label{tools}

\subsection{Contour representation}\label{contour}

We recall briefly how to describe a function $\phi\in \gO_{\gL}$ using only its level lines.
The formalism of this section is identical to the one used in~\cite{cf:part1}, and close to the one displayed in \eg \cite{cf:ADM, cf:CM, cf:PLMST}.

We let $(\bbZ^2)^*$ denote the dual lattice of $\bbZ^2$ (dual edges cross that of $\bbZ^2$ orthogonally in their midpoints).
Two adjacent edges $(\bbZ^2)^*$ meeting at $x^*$ of are said to be \textit{linked} if they both
lie on the same side of the line making an angle $\pi/4$ with the horizontal
and passing through $x$.
(see Figure \ref{linked}.)

We define a \textit{contour sequence} to be a finite sequence $(e_1,\dots,e_n)$ of distinct edges of $(\bbZ^2)^*$ which satisfies:
\begin{enumeratei}
\item
For any $i=\lint 1,n-1\rint$, $e_i$ and $e_{i+1}$ have a common end point $(\bbZ^2)^*$, $e_1$ and $e_{|\gamma|}$ also have a common end point.
\item
If for $i\ne j$, if $e_i$, $e_{i+1}$, $e_j$ and $e_{j+1}$ meet at a common end point then $e_i$, $e_{i+1}$ are linked and so are $e_j$ and $e_{j+1}$ (with the convention that $n+1=1$).
\end{enumeratei}
A \textit{geometric contour} $\tilde \gamma:=\{e_1,\dots,e_{|\tilde \gamma|}\}$ is a set of edges that forms a contour sequence when displayed in the right order.
The cardinality $|\tilde \gamma|$ of $\tilde\gamma$ is called the length of the contour.
A \textit{signed contour} or simply \textit{contour} $\gamma=(\tilde \gamma,\gep)$
is a pair composed of a geometric contour and a sign $\gep\in \{+1,-1\}$. We let $\gep(\gamma)$ denote the sign associated with a contour $\gamma$,
while with a small abuse of notation,
$\tilde \gamma$ will be used for the geometric contour associated to $\gamma$ when needed.
For $x^*\in (\bbZ^2)^*$ we write $x^*\in \gamma$ or $x^*\in \tilde \gamma$ when the point $x^*$ is visited by one edge of the geometric contour.

\begin{figure}[ht]
\begin{center}
\includegraphics[scale=.17]{figures/split}
\end{center}
\caption{\label{linked}\smaller
The rule for splitting a four edges meeting at one points into two pairs of linked edges.
To obtain the set of contours that separates $\{x \sep \phi(x)\ge h\}$ from $\{x \sep \phi(x)< h\}$ for $h\in \bbZ$,
we draw all dual edges separating two sites $x$, $y$ such that $\phi(x)\ge h>\phi(y)$ and apply the above graphic rule for every dual vertex where four edges meet.
When several sets of level lines include the same contour, it corresponds to a cylinder of intensity $2$ or more for $\phi$.}
\end{figure}

We let $\bar \gamma$ denote the set of vertices of $\bbZ^2$ enclosed by $\tilde \gamma$.
We refer to $\bar \gamma$ as the \emph{interior of} $\gamma$ and say that $|\bar \gamma|$ is the volume enclosed in the contour $\gamma$.
We let $\Delta_{\gamma}$, the neighborhood of $\gamma$, be the set of vertices of $\bbZ^2$ located either at a (Euclidean) distance $1/2$ from
$\tilde\gamma$
(when considered as a subset of $\bbR^2$)
or at a distance $1/\sqrt{2}$ from the meeting point of two non-linked edges.
We split the $\Delta_{\gamma}$ into two disjoint sets, the internal and the external neighborhoods of $\gamma$
(see Figure \ref{compa})
\[
\Delta^-_{\gamma}:=\Delta_{\gamma}\cap \bar \gamma \quad \text{ and } \quad \Delta^+_{\gamma}:=\Delta_{\gamma}\cap \bar \gamma^{\cc}.
\]

Given a finite set $\gL\subset \bbZ^2$
a contour $\gamma$ is said to be in $\gL$ is if $\bar \gamma\subset \gL$.
We let $\cC$ denote the set of contours in $\bbZ^2$ and $\cC_{\gL}$ that of contours in $\gL$.

Given $\phi\in \gO_{\gL}$, we say that $\gamma \in \cC_{\gL}$ is a contour
for $\phi$ with boundary condition $n$, if there exists $k\ge 1$ such that
\begin{equation}\label{cont+}
\min_{x\in \Delta^-_{\gamma}} \phi(x)= \max_{x\in \Delta^+_{\gamma}} \phi(x)+k\gep(\gamma),
\end{equation}
where in the above equation by convention we consider that
\[
\phi(x)=n \quad \text{ if } \quad x\in \gL^{\cc}.
\]

The quantity $k$ appearing in \eqref{cont+} is called the \textit{intensity} of the contour and the
triplet $(\gamma,k)=(\tilde \gamma, \gep(\gamma), k)$
with $\gamma \in \cC$ and $k\in \bbN$ an intensity, is called a \textit{cylinder}. We say that $(\gamma,k)$
is a cylinder for $\phi$ (with boundary condition $n$) if $\gamma$ is a contour of intensity $k$.
The cylinder function associated to $(\gamma,k)$ is defined on $\bbZ^2$ by
\begin{equation}\label{cylfunc}
\varphi_{(\gamma,k)}(x)=\gep(\gamma) k \ind_{\bar \gamma}(x).
\end{equation}
We use $\hat \gamma$ to denote a generic cylinder associated with the contour $\gamma$ (we use the notation $k(\hat \gamma)$ to denote its intensity).
We let $\hat \Upsilon_n(\phi)$ denote the set of cylinders for $\phi$ with boundary condition $n$ and $\Upsilon_n(\phi)$ the associated set of contours.

We say that $\gL$ is a \textit{simply connected} subset of $\bbZ^2$, if it can be expressed as the interior of a contour,
that is, if
\begin{equation} \label{connectedness}
\exists \gamma_{\gL}\in \cC, \quad \bar\gamma_{\gL}=\gL.
\end{equation}
Note that, when $\gL$ is simply connected,
an element
$\phi\in \gO_{\gL}$ is uniquely characterized by its cylinders. More precisely, we have
\begin{equation}\label{cylinder}
\forall x\in \gL, \quad \phi(x):= n+\sum_{\hat \gamma \in \hat \Upsilon_n(\phi)} \varphi_{\hat \gamma}(x).
\end{equation}
Furthermore, the reader can check that
\begin{equation}\label{express}
\cH^n_{\gL}(\phi)=\sum_{\hat \gamma \in \hat \Upsilon_n(\phi)} k(\hat \gamma)\,|\tilde \gamma|.
\end{equation}
Of course not every set of cylinder is of the form $\hat \Upsilon_n(\phi)$ and we must introduce a notion of compatibility
which characterizes the ``right'' sets of cylinder.

Two cylinders $\hat \gamma$ and $\hat \gamma'$ are said to be \emph{compatible} if they are
cylinders for the function $\varphi_{\hat \gamma}+\varphi_{\hat \gamma'}$.
This is equivalent to the three following conditions being satisfied
(see Figure \ref{compa}):
\begin{enumeratei}
\item\label{enum:i}
$\tilde \gamma\ne \tilde\gamma'$ and $\bar \gamma \cap \bar \gamma' \in\{\emptyset,\bar \gamma, \bar \gamma'\}$.
\item\label{enum:ii}
If $\gep= \gep'$ and $\bar\gamma\cap \bar \gamma'= \emptyset$, then
then $\bar \gamma'\cap \Delta^+_{\gamma}=\emptyset$.
\item\label{enum:iii}
If $\gep\ne \gep'$ and $\bar \gamma'\subset \bar \gamma$ (\resp $\bar \gamma\subset \bar \gamma'$) then
$\bar \gamma'\cap \Delta^-_{\gamma}=\emptyset$ (\resp $\bar \gamma\cap \Delta^-_{\gamma'}=\emptyset$).
\end{enumeratei}
This first condition simply states that compatible contours do not cross each-other.
The conditions $\bar \gamma'\cap \Delta^+_{\gamma}=\emptyset$ and $\bar \gamma'\cap \Delta^-_{\gamma}=\emptyset$ in \eqref{enum:ii} and \eqref{enum:iii}
can be reformulated as: $\tilde \gamma$ and $\tilde \gamma'$ do not share edges, and if both $\tilde \gamma$ and $\tilde \gamma'$ possess
two edges adjacent to one vertex
$x^*\in (\bbZ^2)^*$
then the two edges in $\gamma$ are linked and so are those in $\gamma'$.

\begin{figure}[ht]
\begin{center}
\includegraphics[scale=.3]{figures/contour}
\end{center}
\caption{\label{compa}\smaller
A contour $\gamma$ represented with its internal (circles) and external (squares) neighborhood.
To be compatible with $\gamma$, a contour $\gamma'$ of the same sign such that $\bar \gamma'\cap \bar \gamma= \emptyset$ cannot enclose any squares.
A compatible contour of opposite sign enclosed in $\gamma$ (such that $\bar \gamma'\subset \bar \gamma$) cannot enclose any circles.
}
\end{figure}

Note that the compatibility of two cylinders does not depend on their respective intensity, so that the notion can naturally be extended to signed contours:
The contours $\gamma$ and $\gamma'$ are said to be compatible (we write $\gamma \mid \gamma'$) if the cylinders $(\gamma,1)$ and $(\gamma',1)$ are.
Two distinct non-compatible contours are said to be \textit{connected} (we write $\gamma \perp \gamma'$).

If $\bC_1$ and $\bC_2$ are two finite collections of contours (compatible or not) we say that~$\bC_1$ is \textit{compatible with} $\bC_2$ and write
$\bC_1 \mid \bC_2$ if
\begin{equation}\label{compatiblesets}
\forall \gamma_1 \in \bC_1, \forall \gamma_2\in \bC_2, \quad \gamma_1 \mid \gamma_2.
\end{equation}
If \eqref{compatiblesets} does not holds we say that $\bC_1$ and $\bC_2$ are connected and write $\bC_1 \perp \bC_2$.
For a contour $\gamma$ and a collection $\bC$ with use the notation $\gamma \perp \bC$ and $\gamma \mid \bC$ for
$\{\gamma\} \perp \bC$ and $\{\gamma\} \mid \bC$

A (finite or countable) collection of cylinders (or of signed contours)
is said to be a compatible collection if its elements are pairwise compatible (see Figure \ref{grid}).
The reader can check by inspection that the following result holds. In
particular it establishes that the set of compatible collections of cylinders is in bijection with $\gO_{\gL}$
(simple connectivity is required to avoid having level lines enclosing holes).

\begin{lemma}\label{toott}
If $\gL$ is simply connected, then for any $\phi\in \gO_{\gL}$, $\hat \Upsilon_n(\phi)$ is a compatible collection of cylinders
and reciprocally, if $\hat \gG \subset \hat \cC_{\gL}$ is a compatible collection of cylinder in $\gL$
then its elements are the cylinders of the function
$\sum_{\hat \gamma\in \hat \gG} \varphi_{\hat \gamma}$.
\end{lemma}

Using \eqref{express} and the contour representation above, we can rewrite the partition function $\cZ_{\gL,\gb}$ in a new form.
We let $\cK(\gL)$ and $\hat \cK(\gL)$ denote the set of compatible collections of contour and cylinders in $\gL$.
We have
\begin{equation*}
\cZ_{\gL,\gb}= \sum_{\hat \gG \in \hat \cK(\gL)} \prod_{\hat \gamma \in \hat \gG} e^{-k(\hat \gamma)\gb|\tilde \gamma|}.
\end{equation*}
Summing over all the possible intensities, we obtain
\begin{equation}\label{contourdecomp}
\cZ_{\gL,\gb}=\sum_{\gG \in \cK(\gL)} \prod_{\gamma \in \gG} \frac{1}{e^{\gb|\tilde \gamma|}-1}.
\end{equation}
This last representation of the partition function is suitable to apply the cluster expansion techniques which we introduce in the next section.

\begin{figure}[ht]
\begin{center}
\includegraphics[scale=.3]{figures/grid}
\end{center}
\caption{\label{grid}\smaller
A compatible collection of contour on the dual lattice (the primal lattice is displayed is dotted lines).
Contours of different signs are displayed in different colors (red-dotted/blue-solid).
The primal lattice is represented in dotted line.
}
\end{figure}

We end this section by introducing a notion which will be of fundamental use in our proofs, and a few notation.
Given $\gG$ a compatible collection of contour and $\gamma \in \gG$, we say that $\gamma$ is an \textit{external contour} in $\gG$
if $\bar \gamma$ is maximal in $\gG$ for the inclusion, that~is,%
\begin{equation}\label{defexternal}
\forall \gamma'\in \gG, \quad \bar \gamma' \subset \bar\gamma \text{ or } \bar \gamma' \cap \bar \gamma= \emptyset.
\end{equation}

We say that $\gG$ is a \textit{compatible collection of external contours} if it is a compatible collection and every contour of $\gG$ is external in $\gG$.
Given $\bL$ a finite set of contours, we let $\cK(\bL)$ denote the set of compatible collections of contours included in $\bL$ and
$\cK_{\ext}(\bL)$ denote the set of compatible collection of external contours
(we use the notation $\cK_{\ext}(\gL)$ for $\bL=\cC_{\gL}$).
Given $\phi\in \gO_{\gL}$, we define
\begin{equation*}%\label{defuext}
\Upsilon^{\ext}_n(\phi):=\{ \gamma \in \Upsilon_n(\phi) \sep \gamma \text{ is external in } \Upsilon_n(\phi) \}.
\end{equation*}
Obviously $ \Upsilon^{\ext}_n(\phi)\in \cK_{\ext}(\gL)$.
We say that two contours $\gamma_1$ and $\gamma_2$ are \textit{externally compatible} if they are compatible and
$\bar \gamma_1\cap \bar \gamma_2=\emptyset$. We use the notation $\gamma_1\parallel \gamma_2$.
We~say that two collections
$\gG_1,\ \gG_2 \in \cK_{\ext}(\bL)$ are externally compatible if
\begin{equation*}%\label{extcompatibl}
\forall \gamma_1\in \gG_1, \ \forall \gamma_2\in \gG_2, \quad \gamma_1 \parallel \gamma_2,
\end{equation*}
or equivalently if $\gG_1\cap \gG_2=\emptyset$ and
$\gG_1\cup \gG_2\in \cK_{\ext}(\bL)$.
We also use the notation $\gG_1 \parallel \gG_2$ for external compatibility between contour collections, and also $\gamma \parallel \gG$ for $\{ \gamma \} \parallel \gG$.

\subsection{Cluster expansion}\label{clusexp}

Partition functions which can be written as a sum over collections of compatible geometric objects such as \eqref{contourdecomp}
appears in a variety of situation in statistical mechanics. A powerful method called \textit{cluster expansion} has been engineered to analyze the associated systems in the
low temperature regime (that corresponds to large $\gb$).
We introduce it here as it appears in \cite{cf:KP}, with a set of notation adapted to our context.

Recall that $\cC$ is the set of contours in $\bbZ^2$ and let $w: \cC \to \bbR_+$ be an arbitrary function (in full generality $w$ could assume complex values, \cf \cite{cf:KP}).
Recall that $\cK(\bL)$ denote the set of compatible collections of contours in $\bL$.
Given a finite subset ${\bL}$ of~$\cC$, the partition function associated to $w$ and $\bL$, $Z[\bL,w]$ is defined
by
\begin{equation}\label{weightpar}
Z[\bL,w]:= \sum_{\gG\in \cK(\bL)} \prod_{\gamma \in \gG} w(\gamma).
\end{equation}
For $\gL$ a subset of $\bbZ^2$ we write $Z[\gL,w]$ for $Z[\cC_{\gL},w]$.

We consider consider also $\bbP^w_{\bL}$ the probability measure on $\cK(\bL)$ corresponding to $Z[\bL,w]$ and call $\Upsilon$ the associated random variable.
The distribution $\bbP^w_{\bL}$ has its support in $\cK(\bL)$ and we have for $\gG\in \cK(\bL)$,
\begin{equation*}%\label{pwl}
\bbP^w_{\bL}(\Upsilon=\gG):=\frac{1}{Z[\bL,w]}\prod_{\gamma\in \gG}w(\gamma),
\end{equation*}
We use the notation $\bbP^w_{\gL}$ when $\bL:=\cC_{\gL}$ for $\gL$ a finite subset of $\bbZ$.

\begin{rem}\label{teorema}
Going back to \eqref{contourdecomp}, the reader can check that the distribution of $\Upsilon(\phi)$ under
$\bP_{\gL,\gb}$, $\gL$ simply connected, is given by $\bbP^{w_{\gb}}_\gL$ where
\begin{equation}\label{specialweight}
w_{\gb}(\gamma):=\frac{1}{e^{\gb|\tilde \gamma|}-1}.
\end{equation}
\end{rem}
\subsubsection{The key result}

The starting point of cluster expansion is to observe that the $\log$ of a partition function
can be expressed as sum over geometric objects called \textit{clusters}.
A cluster of contour $\bC$ in $\bL$ is a finite non-empty subset of $\bL$ which cannot be split into two compatible parts (recall \eqref{compatiblesets})
or more formally which satisfies
\begin{equation}\label{noncluster}
\forall \bB \subset \bC, \quad \bB \perp (\bC \setminus \bB).
\end{equation}

We let $\cQ(\bL)$ denote the set of clusters in $\bL$ and $\cQ$ the set of all clusters (finite subsets of $\cC$).
The starting point of cluster expansion is the observation that $\log Z[\bL,w]$ can be written as a sum over clusters
\begin{equation}\label{expansion}
\log Z[\bL,w]:= \sum_{\bC \in \cQ(\bL)} w^T(\bC),
\end{equation}
where the modified weights $w^T$ are given by
\begin{equation*}%\label{modifi}
w^T(\bC):= \sum_{\bB\in \cP(\bC)} (-1)^{|\bB|+|\bC|} \log Z[\bB,w],
\end{equation*}
where $\cP$ stands for the set of parts. In fact \eqref{expansion} is almost immediate if we consider the sum over all subsets of $\bL$. Then
one can check that $w^T(\bC)=0$ if $\bC$ is not a cluster (we refer to the first lines in \cite[\S 3]{cf:KP} for full details).

The reason why the expansion \eqref{expansion} is relevant is that if the original weights $w$ are small in a certain sense, and in particular decay exponentially with the
length of the contours, then the modified weights $w^T(\bC)$ are also small and decay exponentially fast with the \textit{total length of the cluster} $L(\bC)$,
defined as follows
\begin{equation*}
L(\bC)=\sum_{\gamma\in \bC} |\tilde \gamma|.
\end{equation*}
The powerful estimate displayed below is the main result of \cite{cf:KP}.

\begin{theorema}\label{superclust}
If there exist two functions $a$ and $d$, $\cC\to \bbR_+$, such that for every $\gamma$%
\begin{equation}\label{condit}
\sum_{\{\gamma'\in \cC \sep \gamma' \perp \gamma\}} e^{a(\gamma')+d(\gamma')}w(\gamma')\le a(\gamma)
\end{equation}
then
\begin{equation}\label{dabound}
\sum_{\{ \bC\in \cQ \sep \bC \perp \gamma \}} |w^T(\bC)|\exp\biggl(\sum_{\gamma'\in\bC}d(\gamma')\biggr) \le a(\gamma).
\end{equation}
\end{theorema}

\begin{rem}\label{laremarc}
For simplicity we introduced the result for the notion of contour compatibility/connectedness defined in Section \ref{contour}.
However the result is purely algebraic and is remains valid if compatibility is replaced by another symmetric relation on contours and
an the notion of cluster is defined using this other relation.
In the present paper we use the result with compatibility replaced by external compatibility in the proof of Lemma \ref{dzip}.
\end{rem}

For all practical purpose, in the remainder of the paper, we use the criterion \eqref{condit} for a pair of simple functions\vspace*{-3pt}
\[
a_0(\gamma)=|\tilde \gamma| \quad\text{and}\quad d_0(\gamma):= \left(\gb-5\right) |\tilde \gamma|,
\]
with $\gb>5$.
A simple and practical way of verifying condition \eqref{condit} in that case is to check for every $x^*\in (\bbZ^2)^*$\vspace*{-3pt}
\begin{equation}\label{usualchoice}
\sum_{\gamma\in \cC \sep x^*\in \gamma} e^{(\gb-4)|\tilde \gamma|} w(\gamma)\le 1.
\end{equation}
Given $\bC$ a set of contours, let us use the notation $x^*\in \bC$ for
\begin{equation*}
\exists \gamma \in \bC, \quad x^*\in\gamma.
\end{equation*}
We let the reader check that
that for $x^*\in (\bbZ^2)^*$, any clusters which satisfies $x^*\in \bC$ is incompatible with a contour of length $4$ which displays $x^*$ in its top right corner
(the choice for the sign being left open).
Applying \eqref{dabound} for these two contours of length $4$, provided that \eqref{condit} holds for \eqref{usualchoice}
we obtain thus that
\begin{equation}\label{dabound2}
\sum_{ \bC\in \cQ \sep x^*\in \bC } |w^T(\bC)|e^{(\gb-5)L(\bC)} \le 8.
\end{equation}

For the Solid-On-Solid model without constraint, which corresponds to the weight function \eqref{specialweight},
one can check that
\eqref{usualchoice} holds provided $\gb>5$.

The results mentioned in the rest of the section are classical consequences of Theorem \ref{superclust},
but are sometimes exposed in the literature in a way that does not exactly fit the needs of our paper.
For the sake of completeness, we prove these corollaries in the appendix.

\subsubsection{Free energy and boundary effects}\label{fffa}
In our analysis we will be only interested in the case of \textit{translation invariant} weight functions $w$,
meaning that $w(\gamma+x)=w(\gamma)$ for $x\in \bbZ^2$ where $\gamma+x$ is defined as the contour with the same sign as $\gamma$,
with the set of edges obtained by translating every edges of
$\tilde\gamma$ by $x$.

If the partition function of a statistical mechanics model has an expression of the form \eqref{weightpar},
with translation invariant weights, the cluster expansion yields a simple expression for the free energy of the associated model.
Assuming that $w$ is a translation invariant weight function which satisfies \eqref{usualchoice},
the following limit exists\vspace*{-3pt}
\begin{equation}\label{lafreen}
\limtwo{|\gL| \to \infty}{|\partial\gL|/|\gL|\to 0} \frac{1}{|\gL|} \log Z[ \gL,w]=\tf(w).
\end{equation}
More precisely we have for an arbitrary point $x^*$ in the dual lattice $(\bbZ^2)^*$
\begin{equation}\label{fraconv}
\tf(w)=\sum_{\bC \in \cQ(\bL) \sep x^*\in \bC} \frac{1}{|\bC|}w^T(\bC),
\end{equation}
where $|\bC|:=\{y^*\in (\bbZ^2)^*\sep y^*\in \bC \}$ is the number of points in the dual lattice which are visited by a contour in $\bC$
(note that $|\bC|\le L(\bC)$ and that the inequality can be strict). The above expression does not depend on $x^*$ by translation invariance.

The fact that the sum in \eqref{fraconv} converges is a consequence of \eqref{dabound2}.
This convergence result can simply be obtained by controlling the difference between the expression given for
$|\gL|\tf(w)$ and $\log Z[ \gL,w]$ using \eqref{dabound}. This difference can be shown to be proportional to the size of the boundary.

We do not prove \eqref{lafreen} but present instead a very similar result for another kind of partition function. Given $\gamma\in \cC$, we let $Z[\gamma,w]$ denote the
partition function corresponding to the set of contours in the domain $\bar \gamma$ which are compatible with $\gamma$,
\begin{equation*}%\label{cgammadef}
\cC_{\gamma}:= \{ \gamma'\in \cC \sep \bar \gamma'\subset \bar\gamma \text{ and } \gamma'\mid \gamma\}.
\end{equation*}

\begin{lemma}\label{finalfrontier}
If $w$ is a translation invariant weight function which satisfies Equation~\eqref{usualchoice} for $\gb$ sufficiently large
we have
\begin{equation*}
\bigl |\log Z[\gamma,w]- |\bar\gamma|\tf(w) \bigr|\le \frac 1 4\, |\tilde \gamma|.
\end{equation*}
\end{lemma}
The proof of Lemma \ref{finalfrontier} is displayed in Appendix \ref{pfinal}, and \eqref{lafreen} can be obtained with only minor modifications.

\subsubsection{Correlation decay and infinite volume limits}\label{fffb}

We say that a countable collection of contours $\gG\subset \cC$ is locally finite if
\begin{equation}\label{locallyfinite}
\forall x\in \bbZ^2, \quad \#\{ \gamma \in \gG \sep x\in \bar \gamma\}<\infty.
\end{equation}
We let $\cK$ denote the set of \textit{locally finite} compatible collection of contours on $\bbZ^2$.
We~say that a function $f: \cK\to \bbR$ is a local function if there exists a finite set
$A \subset \bbZ^2$ such that $f(\gG)$ is entirely determined by $\gG\cap \cC'_A$,
where
\[
\cC'_A:= \{ \gamma \in \cC \sep \bar \gamma \cap A\ne \emptyset \}.
\]
This is equivalent to say that there exists
$\hat f: \cK[\cC'_A]\to \bbR$ such that
\[
f(\gG)=\hat f(\gG\cap \cC'_{A}).
\]
Given $\gG\in \cK[\cC'_{A}]$ we obtain as a consequence of \eqref{expansion} that
\begin{equation}\label{godstim}
\bbP^w_{\bL}[\Upsilon \cap \cC'_{A}=\gG]=w_{\bL}(\gG)\,\frac{Z[\bL'_{A,\gG}]}{Z[\bL]}\\
=
w_{\bL}(\gG)\exp\biggl(-\sum_{\bC\in \cQ(\bL,A,\gG)} w^T(\bC)\biggr),
\end{equation}
where\vspace*{-3pt}
\begin{equation*}\begin{split}
w_\bL(\gG)&:=\ind_{\{\gG\subset \bL\}}\prod_{\gamma\in \gG}w(\gamma),\\
\bL'_{A,\gG}&:=\{ \gamma \in \bL \setminus \cC'_A \sep \gamma \text{ is compatible with } \gG\},
\end{split}
\end{equation*}
and $\cQ(\bL,A,\gG)$ is the set of clusters that either intersect $A$ or are connected with $\gG$
\begin{equation*}
\begin{split}
\cQ(\bL,A,\gG)&:= \{\bC\in \cQ(\bL) \sep \bC\cap (\bL'_{A,\gG})^\cc\ne \emptyset \}\\
&\phantom{:}=\{\bC\in \cQ(\bL) \sep \exists \gamma\in \bC, \ \gamma\in \cC'_A \text{ or } \gamma \perp \gG \},\\
\end{split}
\end{equation*}
When \eqref{usualchoice} holds, using Equation \eqref{godstim},
we can prove two important consequences: Firstly, $\bbP^w_{\bL}$ converges to an infinite volume $\bbP^w$ limit when
$\bL$ exhaust $\cC$ with the convergence holding in the local sense (the expectation of every local function converges).
Secondly, the correlation between two local functions decays exponentially with the distance of their support.

The infinite volume limit $\bbP^w$ is defined via its finite dimensional projection (using Kolmogorov's extension theorem). It is the
unique probability on $\cK$ which satisfies for every finite subset $A$
\begin{equation}\label{godstiminf}
\bbP^w[\Upsilon \cap \cC'_{A}=\gG]=
w(\gG)\exp\biggl(-\sum_{\bC\in \cQ(A,\gG)} w^T(\bC)\biggr),
\end{equation}
where $w(\gG):=\prod_{\gamma\in \gG}w(\gamma)$ and
\begin{equation*}
\cQ(A,\gG):=\{\bC\in \cQ \sep \exists \gamma\in \bC, \ \gamma\in \cC'_A \text{ or } \gamma \perp \gG \}.
\end{equation*}
The convergence of the sum in the exponential in \eqref{godstiminf} is ensured by \eqref{dabound2}.
We show that the convergence occurs in an exponential fashion and that spatial correlation decay exponentially.
To state the result, we need to introduce the following notion of distance between finite subset of $\bbZ^2$
and the complement of a finite set of contours.
\begin{equation*}
d(A,\bL^{\cc}):= \min\{ x\in A, \gamma\in \cC\setminus \bL, \ \max_{y\in \bar \gamma} |x-y|\}.
\end{equation*}
Note that when $A$ is fixed, this distance grows to infinity when $\bL$ exhausts $\cC$.

\begin{proposition}\label{propinfi}
If $w$ is a translation invariant weight function which satisfies \eqref{usualchoice},
Then
for every pair of local functions $f$ and $g$, $\cK \to [0,1]$ with respective supports $A$ and $B$ and every $\bL$, $\bL'$ such that
$ d(A,\bL^\cc) \le d(A,(\bL')^\cc)$ we have
\begin{equation}\label{convergenz}
\begin{split}
\bigl| \bbP^w_{\bL}[f(\gG)]-\bbP^w_{\bL'}[f(\gG)]\bigr|\le &|A|e^{-(\gb/100)d(A,\bL^\cc)},\\
\bigl| \bbP^w_{\bL}[f(\gG)]-\bbP^w[f(\gG)]\bigr|\le &|A|e^{-(\gb/100)d(A,\bL)},
\end{split}
\end{equation}
and also
\begin{equation}\label{decayze}
\bigl|\bbP^w_{\bL}[f(\gG)g(\gG)]- \bbP^w_{\bL}[f(\gG)]\bbP^w_{\bL}[g(\gG)]\bigr|\le |A|e^{-(\gb/100)d(A,B)}.
\end{equation}
\end{proposition}

The proof of this result is displayed in Appendix \ref{ppropinf} for completeness.

\begin{rem}
Let us remark that the result can be applied to the weights given by~\eqref{specialweight} in order to obtain the convergence of the distribution
of contours associated with the measure $\bP_{\gL,\gb}$ when $\gL$ is simply connected.
A proof of Theorem \ref{infinitevol} can then be deduced from this result by noticing that conditioned to the set of contours, the heights of the cylinders are
independent geometric variables (see \eg \cite[Lem.\,4.3]{cf:part1}). The~reader can refer to the proof of Proposition \ref{conviark} to see
how results on the distribution of the field $\phi$ can be deduced from a result about the contour distribution.
\end{rem}

\subsection{Contour decomposition for the wetting problem} \label{contsec}

We face various obstacles when trying to obtain a decomposition similar to \eqref{contourdecomp}, for $\cZ^{n,h}_{\gL,\gb}$.
First because the function $\phi$ cannot be expressed directly from the contour collection $\Upsilon_n(\phi)$.
Opting for a representation using cylinders does not fully solve the problem,
since the quantities $\ind_{\{\phi(x)\ge 0\}}$ and $\ind_{\{ \phi(x)=0\}}$ which appear in the Hamiltonian cannot be fitted in the expansion,
because they depend on the set of contours in a highly non-local way.

The way out is to opt for a more abstract representation, where the contours in the sum do not correspond to the level line of $\phi$.
We obtain one such representation for each choice of boundary condition $n$
\begin{equation}\label{reprezent}
\cZ^{n,h}_{\gL,\gb}= \sum_{\gG\in \cK(\gL)} \prod_{\gamma\in \gG} w^h_n(\gamma).
\end{equation}
Let us stress that using this type of contour decomposition is not a new idea, and that our
weight function is very similar to the ones used \eg in \cite{cf:ADM, cf:CM}.

In order to provide the expression of the weights $w^h_n$ (displayed in \eqref{defw}) we need to introduce a few notation.
Given $\gamma$ a contour, $n\in \bbZ_+$, we let $\gO[\gamma,n]$ and $\bar \gO[\gamma,n]$ denote the sets of functions in $\bar \gamma$ defined as follows
\begin{equation}\label{defoom}
\begin{split}
\gO[\gamma,n]&:=\{\phi,\ \bar \gamma \rightarrow \bbZ \sep \forall x\in \Delta^-_{\gamma},\ \gep(\gamma)(\phi(x)-n)\ge 0\},\\
\bar \gO[\gamma,n]&:=\gO[\gamma,n]\setminus \gO[\gamma,n+\gep(\gamma)]\\
&\phantom{:}= \{ \phi \in \gO[\gamma,n] \sep \exists x\in \Delta^-_{\gamma},\ \phi(x)=n\}.
\end{split}
\end{equation}
We define $\gO^+[\gamma,n]$, $\bar \gO^+[\gamma,n]$ as the restrictions of $\gO[\gamma,n]$ and $\bar \gO[\gamma,n]$ to the set of non-negative functions
(recall the convention adopted in
\eqref{positivity}).
The set $\bar \gO[\gamma,n]$ and $\gO[\gamma,n]$ can respectively be described as the sets of functions
$\phi$ such that $\Upsilon_n(\phi)$ resp. $\Upsilon_n(\phi)\setminus \{\gamma\}$
is compatible with $\gamma$.

Given $\gamma$, $n$ and $h>0$, we define $z^{h}_n(\gamma)$ and $\bar z^h_n(\gamma)$
to be the two partition functions associated with the sets $\gO^+[\gamma,n]$ and $\bar \gO^+[\gamma,n]$ and the energy functional
$\gb \cH^n_{\bar \gamma}(\phi)-h|\phi^{-1}\{0\}|$ (recall \eqref{defhamil})
\begin{equation}\label{partisioux}\begin{split}
z^h_n(\gamma)&:= \sum_{\phi\in \gO^+[\gamma,n]} e^{-\gb \cH^n_{\bar \gamma}(\phi)+h |\phi^{-1}\{0\}|},\\
\bar z^h_n(\gamma)&:= \sum_{\phi\in \bar \gO^+[\gamma,n]} e^{-\gb \cH^n_{\bar \gamma}(\phi)+h|\phi^{-1}\{0\}|}.
\end{split}\end{equation}
We extend the definition to the case of negative $n$ by setting $z^h_n=\bar z^h_n=0$ for $n<0$.
The reader can check that $\gO^+[\gamma,n]:=\bigcup_{k\ge 0}\bar\gO^+[\gamma,n+\gep(\gamma)k]$ and thus that
\begin{equation}\label{relatz}
z^h_n(\gamma)= \sum_{k\ge 0} e^{-k \gb|\tilde \gamma|}\bar z^h_{n+\gep(\gamma)k}(\gamma).
\end{equation}
We are now ready to define our contour weight $w^h_n(\gamma)$ (for $n\ge 0$) as follows
\begin{equation}\label{defw}
w^h_n(\gamma)= \frac{ e^{-\gb|\tilde \gamma|} z^h_{n+\gep(\gamma)}(\gamma)}{\bar z^h_n(\gamma)}.
\end{equation}
Note that, with our convention, negative contours have weight zero for $n=0$.
This definition turns out to be the most natural to obtain a contour representation for the partition function.

\begin{proposition}\label{trax}
The contour representation \eqref{reprezent} of the partition $\cZ^{n,h}_{\gL,\gb}$ holds true for the weights defined in \eqref{defw}
when $\gL$ is simply connected (recall \eqref{connectedness}).
\end{proposition}

While it involves some notation, the proof is not conceptually difficult.
The idea is to process recursively starting with external contours of the field $\phi$ \eqref{defexternal} and iterating the procedure.

\begin{proof}

The starting point of our proof is the observation that the complete descrip\-tion of
$\phi\in \gO^+_{\gL}$ can be obtained by knowing the set of external contours $\Upsilon^{\ext}_n(\phi)$ together with the associated intensity,
and the value of the restriction $\phi\restrict_{\bar \gamma}$ for every $\gamma\in \Upsilon^{\ext}_n(\phi)$.
When $\gamma$ is an external contour associated with boundary condition $n$ we have
$\phi\restrict_{\bar \gamma}\in \gO^+[\gamma,n+\gep(\gamma)]$ (recall \eqref{defoom}), and this is the only requirement that $\phi\restrict_{\bar \gamma}$
must satisfy.
Hence we obtain directly from \eqref{partisioux}
\begin{equation*}%\label{tapzero}
Z^{n,h}_{\gL,\gb}
=\sum_{\gG\in \cK_{\ext}(\gL)}\prod_{\gamma\in \gG} e^{-\gb |\tilde \gamma|} z^h_{n+\gep(\gamma)}(\gamma).
\end{equation*}
Using the definition \eqref{defw} we can rewrite the sum as
\begin{equation}\label{tapouz}
Z^{n,h}_{\gL,\gb}
=\sum_{\gG\in \cK_{\ext}(\gL)}\prod_{\gamma\in \gG} w^h_n(\gamma) \bar z^h_n(\gamma).
\end{equation}
Let us now introduce $\bar \cK_{\ext}(\gamma)$ which is the space in which the set of external contours associated to an element of $\bar \gO^+[n,\gamma]$ lies
\begin{equation}\label{daext}
\bar \cK_{\ext}(\gamma):=\{ \gG \in \cK_{\ext}(\bar \gamma) \sep \gG \mid \gamma \}.
\end{equation}
Decomposing according to the external contours of $\phi\restrict_{\bar \gamma}$ we obtain similarly to \eqref{partisioux} that
\begin{equation}\label{tapun}
\bar z^h_n(\gamma)= \sum_{\gG_1\in \bar \cK_{\ext}(\gamma)}\prod_{\gamma_1\in \gG_1}
w^h_n(\gamma_1) \bar z^h_n(\gamma_1).
\end{equation}
Injecting \eqref{tapun} in \eqref{tapouz} and iterating the procedure, we obtain \eqref{reprezent}.
\end{proof}

\subsection{Rewriting partition functions}

In order to obtain bounds on the contour weights $w^h_n(\gamma_1)$
which are sufficient to prove \eqref{usualchoice}, we have to use alternative expressions for the partition functions in order to facilitate
the comparison between
$z^h_{n+\gep(\gamma)}(\gamma)$ and $\bar z^h_n(\gamma)$.
One of the objective is to get rid the positivity constraint for~$\phi$.
Let us define for $\gG$ a finite subset of $\bbZ^2$,
\begin{equation*}%\label{ZZplus}
\cZ^+_{\gG}:=\cZ^0_{\gG,\gb}=\sum_{\phi\in \gO^+_{\gG}} \exp\left(-\gb \cH_\gL(\phi)\right),
\end{equation*}
and set
\begin{equation*}
\bar H(\gG):= \log \cZ^+_{\gG}-|\gG|\log \Bigl(\frac{e^{4\gb}}{e^{4\gb}-1}\Bigl).
\end{equation*}
We introduce the partition functions $z(\gamma)$, $\bar z(\gamma)$ which corresponds to the model without positivity constraint or interaction at level zero
\begin{equation}\label{tambem}
z(\gamma):=\sum_{\phi\in \gO[\gamma,n]}e^{-\gb\cH^n_{\bar \gamma}(\phi)}, \quad \bar z(\gamma):=\sum_{\phi\in \bar \gO[\gamma,n]}e^{-\gb\cH^n_{\bar \gamma}(\phi)},
\end{equation}
which, by translation invariance, do not depend on $n$. We consider the associated probability distributions $\bP^n_{\gamma}$ and $\bar \bP^n_{\gamma}$ on $\gO[\gamma,n]$ and $\bar \gO[\gamma,n]$.

\begin{lemma}\label{rwo}
We have for any $n\ge 1$
\begin{equation}\label{swoutch}\begin{split}
z^h_n(\gamma)&=\sum_{\phi\in \gO[\gamma,n]} e^{-\gb \cH^n_{\bar \gamma}(\phi)+ u|\phi^{-1}(\bbZ_-)|-\bar H(\phi^{-1}(\bbZ_-))},\\
\bar z^h_n(\gamma)&=\sum_{\phi\in \bar \gO[\gamma,n]} e^{-\gb \cH^n_{\bar \gamma}(\phi)+ u|\phi^{-1}(\bbZ_-)|-\bar H(\phi^{-1}(\bbZ_-))},
\end{split}\end{equation}
where in the formula above $u=u_h:= h- \log \left(\spfrac{e^{4\gb}}{e^{4\gb}-1}\right)$.
Alternatively we can write
\begin{equation}\label{dooz}\begin{split}
z^h_n(\gamma)&= z(\gamma)\bE^n_{\gamma}\bigl[e^{u|\phi^{-1}(\bbZ_-)|-\bar H(\phi^{-1}(\bbZ_-))}\bigr],\\
\bar z^h_n(\gamma)&= z(\gamma)\bE^n_{\gamma}\bigl[e^{u|\phi^{-1}(\bbZ_-)|-\bar H(\phi^{-1}(\bbZ_-))}\ind_{\{\exists x\in \Delta^-_{\gamma}, \phi(x)=n\}} \bigr],\\
&=\bar z(\gamma)\bar \bE^n_{\gamma}\bigl[e^{u|\phi^{-1}(\bbZ_-)|-\bar H(\phi^{-1}(\bbZ_-))}\bigr].
\end{split}\end{equation}
\end{lemma}

\begin{proof}
The statement \eqref{swoutch} can be proved in the same manner as \cite[Lem.\,3.1 \& Lem.\,3.2]{cf:part1}, and \eqref{dooz} is an obvious consequence of it.
\end{proof}

It follows from the definition of $\bar H$ that if
$\gG=\bigcup_{i\in \lint 1, m\rint} \gG_i$ is the decomposition of~$\gG$ into maximal connected components (in $\bbZ^2$) then
\begin{equation}\label{maxdecomp}
\bar H(\gG)=\sum_{i\in \lint 1, m\rint} \bar H(\gG_i).
\end{equation}
For our purpose we need in fact to estimate sharply the value of $\bar H$ only for connected components of size one and two,
and to have a rougher estimate for other connected sets.

\begin{lemma}\label{zoomats}

We have for any two neighboring points $x\sim y$ in $\bbZ^2$
\begin{equation}\label{sitmats}
\bar H\{x\}=0 \quad \text{ and } \quad \bar H\{x,y\}=\log \Bigl(\frac{1-J^4}{1-J^3}\Bigr).
\end{equation}
For $\gb$ sufficiently large,
for all $|\gG|$ connected and larger than $2$ we have
\begin{equation}\label{morseu}
0\le \bar H(\gG)\le 2J^2|\gG|.
\end{equation}
\end{lemma}

\begin{proof}
The equalities in \eqref{sitmats} are the result of a direct computation whose details are given in the proof of \cite[Lem.\,3.2]{cf:part1}.
For \eqref{morseu} the lower bound is a consequence of the super-additivity of $\bar H$ (see also \cite[Lem.\,3.2]{cf:part1}). For the upper bound we use the expansion \eqref{expansion} to evaluate the partition function of SOS which corresponds to weight function $w_{\gb}$ given in \eqref{specialweight}.
We have
\begin{equation*}
\log \cZ^+_{\gG}\le \log \cZ_{\gG,\gb}:=\sum_{\bC\in \cQ(\gG)} w_{\gb}^T(\bC).
\end{equation*}
As \eqref{usualchoice} is valid for $w_{\gb}$ if $\gb$ is sufficiently large, \eqref{dabound2} implies that
\begin{equation}\label{treck}
\sum_{\bC\in \cQ(\gG)} w_{\gb}^T(\bC)\le |\gG|\left[ 2e^{-4\gb}+O(e^{-6\gb})\right],
\end{equation}
which is sufficient to conclude.
\end{proof}

\subsection{Peak probabilities}\label{peaksec}

We recall here a result concerning the asymptotic probability of observing ``peaks'' of a given shape for $\phi$ under the measure $\bP_{\gL,\gb}$.
We provide a result which is slightly more general than the one proposed in \cite[Prop.\,4.5]{cf:part1}.
Given a finite set $\bL$ of contour included in $\cL_\gL$, and $\gb>0$, we define $\bP_{\bL,\gL,\gb}$ to be a measure on $\gO_{\gL}$
which can be sampled as follows.
\begin{enumerate}
\item [(A)] Sample a set of contours $\Upsilon$ according to the measure $\bbP^{w_{\gb}}_{\bL}$ (recall \eqref{specialweight}).
\item [(B)] For each contour $\gamma\in \Upsilon$ sample independently a geometric variable $k(\gamma)$ satisfying
$\bbP[k(\gamma)=i]=[w_{\gb}(\gamma)]^{-1} e^{-\gb|\tilde \gamma|i}$.
\item [(C)] Set (recall \eqref{cylfunc})
\[
\phi:=\sum_{\gamma\in \Upsilon} \varphi_{(\gamma,k(\gamma))}.
\]
\end{enumerate}
Note that when $\bL=\cC_{\gL}$ we have $\bP_{\bL,\gL,\gb}=\bP_{\gL,\gb}$. The probability distribution~$\bP^0_{\gamma}$ and~$\bar \bP^0_{\gamma}$ defined below
Equation \eqref{tambem} are also of the form $\bP_{\bL,\bar \gamma,\gb}$ for adequate choices of~$\bL$.
This definition thus enables us us to treat measures which include special boundary condition or contour restriction.

\begin{proposition}\label{rouxrou}
If $\gb$ is sufficiently large, then for any choice $\gL$, $\bL$ and $n$ and any triple of distinct vertices $(x,y,z)\in \gL^3$
such that $x\sim y \sim z$ we have
\begin{equation}\label{zups}
\begin{split}
\bP_{\bL,\gL,\gb}[\phi(x)\ge n]&\le 2 e^{-4\gb n},\\
\bP_{\bL,\gL,\gb}[\min(\phi(x),\phi(y)) \ge n]&\le 2 e^{-6\gb n},\\
\bP_{\bL,\gL,\gb}[\min(\phi(x),\phi(y),\phi(z))\ge n]&\le 2 ne^{-8\gb n}.
\end{split}
\end{equation}
If we assume in addition that $\bL$ contains the positive contour of length $4$ enclosing $x$, then
\begin{equation}\label{infin}
\bP_{\bL,\gL,\gb}[\phi(x)\ge n]\ge \frac{1}{2}\, e^{-4\gb n}.
\end{equation}
If we assume that $\bL$ contains the positive contour of length $6$ enclosing $x$ and $y$, then
\begin{equation}\label{infde}
\bP_{\bL,\gL,\gb}[\min (\phi(x),\phi(y))\ge n]\ge \frac{1}{2}\, e^{-6\gb n}.
\end{equation}
\end{proposition}

The proof of \eqref{zups} is identical to that of \cite[Prop.\,4.5]{cf:part1}. The proofs of \eqref{infin} and \eqref{infde} are detailed in Appendix \ref{lazt}.

\subsection{Monotonicity and the FKG inequality}

The set $\gO_{\gL}$ as well as its variants ($\gO^+_{\gL}$ and others introduced later in the paper) are naturally equipped with an order defined as follows:
\[
\phi\le \phi' \iff \forall x\in \gL,\ \phi(x)\le \phi'(x).
\]
Using this order we can define a notion of increasing function ($f$ is increasing if $\phi\le \phi'\Rightarrow f(\phi)\le f(\phi')$) and of increasing event
($A$ is increasing if the function $\ind_A$~is).
We say that a probability measure $\mu$ on $\gO_{\gL}$ \textit{stochastically dominates} another one~$\mu'$ (we write $\mu \succcurlyeq \mu'$) if for any increasing function $f$,
\[
\mu(f(\phi))\ge \mu'(f(\phi)).
\]

The FKG inequality allows us to say that if a probability measure $\mu$ supported on a subset of $\gO_{\gL}$ satisfies a certain condition,
increasing functions are positively correlated.
For the inequality to be satisfied \cite{cf:Holley}, we need the support of $\mu$ to be a \textit{distributive lattice}, that is, to
be stable over the operations $\vee$ and $\wedge$
defined by
\[
(\phi_1\vee \phi_2)(x):=\max(\phi_1(x),\phi_2(x)) \quad\text{and}\quad(\phi_1\wedge \phi_2)(x):=\min(\phi_1(x),\phi_2(x)).
\]
Moreover the probability considered needs to verify Holley's condition \cite[Eq.\,(7)]{cf:Holley},
\begin{equation*}%\label{hcond}
\mu(\phi_1\vee \phi_2) \mu(\phi_1\wedge \phi_2)\ge \mu(\phi_1)\mu(\phi_2).
\end{equation*}
If this is satisfied then for any pair of increasing functions $f$ and $g$ we have
\begin{equation*}%\label{FKeq}
\mu(f(\phi)g(\phi))\ge \mu(f(\phi)) \mu(g(\phi)).
\end{equation*}

We obtain as immediate consequences of the FKG inequality, several stochastic domination results.
Given $\psi\in \gO_{\infty}$ and $\gL\subset \bbZ^2$, $\gb>0$ and $h\in \bbR$,
we let~$\tilde \bP^{\psi,h}_{\gb,\gL}$ denote a measure defined on a subset $\tilde \gO_\gL \subset \gO^+_{\gL}$ (recall \eqref{positivity}) which is a
distributive lattice, with the probability of each state proportional to the Gibbs weight
\hbox{$\exp\bigl(-\gb\cH^\psi_\gL(\phi)+ h|\phi^{-1}\{0\}|\bigr)$}.

\begin{cor}\label{FKalt}
The following holds.
\begin{enumeratei}
\item
For any increasing event $A$
\begin{equation*}%\label{FKG}
\tilde \bP^{\psi,h}_{\gb,\gL}[ \ \cdot \ | \ A] \succcurlyeq \tilde \bP^{\psi,h}_{\gb,\gL}.
\end{equation*}
\item
For any $h'>h$
\begin{equation*}
\tilde \bP^{\psi,h'}_{\gb,\gL} \preccurlyeq \tilde \bP^{\psi,h}_{\gb,\gL}.
\end{equation*}
\item
For any $\psi'\ge \psi$
\begin{equation*}
\tilde \bP^{\psi',h}_{\gb,\gL} \succcurlyeq \tilde \bP^{\psi,h}_{\gb,\gL}.
\end{equation*}
\end{enumeratei}
\end{cor}

\begin{proof}
The first point is immediate, for the other ones we simply have to notice that
$\exp((h-h')|\phi^{-1}\{0\}|)$ and $\exp\bigl(\gb \bigl(\cH^\psi_\gL(\phi)-\cH^{\psi'}_\gL(\phi) \bigr) \bigr)$ are increasing functions.
\end{proof}

\section{Organization of the proof of Theorem \ref{main}}\label{zorga}

We start with a small notational remark. As our main result concerns the behavior of the free energy close to $h_w(\gb)$,
it is more convenient for us to work as in the statement of Theorem \ref{main} with the parameter $u=h-h_w(\gb)$ than with $h$.
Therefore in most cases we work with all quantities defined as functions of $u$ rather than $h$.
When $h$ appear in a computation, we always assume that
\[
h=h_u=: h_w(\gb)+u.
\]

\subsection{Contour stability and consequences}\label{stabcon}

If we want \eqref{reprezent} to yield information about the free energy, we need the contour weights $w^u_n(\gamma)$ to be small, or more precisely
we want \eqref{usualchoice} to be satisfied.
We say that a contour is $n$-stable for $u$ if
\begin{equation}\label{stability}
w^u_n(\gamma)\le e^{-(\gb-1) |\tilde \gamma|}.
\end{equation}

The most important part of our proof is to show that we can partition $\bbR_+$
into intervals $([u^*_{n+1},u^*_{n}])_{n\ge 0}$ (with the convention than $u^*_0=\infty$) in which all the contours are $n$-stable.
This result also plays a central role in our proof of Theorem \ref{Gibbs}.

\begin{theorem}\label{converteo}
When $\gb$ is sufficiently large,
there exists a decreasing sequence $(u^*_n)_{n\ge 1}$ satisfying
\[
u^*_n\in \Bigl[\frac{1}{200} J^{n+2}, 200 J^{n+2}\Bigr]
\]
such that all contours are $n$ stable for $u\in[u^*_{n+1},u^*_{n}]$.
\end{theorem}

We give a road-map for the proof of Theorem \ref{converteo} in Section \ref{sketch}, by presenting the main steps.
The detailed proof is then given in Section \ref{ponverteo}.

Note that the $n$-stability of all contour implies that \eqref{usualchoice} is satisfied for $w^u_n$.
Indeed, a classical counting argument shows that for $k\in \bbN$ even and $x^*\in (\bbZ^2)^*$
\begin{equation*}%\label{counting}
\#\{ \gamma\in \cC \sep |\tilde \gamma|=k,\ x^*\in \gamma\} \le 8.3^{k-2}\le 3^{k}
\end{equation*}
(starting from $x^*$ we have $2$ choices for the sign, $4$ choices for the first step, at most~$3$ for the other steps, and the last step is determined by the fact that $\gamma$ is a loop).

Thus combining Theorem \ref{converteo} with the results introduced in Sections \ref{fffa} and \ref{fffb}
we can derive consequences for the free energy and the measure $\bP^{n,h}_{\gL,\gb}$.
These consequences are detailed in Section \ref{dalast}.
We state here two statements which are of interest in the proof of Theorem \ref{main} which are respectively proved in Sections \ref{laprov} and~\ref{lapriv}.
First, we obtain a result concerning the regularity of the free energy.

\begin{proposition}\label{infinitz}
The free energy $u\mapsto \bar \tf(\gb,u)$ is infinitely differentiable on $(u^*_{n+1},u^*_{n})$ for $n\ge 0$
for all $n\ge 0$. Moreover all derivatives of $\tf(\gb,h)$ are uniformly bounded on $(u^*_{n+1},u^*_{n})$.
\end{proposition}

Secondly, we obtain an a priori bound on the derivative which, together with Theorem~\ref{oldmain}, leads to a
sharp asymptotic estimate on the layering transition points $u^*_n$.

\begin{proposition}\label{convirj}
Given $\gb>\gb_0$ sufficiently large, there exists a constant such that for every $n$ and every $u\in (u^*_{n+1},u^*_n)$ (where by convention $u^*_0=+\infty$)
we have
\begin{equation*}
\frac{1}{10}J^{2n}\le \partial_u\bar \tf(\gb,u)\le 10 J^{2n}.
\end{equation*}
In particular, $\bar \tf(\gb,u)$ is not differentiable at $u^*_n$.
\end{proposition}
To conclude the proof of Theorem \ref{main}, we provide a proof of \eqref{asimptic}.

\begin{proof}[Proof of \eqref{asimptic}]

We make use of Theorem \ref{oldmain}.
Recalling \eqref{defun}, Equation \eqref{lequivdelamor} implies that for every $u,v\in [u_{n+1}, u_n]$
\begin{equation*}
\bar \tf(\gb,v)-\bar\tf(\gb,u)= \alpha_1 J^{2n} (v-u)+o(J^{3n}).
\end{equation*}
By convexity, this implies that for $\delta>0$, for all $n$ sufficiently large
we have
\begin{equation*}
\forall u\in (u_{n+1}(1+\delta), u_{n}(1-\delta)), \quad \left| \partial_u \bar \tf(\gb,u)- \alpha_1 J^{2n}\right| \le 2\delta J^{2n}.
\end{equation*}
In view of Proposition \ref{convirj} and the of the fact that $\alpha_1, \alpha_2\!\in\![1/2,2]$ (\cf Proposi\-tion~\ref{rouxrou}), we can conclude that
for $n$ sufficiently large
\[
u^*_n\in [u_n(1-\delta), u_n(1+\delta)].\qedhere
\]
\end{proof}

\subsection{Truncated weights and road map to Theorem \ref{converteo}}\label{sketch}

Our first step is to prove $n$-stability in a reduced intervals.
We define for $n\ge 1$,
\begin{equation*}
u^{+}_n:= 200 J^{n+2}, \quad u^-_n:= \frac{1}{200} J^{n+2},
\end{equation*}
and also $u^{\pm}_0=\infty$. And we prove the following.
\begin{proposition}\label{truta}
For all $\gb$ sufficiently large
every contour is $n$-stable for $u\in [u^{+}_{n+1}, u^{-}_n]$.
\end{proposition}
Using this partial result, we can obtain a characterization of $u^*_n$. This requires introducing the notion of truncated weights and free energy
(we follow here ideas which were developed in \cite{cf:CM}).
We define the truncated weights $ w^{u,\tr}_n$ by
\begin{equation}\label{trunk}
w^{u,\tr}_n(\gamma):= \max\bigl(e^{-(\gb-1)|\tilde \gamma|}, w^u_n(\gamma)\bigr).
\end{equation}
We define in the same manner
\begin{equation*}
\cZ^{n,u,\tr}_{\gL,\gb}:= \sum_{\gG\in \cK(\gL)} \prod_{\gamma\in \gG} w^{\tr,u}_n(\gamma),
\end{equation*}
and the corresponding free energy
\begin{equation*}
\bar \tf^{\tr}_n(\gb,u):= \limtwo{|\gL|\to \infty}{|\partial\gL|/|\gL|\to 0} \frac{1}{|\gL|}\log \cZ^{n,u,\tr}_{\gL,\gb}-\tf(\gb).
\end{equation*}
In view of \eqref{reprezent}, we have for every $n$ and $u$
\begin{equation*}%\label{zinak}
\bar \tf^{\tr}_n(\gb,u)\le \bar\tf(\gb,u),
\end{equation*}
and equality is achieved if and only if all contours are stable (the \textit{only if} part of the statement may appear less obvious, but as we
do not use that fact in our proof, we~leave it as an exercise to the interested reader).
In particular, a simple consequence of Proposition \ref{truta} is the following.
\begin{cor}\label{simplecor}
For every $n\ge 0$ and every $u\in [u^{+}_{n+1}, u^-_n]$, we have
\[
\tf^{\tr}_{n}(\gb,u)=\tf(\gb,u).
\]
\end{cor}
Another important observation, that as the weights $w^{u,\tr}_n(\gamma)$ are continuous in $u$, so are the weights $w^T(\bC)$ associated to clusters.
Thus as the convergence \eqref{fraconv} is uniform, the function $u\mapsto \bar \tf^{\tr}_n(\gb,u)$ is continuous for every $n$.
Now, from Corollary~\ref{simplecor}, we have for any $n\ge 1$
\begin{equation*}\begin{split}
\tf^{\tr}_{n}(\gb,u^-_n)&=\tf(\gb,u^-_n)\ge \tf^{\tr}_{n-1}(\gb,u^-_n),\\
\tf^{\tr}_{n-1}(\gb,u^+_n)&=\tf(\gb,u^+_n)\ge \tf^{\tr}_{n}(\gb,u^+_n).
\end{split}\end{equation*}
Using the continuity of $[\tf^{\tr}_{n-1}-\tf^{\tr}_{n}](\gb,u)$ we define
\begin{equation}\label{caraca}
u^*_n:=\min\left\{v\in [u^-_n,u^+_n] \sep \tf^{\tr}_{n-1}(\gb,u)=\tf^{\tr}_{n}(\gb,u) \right\}.
\end{equation}
To complete the proof of Theorem \ref{converteo}, we need to extend the stability result to the interval $[u^*_n,u^*_{n+1}]$.
\begin{proposition}\label{mejor}
For all $\gb$ sufficiently large,
every contour is $n$-stable for $u\in[u^*_{n+1},u^*_n]$.
\end{proposition}

\begin{rem}
The characterization of $u^*_n$ as a $\min$ in \eqref{caraca}, is a bit arbitrary in the sense that the only requirements of the proof are $u\in [u^-_n,u^+_n]$ and
$\tf^{\tr}_{n-1}(\gb,u)=\tf^{\tr}_{n}(\gb,u)$. It does not mean however that there is any freedom in the choice of $u^*_n$,
as further results imply that
\begin{equation*}%\label{zuniks}
\left\{v\in [u^-_n,u^+_n] \sep \tf^{\tr}_{n-1}(\gb,u)=\tf^{\tr}_{n}(\gb,u) \right\}=\{u^*_n\}.
\end{equation*}
While the definition of the truncated potential also offers some degree of freedom, a~consequence of latter results is that the value of $u^*_n$ does not depend on the particular choice which is made for truncation.
\end{rem}

Proposition \ref{truta} turns out to be the more difficult statement as it requires quantitative estimates which prove to be quite technical.
Its extension, Proposition \ref{mejor}, is proved using softer arguments
combining a monotonicity statement (Lemma \ref{mozon} below) together with Lemma \ref{finalfrontier}.
For the proof of both Proposition \ref{truta} and Proposition~\ref{mejor}, an important building brick is the following monotonicity consideration.
\begin{lemma}\label{mozon}
For any $n\in \bbN$, we have:
\begin{enumeratei}
\item
For any positive contour, $u\mapsto w^u_n(\gamma)$ is decreasing in $u$.
\item
For any negative contour $u\mapsto w^u_n(\gamma)$ is increasing in $u$.
\end{enumeratei}
\end{lemma}
\begin{rem} \label{mouton}
As a consequence of the above statement for each contour, the proof of Propositions \ref{truta} and \ref{mejor} reduces to checking
stability for one value of $u$ which is chosen to be an extremity of the interval (the right extremity for negative contour, the left one for positive contour).
\end{rem}

\section{Proof of Theorem \ref{converteo}}\label{ponverteo}

In this section, we prove all the statements exposed in Section \ref{stabcon}.
We first prove Lemma~\ref{mozon} in Section \ref{mooz}, while the other subsections are devoted to the proof of Propositions \ref{truta} and \ref{mejor}.
Our proof for the contour's stability depends on the size of the contour.
This gives a utility to the following definition.
Here and in the remainder of the paper, $\Diam(\gamma)$ denotes the Euclidean diameter of the geometric contour $\tilde \gamma$ considered as a subset of $\bbR^2$.

\begin{definition}
A contour is said to be $n$-\textit{small} if $\Diam(\gamma)\le \lfloor \max(\gb,\gb n)^2 \rfloor$.
A~contour which is not small is said to be \textit{large}.
\end{definition}

The stability of small contours can relatively is proved directly ``by hand'' in Section~\ref{smallcproof}, using directly the estimates we have for the Solid-On-Solid measures.
The stability of large contours is proved in two steps. First, we restrict our proof to the interval $[u^{+}_{n+1},u^{-}_n]$ to prove Proposition \ref{truta}.
This is the most delicate part and it spreads from Section \ref{prezent} to Section \ref{dcimus} with the more technical computation postponed to Section
\ref{secrestrict}.
The last step of the proof of large contour stability is the extension to the full interval $[u^{*}_{n+1},u^{*}_n]$ to complete the proof of Proposition \ref{mejor}. This is done in Section~\ref{mejorplus}.

Note that whenever a contour is $n$-small we also have a bound on the enclosed area, which we are to use in most computations:
\begin{equation}\label{areabound}
|\bar \gamma| \le \Diam(\gamma)^2\le \max(\gb,\gb n)^4.
\end{equation}

\subsection{Proof of Lemma \ref{mozon}}\label{mooz}
Let us assume for simplicity that $\gamma$ is a positive contour (the proof for the negative case being identical).
We let $\bP^{n,u}_{\gamma}$ and $\bar \bP^{n,u}_{\gamma}$ be the respective probability on $\gO^+[\gamma,n]$, and $\bar \gO^+[\gamma,n]$ corresponding to the partition functions
$z^u_n(\gamma)$ and $\bar z^u_n(\gamma)$, that is, (recall $h=u+h_w(\gb)$)
\begin{equation}\label{defpnu}
\begin{split}
\bar \bP^{n,u}_{\gamma}(\phi)&:=\frac{1}{z^u_n} e^{-\gb \cH^n_{\bar \gamma}(\phi)+h |\phi^{-1}\{0\}|}, \quad \forall \phi \in \bar \gO^+[\gamma,n],\\
\bP^{n,u}_{\gamma}(\phi)&:=\frac{1}{z^u_n} e^{-\gb \cH^n_{\bar \gamma}(\phi)+h |\phi^{-1}\{0\}|}, \quad \forall \phi \in \gO^+[\gamma,n].
\end{split}
\end{equation}
Using these definitions, the reader can check that the logarithmic derivative of $w^u_n(\gamma)$ can be expressed in the following manner:
\begin{equation}\label{dirivatz}
\partial_u \log w^u_n(\gamma)= \bE^{n+1,u}_{\gamma} [ |\phi^{-1}(0)|]- \bar \bE^{n,u}_{\gamma} [|\phi^{-1}(0)|].
\end{equation}
As $|\phi^{-1}(0)|$ is a decreasing function of $\phi$,
if we show that $\bP^{n+1,u}_{\gamma}$ stochastically dominates $\bar \bP^{n,u}_{\gamma}$, then it implies that the r.h.s.\ in
\eqref{dirivatz} is negative which concludes the proof.
Let us introduce the events
\begin{equation*}
\begin{split}
\cA&:= \{ \forall x\in \Delta^-_{\gamma},\ \phi(x)\le n\},\\
\cB&:= \{ \forall x\in \Delta^-_{\gamma}, \ \phi(x)\ge n+1\}.
\end{split}
\end{equation*}
Observe that $\bar \bP^{n,u}_{\gamma}$ and $\bP^{n+1,u}_{\gamma}$ can both be defined as a conditioned variant of $\bP^{n,u}$. We have
\begin{equation*}
\bar \bP^{n,u}_{\gamma}=\bP^{n,u}_{\gamma}[ \ \cdot \ | \ \cA ] \quad\text{and}\quad \bP^{n+1,u}_{\gamma}=\bP^{n,u}_{\gamma}[ \ \cdot \ | \ \cB ].
\end{equation*}
Noting that $\gO^+[\gamma,n]$ is a distributive lattice, that $\cA$ is a decreasing event and that~$\cB$ is an increasing event we deduce from Corollary \ref{FKalt} that
\begin{equation*}
\bar \bP^{n,u}_{\gamma} \ldom \bP^{n,u}_{\gamma}\quad \text{ and } \quad \bP^{n,u}_{\gamma}\ldom \bP^{n+1,u}_{\gamma},
\end{equation*}
which is sufficient to conclude.

\subsection{Stability of $n$-Small contours}\label{smallcproof}

We prove the stability directly on a larger interval for the parameter $u$, so that it can be used for both Propositions \ref{truta} and \ref{mejor}.
More precisely the main statement proved in this section is the following.

\begin{proposition}\label{lesmall}
For $\gb$ sufficiently large, we have:
\begin{enumeratei}
\item
Every positive $n$-small contour is $n$-stable for $u=u^{-}_{n+1}$.
\item
Every negative $n$-small contour is $n$-stable for $u=u^{+}_n$.
\end{enumeratei}

\end{proposition}
This proposition combined with Lemma \ref{mozon} implies stability of positive and negative contours on the intervals $[u^{-}_{n+1},\infty)$ and
$(-\infty,u^{-}_{n}]$ respectively.
Both intervals include $[u^*_{n+1},u^*_n]$, which is sufficient
for Proposition \ref{mejor} and a fortiori for Proposition~\ref{truta}.

\begin{proof}
Here and a few other instances,
we have to treat separately the cases where level zero is involved: $n=0$, $\gamma$ positive, and $n=1$, $\gamma$ negative.
We need to show that when $\Diam(\gamma)\le \gb^2$ and $\gamma$ is positive we have
\begin{equation*}
\frac{z^{u^-_1}_{1}(\gamma)}{\bar z^{u^-_1}_0(\gamma)}\le e^{|\tilde \gamma|}.
\end{equation*}
Considering the contribution of the ground state $\phi\equiv 1$, we have $z^{u^-_1}_{0}(\gamma)\ge e^{h|\bar \gamma|}\ge 1$.
On the other hand, setting $h^-_1=u^-_1+h_w(\gb)$,
we have for $\gb$ sufficiently large (recall $J=e^{-2\gb}$)
\begin{multline}\label{truxc}
z^{u^-_1}_{1}(\gamma)\le \sum_{\phi\in \gO_{\bar \gamma}} e^{-\gb\cH^1_n(\phi)+h^-_1|\phi^{-1}(0)|}\\[-10pt]
\le
e^{h^-_1|\bar \gamma|} \cZ_{\bar \gamma,\gb}\le \bigl(e^{h^-_1}(1+3J^2) \bigr)^{|\bar \gamma|}\le
e^{J |\bar \gamma|}\le e^{J\gb^4}\le e.
\end{multline}
The third inequality is obtained by using \eqref{treck}.
The fourth and fifth inequality use \eqref{areabound} and are valid for $\gb$ sufficiently large (note that $h^-_1$, like $h_w(\gb)$ is of order $J^2$).
Similar computations can be used to prove that $z^{u^+_1}_{0}(\gamma)\le e $ and $z^{u^+_1}_{1}(\gamma)\ge 1$.

In all other cases ($n\ge 1$, $\gamma$ positive and $n\ge 2$, $\gamma$ negative)
we can rewrite the ratio of partition function, using \eqref{dooz} from Lemma \ref{rwo}. We obtain
\begin{equation}\label{topi}
\frac{z^u_{n+\gep(\gamma)}(\gamma)}{ \bar z^u_n(\gamma)}
= \frac{ \bE^{n+\gep(\gamma)}_{\gamma}\bigl[e^{u |\phi^{-1} (\bbZ_-)| -\bar H(\phi^{-1}(\bbZ_-))}\bigr]}{\bE^n_{\gb,\gamma}\bigl[e^{u |\phi^{-1}(\bbZ_-))-\bar H(\phi^{-1}(\bbZ_-))}\ind_{\{ \exists x\in \Delta^-_{\gamma},\phi(x)=n\}}\bigr]}.
\end{equation}
As $\bar H$ is positive (recall \eqref{morseu}), using \eqref{areabound}, we see that the numerator of the r.h.s.\ satisfies
\begin{equation*}%\label{gonza}
\bE^{n+\gep(\gamma)}_{\gamma}\bigl[e^{u |\phi^{-1} (\bbZ_-)| -\bar H(\phi^{-1}(\bbZ_-))}\bigr] \le e^{u|\bar \gamma|}\le e^{u \gb^4n^4},
\end{equation*}
where $n$-smallness of $\gamma$ is used in the last inequality.
Considering $u$ being either equal to $u^-_{n+1}$ or $u^+_n$ (depending on the value of $\gep(\gamma)$) we conclude that provided $\gb$ is sufficiently large
\begin{equation}\label{bip}
\bE^{n+\gep(\gamma)}_{\gamma}\bigl[e^{u |\phi^{-1} (\bbZ_-)| -\bar H(\phi^{-1}(\bbZ_-))}\bigr]\le e.
\end{equation}
For the denominator,
considering only the contribution of the event
\[
\{\phi \in \gO[n,\gamma] \sep \forall x\in \bar \gamma,\ \phi(x)\ge 1\},
\]
we obtain that
\begin{multline}\label{conque}
\bE^n_{\gamma}\bigl[e^{u |\phi^{-1}(\bbZ_-))-\bar H(\phi^{-1}(\bbZ_-))}\ind_{\{\exists x\in \Delta^-_{\gamma},\phi(x)=n\}}\bigr]\\
\ge
\bP^n_{\gamma}\left[\{\exists x\in \Delta^-_{\gamma},\phi(x)=n\}\cap \{\forall x\in \bar \gamma,\ \phi(x)\ge 1 \right].
\end{multline}
Using \eqref{zups} we obtain that for any $x_0\in \Delta^-_{\gamma}$
\begin{equation}\label{loot}
\bP^n_{\gamma}\left[\exists x\in \Delta^-_{\gamma},\phi(x)=n \right]\ge \bP^n_{\gamma}\left[\phi(x_0)=n\right]\ge 1-4e^{-4\gb}.
\end{equation}
Using \eqref{infin} we have
\begin{equation}\label{looz}
\bP^n_{\gamma} \left[ \exists x\in \bar \gamma, \phi(x)\le 0 \right]
\le \sum_{x\in \bar \gamma} \bP^n_{\gamma} [ \phi(x)\le 0]\le 2|\bar \gamma| e^{-4n\gb}\le 4n^4\gb^4 e^{-4n\gb}.
\end{equation}
Combining \eqref{conque},\eqref{loot} and \eqref{looz}, we obtain that for $\gb$ sufficiently large
\begin{equation}\label{bap}
\bE^n_{\gamma}\bigl[e^{u |\phi^{-1}(\bbZ_-)|-\bar H(\phi^{-1}(\bbZ_-))}\ind_{\{\exists x\in \Delta^-_{\gamma},\phi(x)=n\}}\bigr]\ge e^{-1},
\end{equation}
and thus we conclude from \eqref{topi}, \eqref{bip} and \eqref{bap} that
\[
\frac{z^u_{n+\gep(\gamma)}(\gamma)}{\bar z^u_n(\gamma)}\le e^2.\qedhere
\]
\end{proof}

\subsection{Larger contours: presenting the induction}\label{prezent}

To prove the stability of larger contour, we proceed with a double induction. A first induction based on the inclusion order for contours,
and a second one on the level $n$ for which stability is tested. To~avoid any confusion, before going into the details of the proof,
we provide the structure of this inductive reasoning.
For $n\ge 0$ and $\gamma$ a contour, we define the property
\begin{equation*}
\cP(n,\gamma):=\begin{cases}
\big[ \text{the contour } \gamma \text{ is stable at level } n \text{ for } u\ge u^+_{n+1} \big] &\text{if } \gep(\gamma)=+,\\
\big[ \text{the contour } \gamma \text{ is stable at level } n \text{ for } u\le u^-_{n} \big] &\text{if } \gep(\gamma)=-,
\end{cases}
\end{equation*}
and
\begin{equation*}
\cP(\gamma):=\big[ \ \cP(n,\gamma) \text{ is satisfied for all } n\ge 0 \big].
\end{equation*}
We are going to prove that $\cP(\gamma)$ holds for every contour using an induction on $\gamma$.
From Proposition \ref{lesmall}, we know that $\cP(\gamma)$ holds true when $\Diam(\gamma)\le \gb^2$.
Thus we only need to perform the induction step, which is proving $\cP(\gamma)$ assuming that
$\cP(\gamma')$ holds for all contours $\gamma'$ ``included in $\bar \gamma$'' \ie such that $\bar \gamma' \subset \bar \gamma$, $\bar \gamma'\ne \bar\gamma$.

To prove $\cP(\gamma)$ itself we use an induction on $n$.
The direction of the induction depends on the sign of $\gamma$:
\begin{itemize}
\item if $\gep(\gamma)=-$ then we prove $\cP(n,\gamma)$ assuming that $\cP(m,\gamma)$ holds for all $m\le n-1$, for all $n\ge 1$,
\item if $\gep(\gamma)=+$ then we prove $\cP(n,\gamma)$ assuming that $\cP(m,\gamma)$ holds for all $m\ge n+1$, for all $n\ge 0$.
\end{itemize}
The descending induction for positive contours works for positive contours because we already know from Proposition \ref{lesmall} that $\cP(\gamma,n)$ holds for
$n\ge \gb^{-1} \sqrt{\Diam(\gamma)}$. The ascending induction for negative contours is initiated for $n=1$ (for which the induction hypothesis
``$\cP(m,\gamma)$ holds for all $m\le 0$'' is empty).

In the remainder of the proof we always assume that $n\ge 0$ for $\gamma$ positive and $n\ge 1$ for $\gamma$ negative.
For readability we also adopt the following convention within the proof%
\begin{equation}\label{uvalu}
u=u(n,\gamma):=\begin{cases}
u^+_{n+1}=200 J^{n+3} \quad & \text{ if } \gep(\gamma)=+,\\
u^-_{n}=\frac{1}{200} J^{n+2} \quad & \text{ if } \gep(\gamma)=-.
\end{cases}
\end{equation}
By Lemma \ref{mozon} it is indeed sufficient to check stability for this value of $u$.
It turns out that the ratio $\bar z^{u}_{n+\gep(\gamma)}/\bar z^{u}_n$ is easier to work with than the quantity $z^{u}_{n+\gep(\gamma)}/\bar z^{u}_n$ which appears in the definition of $w^u_n$.
Thus our first task is to prove the following estimate.

\begin{lemma}\label{reduc}
If $\cP(n+1,\gamma)$ holds then we have, for $u$ defined in \eqref{uvalu}
\begin{equation*}
w^u_n(\gamma)\le 2e^{-\gb |\tilde \gamma|}\,\frac{\bar z^{u}_{n+\gep(\gamma)}(\gamma)}{\bar z^n_u(\gamma)}.
\end{equation*}
\end{lemma}
A consequence of this result is that to prove the $n$-stability of $\gamma$ we only need to show that
\begin{equation}\label{core}
\frac{\bar z^{u}_{n+\gep(\gamma)}(\gamma)}{\bar z^{u}_n(\gamma)}\le \frac{1}{2}\,e^{|\tilde \gamma|}.
\end{equation}

\begin{proof}[Proof of Lemma \ref{reduc}]
Let us assume for simplicity that $\gamma$ is a positive contour (the adaptation for the negative case is straightforward).
From the definition \eqref{defoom} of $\gO^+[\gamma,n]$ and $\bar \gO^+[\gamma,n]$ we have
\[
z^u_{n+1}(\gamma)=\bar z^u_{n+1}(\gamma)+e^{-\gb |\tilde \gamma|} z^u_{n+2}(\gamma).
\]
Thus using $\cP(n+1,\gamma)$ and the fact that $u_{n+1}^+\ge u^+_{n+2}$, we have
\begin{equation*}
w^u_n(\gamma)=e^{-\gb |\tilde \gamma|}\, \frac{\bar z^{u}_{n+1}}{\bar z^n_u}(\gamma) \bigl[ 1+ w^u_{n+1}(\gamma)\bigr]
\le \bigl(1+e^{-(\gb-1)|\tilde \gamma|}\bigr)e^{-\gb |\tilde \gamma|}\,\frac{\bar z^{u}_{n+1}}{\bar z^n_u}(\gamma).\qedhere
\end{equation*}
\end{proof}

The strategy to prove \eqref{core} is to decompose $\bar z^u_{n+\gep(\gamma)}(\gamma)$ according to the set
of large external contours present in the field $\phi \in \bar \gO(n+\gep(\gamma),\gamma)$.
Here and in the remainder of the proof, large means $n$-large (diameter larger than $\gb^2n^2$). Let us introduce some notation to perform this decomposition.
Let $\cK^{\larg}_{\ext}(\gamma,n)$ be the set of compatible collections of large external contours and $\cK^{\sm}_{\ext}(\gamma,n)$ be
the set of compatible collections of small external contours (recall \eqref{daext}):
\begin{equation}\label{smallandbig}
\begin{split}
\cK^{\larg}_{\ext}(\gamma,n)&:= \{ \gG_1 \in \bar \cK_{\ext}(\gamma) \sep \forall \gamma_1\in \gG_1,\ \Diam(\gamma_1)>n^2\gb^2\},\\
\cK^{\sm}_{\ext}(\gamma,n)&:= \{ \gG_2 \in \bar \cK_{\ext}(\gamma) \sep \forall \gamma_2\in \gG_2,\ \Diam(\gamma_2)\le n^2\gb^2\},
\end{split}
\end{equation}
where $n^2\gb^2$ is replaced by $\gb^2$ when $n=0$.
We let $\cK^{\larg,+}_{\ext}(\gamma,n)$ and $\cK^{\sm,+}_{\ext}(\gamma,n)$ the subsets of $\cK^{\larg}_{\ext}(\gamma,n)$ and $\cK^{\sm}_{\ext}(\gamma,n)$ respectively
which contains only positive contours.

For $\gG_1\in \bar \cK_{\ext}(\gamma)$, we let $\bar \gG_1$ and $L(\gG_1)$ denote respectively the set of $\bbZ^2$ sites
enclosed by contours in $\gG_1$ and the total length
of the contours in $\gG_1$
\begin{equation}\label{interiors}
\bar \gG_1:=\bigcup_{\gamma_1\in \gG_1} \bar \gamma_1, \quad \text{ and }\quad L(\gG_1):=\sum_{\gamma_1\in \gG_1} |\tilde \gamma_1|.
\end{equation}
Finally, we define $Z^{u,\sm}_{m}[\gamma,\gG_1,n]$ which corresponds to a partition function on the domain
$\bar \gamma \setminus \bar \gG_1$,
which displays only $n$-small contours which are
compatible with $\gG_1\cup\nobreak \{\gamma\}$,%
\begin{equation*}%\label{defrelou}
Z^{u,\sm}_{m}[\gamma,\gG_1,n]:= \sum_{ \{\gG_2 \in \cK^{\sm}_{\ext}(\gamma,n) \sep \gG_2 \parallel \gG_1\}} \prod_{\gamma_2\in \gG_2 } e^{-\gb |\tilde\gamma_2|}
z^u_{m+\gep(\gamma_2)}(\gamma_2).
\end{equation*}
Note that when $m=0$, the contribution of $\gG_2$ is non-zero only if $\gG_2 \in \cK^{\sm,+}_{\ext}(\gamma,n)$.
The aim of our decomposition procedure is to prove the following result.

\begin{lemma}\label{decopoz}
Assuming that $\cP(\gamma')$ holds when $\bar \gamma' \subset \bar \gamma$, $\bar \gamma'\ne \bar \gamma$,
we have, for $u$ defined in \eqref{uvalu}
\begin{equation}\label{bluz}
\frac{\bar z^u_{n+\gep(\gamma)}(\gamma)}{\bar z^u_{n}(\gamma)}\le \sum_{\gG_1\in \cK^{\larg}_{\ext}(\gamma,n)}
e^{-(\gb-2) L(\gG_1) } \frac{Z^{u,\sm}_{n+\gep(\gamma)}[\gamma,\gG_1,n]}{Z^{u,\sm}_{n}[\gamma,\gG_1,n]}.
\end{equation}
\end{lemma}

To conclude the proof of the result, we need two technical estimates to control the sum in the r.h.s.\ of \eqref{bluz}.
The first one makes it possible to bounds the ratio
\[
(Z^{u,\sm}_{n+\gep(\gamma)}/ Z^{u,\sm}_{n})[\gamma,\gG_1,n]
\]
by a simpler quantity for which one can have a
geometric intuition. It is proved in Section \ref{dcimo}.
\begin{proposition}\label{difficilissimo}
For $\gb$ sufficiently large, we have, for any $\gG_1\in \cK^{\larg}_{\ext}(\gamma,n)$ and~$u$ defined in \eqref{uvalu}
\begin{equation}\label{cookries}
\frac{Z^{u,\sm}_{n+\gep(\gamma)}[\gamma,\Gamma_1,n]}{Z^{u,\sm}_{n}[\gamma,\Gamma_1,n]}\le \frac{1}{2}
\exp\left(-J^{3n+3}|\bar \gamma\setminus \bar \gG_1|+(L(\gG_1)+|\tilde \gamma|)\right).
\end{equation}
\end{proposition}
The second estimate which leads to the conclusion is a control of the simplified sum.
We prove it in Section \ref{dcimus}.
\begin{proposition}\label{dificilissimus}
We have for $\gb$ sufficiently large for every $n\ge 0$
\begin{equation}\label{crookies}
\sum_{\gG_1\in \cK^{\larg}_{\ext}(\gamma,n)} \exp\left(-J^{3n+3}|\bar \gamma\setminus \bar \gG_1|-(\gb-3) L(\gG_1) \right)\le 1.
\end{equation}
\end{proposition}
Combining \eqref{bluz}, \eqref{cookries}and \eqref{crookies}, we deduce that
\begin{equation*}
\frac{\bar z^{u}_{n+\gep(\gamma)}(\gamma)}{\bar z^{u}_n(\gamma)}\le \frac{1}{2}\,e^{|\tilde \gamma|}
\sum_{\gG_1\in \cK^{\larg}_{\ext}(\gamma)} e^{-J^{3n+3} |\bar \gamma\setminus \bar \Gamma_1|+(\gb-3)L(\gG_1)}\le \frac{1}{2}\, e^{|\tilde \gamma|},
\end{equation*}
which ends our proof by induction (\cf \eqref{core}).

\subsection{Proof of Lemma \ref{decopoz}}

We split our reasoning into two lemmas, one providing an upper bound on $\bar z^{u}_{n+\gep(\gamma)}(\gamma)$ and the other providing a lower bound on
$\bar z^{u}_{n}(\gamma)$.

\begin{lemma}\label{plif}
Assuming that $\cP(\gamma')$ holds whenever $\bar \gamma'\subset \gamma$ and $\bar \gamma'\ne \bar \gamma$,
for $u$ defined in \eqref{uvalu} we have, when $n+\gep(\gamma)\ge 1$
\begin{equation}\label{floz}
\bar z^{u}_{n+\gep(\gamma)}(\gamma)\le \sum_{\gG_1\in \cK^{\larg}_{\ext}(\gamma)}
\biggl(\prod_{\gamma_1\in \gG_1} e^{- (\gb-2) |\tilde \gamma_1|} z^u_{n}(\gamma_1) \biggr)
Z^{u,\sm}_{n+\gep(\gamma)}[\gamma, \gG_1,n].
\end{equation}
Furthermore for $n=1$, $\gep(\gamma)=-1$ we have
\begin{equation*}
\bar z^{u}_{0}(\gamma)= \sum_{\gG_1\in \cK^{\larg,+}_{\ext}(\gamma)}
\biggl(\prod_{\gamma_1\in \gG_1} e^{- \gb |\tilde \gamma_1|} z^u_{n}(\gamma_1) \biggr)
Z^{u,\sm}_{n+\gep(\gamma)}[\gamma, \gG_1,n].
\end{equation*}
\end{lemma}
To conclude the proof of \eqref{bluz} we also need a lower bound for $\bar z^u_n(\gamma)$, which is provided by the following lemma.

\begin{lemma}\label{plof}
For any $\gG_1\in \cK^{\larg}_{\ext}(\gamma)$, we have
\begin{equation}\label{flaz}
\bar z^{u}_{n}(\gamma)\ge \prod_{\gamma_1\in \gG_1} z^u_{n}(\gamma_1)
Z^{u,\sm}_{n}[\gamma, \gG_1,n].
\end{equation}
\end{lemma}
The inequality \eqref{bluz} is obtained combining \eqref{floz} and \eqref{flaz}.

\begin{proof}[Proof of Lemma \ref{plif}]
We assume for notational simplicity that $\gamma$ is a positive contour.
Recalling Equation \eqref{tapun} and splitting the set $\gG'$ of external contour between large ($\gG_1$) and small ($\gG_2$) contours we obtain
\begin{multline*}
\bar z^{u}_{n+1}(\gamma)= \sum_{\gG'\in \bar \cK_{\ext}(\gamma)}\prod_{\gamma'\in \gG'} e^{-\gb |\tilde \gamma'|} z^h_{n+1+\gep(\gamma')}(\gamma') \\
=
\!\!\!\!\!\sum_{\gG_1\in \cK^{\larg}_{\ext}(\gamma,n)}
\prod_{\gamma_1\in \gG_1} e^{-\gb |\tilde \gamma_1|} z^u_{n+1+\gep(\gamma_1)}(\gamma_1) \!\!\!\!\!\!\!\!
\sum_{\{ \gG_2 \in \cK^{\sm}_{\ext}(\gamma,n) \sep \gG_2 \parallel \gG_1\}} \prod_{\gamma_2\in \gG_2 } e^{-\gb |\tilde\gamma_2|} z^u_{n+\gep(\gamma_2)}(\gamma_2)\\
= \prod_{\gamma_1\in \gG_1} e^{-\gb |\tilde \gamma_1|} z^u_{n+1+\gep(\gamma_1)}(\gamma_1) Z^{u,\sm}_{n}[\gamma, \gG_1].
\end{multline*}
To conclude, we need to check that
\begin{equation*}
z^u_{n+1+\gep(\gamma_1)}(\gamma_1)\le e^{2|\tilde \gamma_1|}z^u_{n}(\gamma_1).
\end{equation*}
This is of course obvious when $\gep(\gamma_1)=-1$. For positive contours on the other hand we have
\begin{multline}\label{lezinc}
z^u_{n+2}(\gamma_1)= e^{\gb | \tilde \gamma_1|}w^u_{n+1}(\gamma_1) \bar z^u_{n+1}(\gamma_1)\\
\le
e^{\gb | \tilde \gamma_1|}w^u_{n+1}(\gamma_1) z^u_{n+1}(\gamma_1)= e^{2\gb | \tilde \gamma_1|}w^u_{n+1}(\gamma_1) w^u_{n}(\gamma_1) z^u_{n}(\gamma_1).
\end{multline}
Using the induction hypothesis, or more precisely $\cP(n+1,\gamma_1)$ (recall that $u\ge u^+_{n+2}$) and $\cP(n,\gamma_1)$, we deduce from \eqref{lezinc} that
\begin{equation*}
z^u_{n+2}(\gamma_1)\le e^{2|\tilde \gamma_1|}z^u_{n}(\gamma_1).
\end{equation*}
The same proof goes when $\gep(\gamma)=-1$, if we restrict the sum to the set of positive contours in the special case $n=1$.
\end{proof}

\begin{proof}[Proof of Lemma \ref{plof}]

Instead of proving \eqref{flaz}, we prove a stricter inequality where the contours in $\gG_2$ are not required to be small, and which is valid
for all $\gG_1\in \bar \cK_{\ext}(\gamma)$ (recall \eqref{daext})
\begin{equation}\label{fliz}
\bar z^{u}_{n}(\gamma)\ge \prod_{\gamma_1\in \gG_1} z^u_{n}(\gamma_1) \sum_{\{ \gG_2 \in \bar \cK_{\ext}(\gamma) \sep \gG_1 \parallel \gG_2\}}
\prod_{\gamma_2\in \gG_2 } e^{-\gb |\tilde\gamma_2|} z^u_{n+\gep(\gamma_2)}(\gamma_2)
\end{equation}
(in the case $n=0$ only $\gG_2$ with all contour positive give a contribution to the sum).
We shall show that the l.h.s. in \eqref{fliz} corresponds to the contribution to the sum \eqref{partisioux} of
the set of $\phi$s whose external contours are either in $\gG_1$ or compatible with $\gG_1$
\[
\cA:=\left\{ \phi\in \bar \gO^+[\gamma,n] \sep \forall \gamma' \in \Upsilon^{\ext}_n(\phi),\ \gamma' \mid \gG_1 \text{ or } \gamma'\in \gG_1\right\}.
\]
To make our decomposition we use the notation
\begin{equation*}\begin{split}
\gG_{1,1}(\gamma_1)&=\{ \gamma_{1,1}\in \Upsilon^{\ext}_n(\phi)\setminus \{\gamma_1\} \sep \bar \gamma_{1,1}\subset \bar \gamma_1\},\\
\gG_2&=\Upsilon^{\ext}_n(\phi) \setminus \bigcup_{\gamma_1\in \gG_1}\gG_{1,1}(\gamma_1).
\end{split}\end{equation*}
Note that for $\phi\in \cA$ we have $\gG_{1,1}(\gamma_1)\subset \bar \cK_{\ext}(\gamma_1)$.
In analogy with \eqref{tapun}, we can thus write
\begin{multline*}
\sum_{\phi\in \cA} e^{-\gb \cH^n(\phi)+h|\phi^{-1}(0)|}\\[-8pt]
=
\prod_{\gamma_1\in \gG_1}\biggl[ e^{-\gb|\tilde \gamma_{1}|} z^u_{n+\gep(\gamma_1)}(\gamma_1)+ \sum_{\gG_{1,1}\in \bar \cK_{\ext}(\gamma_1)}
\prod_{\gamma_{1,1}\in \gG_{1,1}} e^{-\gb|\tilde \gamma_{1,1}|}z^u_{n+\gep(\gamma_{1,1})} \biggr]\\
\times\sum_{\{ \gG_2 \in \bar\cK_{\ext}(\gamma) \sep \gG_1 \parallel \gG_2\}} \prod_{\gamma_2\in \gG_2 } e^{-\gb |\tilde\gamma_2|} z^u_{n+\gep(\gamma_2)}(\gamma_2).
\end{multline*}
In each factor of the product over $\gamma_1$, the first term corresponds to the contribution of $\phi$s
for which $\gamma_1$ is a contour.
Finally recalling Equations \eqref{relatz} and \eqref{tapun} and
\begin{multline*}
e^{-\gb|\tilde \gamma_{1}|} z^u_{n+\gep(\gamma_1)}(\gamma_1)+ \sum_{\gG_1\in\bar\cK^{\ext}(\gamma_1)} \prod_{\gamma_{1,1}\in \gG_1} e^{-\gb|\tilde \gamma_{1,1}|}
z^u_{n+\gep(\gamma_{1,1})}\\[-5pt]
=
e^{-\gb|\tilde \gamma_{1}|} z^u_{n+\gep(\gamma_1)}(\gamma_1)+ \bar z^u_n(\gamma_1)= z^{u}_n(\gamma_1),
\end{multline*}
which yields \eqref{fliz}.
\end{proof}

\subsection{Proof of Proposition \ref{difficilissimo}}\label{dcimo}

To prove the inequality \eqref{cookries}, we prove separately bounds for the numerator and for the denominator.
As for Proposition \ref{lesmall}
we have to treat separately the cases $n=1$, $\gamma$ negative and $n=0$, $\gamma$ positive, which we do in Lemma \ref{dzip}. The general case is dealt with using Lemma
\ref{dalo}.
The proof of these two results is technically involved, and for that reason, postponed to Section \ref{secrestrict}.

\begin{lemma}\label{dzip}
There exists a constant $C$ (independent of $\gb$) such that for all $\gb$ sufficiently large for every every $\gamma$ with $\Diam(\gamma)> \gb^2$, and every
$\gG_1\in \cK^{\larg}_{\ext}(\gamma,1)$,
$h\in \left[0,2J^2\right]$, we have
\begin{equation}\label{lezink}
\begin{split}
\log Z^{u,\sm}_{0}[\gamma,\Gamma_1,1]&\le|\bar \gamma \setminus \bar \Gamma_1|\left(h+J^2+2J^3+C J^4\right),\\
\log Z^{u,\sm}_{0}[\gamma,\Gamma_1,1]&\ge |\bar \gamma \setminus \bar \Gamma_1|\left(h+2J^2+2J^3-CJ^4\right)+\frac{1}{2}(|\tilde \gamma|+L(\gG_1)),\\
\log Z^{u,\sm}_{1}[\gamma,\Gamma_1,1]&\le|\bar \gamma \setminus \bar \Gamma_1|\left(2J^2+4J^3-C J^4\right),\\
\log Z^{u,\sm}_{1} [\gamma,\Gamma_1,1]&\ge |\bar \gamma \setminus \bar \Gamma_1| \left(2J^2+4J^3-CJ^4\right) -\frac{1}{2}(|\tilde \gamma|+L(\gG_1)).
\end{split}
\end{equation}
\end{lemma}

\begin{rem}
Note that the inequalities of \eqref{lezink} also hold if $[\gamma,\Gamma_1,1]$ is replaced by $[\gamma,\Gamma_1,0]$ because the associated notions of small contour are the same.
The range we have chosen for $h$ is sufficient to treat the case of $u=u^{\pm}_1$ for $\gb$ sufficiently large as for these value we have $h=J^2+O(J^3)$.
\end{rem}

To treat the other cases, we define $\gO^+[m,\gamma,\gG_1,n]$ to be the set of interface realizations naturally associated with the partition function
$Z^{u,\sm}_m[\gamma,\gG_1,n]$.
We define first $\cC[\gamma,\gG_1,n]$ the set of contours which can appear in $\Upsilon_m(\phi)$
\begin{equation*}
\cC[\gamma,\gG_1,n]:=\{ \gamma_2\in \cC \sep \bar \gamma_2\subset \bar \gamma\setminus \bar \gG_1,\ \Diam(\gamma_2)\le (n\gb)^2 \text{ and }
\gamma_2 \mid \left(\{\gamma\} \cup \gG_1\right) \}.
\end{equation*}
We set\vspace*{-7pt}%
\begin{multline*}
\gO[m,\gamma,\gG_1,n]:=\Big\{ (\phi \sep \bar \gamma \setminus \bar \gG_1 \rightarrow \bbZ) \sep \Upsilon^{\ext}_m(\phi)\subset \cC[\gamma,\gG_1,n]
\\[-5pt]
\text{ and } \phi= m+ \textstyle\sum_{\hat \gamma_2 \in \hat \Upsilon_m(\phi)} \varphi_{\hat \gamma_2}. \Big\},\vspace*{-7pt}
\end{multline*}
and as usual
\begin{equation*}
\gO^+[m,\gamma,\gG_1,n]:= \Big\{ \phi\in \gO[m,\gamma,\gG_1,n] \sep \forall x\in \bar \gamma \setminus \bar \gG_1,\; \phi(x)\ge 0 \Big\}.
\end{equation*}
The condition $\phi= m+ \sum_{\hat \gamma_2 \in \hat \Upsilon_m(\phi)} \varphi_{\hat \gamma_2}$ corresponds to \eqref{cylinder}, and is violated when $\phi$ presents some level lines which surround holes in $\bar \gamma \setminus \bar \gG_1$.
With this definition, the reader can check that (recall our convention $h=u+h_w(\gb)$)
\begin{equation}\label{winouze}
Z^{u,\sm}_{m}[\gamma,\Gamma_1,n]= \sum_{\phi \in \gO^+[m,\gamma,\gG_1,n]}e^{-\gb \cH^m_{\bar \gamma \setminus \bar \gG_1}(\phi)+h|\phi^{-1}\{0\}|}.
\end{equation}
We let $\bP^{m,\sm,n}_{\gamma,\gG_1}$ be the SOS measure restricted to $\gO[m,\gamma,\gG_1,n]$
\begin{equation*}
\bP^{n,\sm}_{\gamma,\gG_1}(\phi):=\frac{1}{Z^{\sm}[\gamma,\gG_1,n]}e^{-\gb \cH^m_{\bar \gamma \setminus \bar \gG_1}(\phi)},
\end{equation*}
where
\begin{equation*}
Z^{\sm}[\gamma,\gG_1,n]:= \sum_{\phi \in \gO[m,\gamma,\gG_1,n]} e^{-\gb \cH^m_{\bar \gamma \setminus \bar \gG_1}(\phi)}
\end{equation*}
(again by translation invariance, the partition function does not depend on the boundary condition $m$).
We state a result which is similar to Lemma \ref{rwo} and is useful in our proofs.

\begin{lemma}\label{zilak}
For any $m\ge 1$ and any $\gamma\in \cC$ and $\gG_1\in \cK^{\larg}_{\ext}(\gamma)$
\begin{equation*}%\label{groot}
Z^{u,\sm}_{m}[\gamma,\Gamma_1,n]= Z^{\sm}[\gamma,\gG_1,n] \bE^{m,\sm,n}_{\gamma,\gG_1}\bigl[ e^{u |\phi^{-1}(\bbZ_-)|- \bar H(\phi^{-1})(\bbZ_-)}\bigr].
\end{equation*}
\end{lemma}
\begin{proof}
We have to show that
\begin{equation*}
Z^{u,\sm}_{m}[\gamma,\Gamma_1,n]= \sum_{\phi \in \gO^+[m,\gamma,\gG_1,n]}e^{-\gb \cH^m_{\bar \gamma \setminus \bar \gG_1}(\phi)+u|\phi^{-1}(\bbZ_-)|-\bar H(\phi^{-1}(\bbZ_-))}.
\end{equation*}
The proof can be adapted from that of \cite[Lem.\,3.1]{cf:part1}. The sum over all the possible options for
the negative parts of $\phi$ cancels the term $\bar H$ and changes $u$ into $h$, so that one recovers \eqref{winouze}.
The key observation to check that the proof adapts is that
the contour restriction does not bring any constraint on the choice of $\phi_-=\max(0,-\phi)$ once $\phi^{-1}(\bbZ_-)$ is fixed.
This is the case
because the contour restriction forces the diameter of maximal connected components of $\phi^{-1}(\bbZ_-)$
to be smaller than $(\gb n)^2$.
\end{proof}

As a consequence of Lemma \ref{zilak},
when neither $n$ nor $n+\gep(\gamma)$ are zero, the $\log$ of the estimated ratio can be rewritten in the following form
\begin{multline}\label{crounch}
\log \biggl(\frac{ Z^{u,\sm}_{n+\gep(\gamma)}[\gamma,\Gamma_1,n]}{ Z^{u,\sm}_{n}[\gamma,\Gamma_1,n]}\biggr)=
\log \bE^{n+\gep(\gamma),\sm,n}_{\gamma,\gG_1}\bigl[ e^{u |\phi^{-1}(\bbZ_-)|- \bar H(\phi^{-1})(\bbZ_-)} \bigr] \\[-3pt]
- \log
\bE^{n,\sm,n}_{\gamma,\gG_1}\bigl[ e^{u |\phi^{-1}(\bbZ_-)|- \bar H(\phi^{-1})(\bbZ_-)} \bigr].
\end{multline}
We need the following statements
\begin{lemma}\label{dalo}
The following estimates hold:
\begin{enumeratei}
\item\label{dalo1}
For positive $\gamma$ and $u$ as in \eqref{uvalu}
\begin{equation*}
\log \bE^{n+1,\sm,n}_{\gamma,\gG_1}\bigl[ e^{u |\phi^{-1}(\bbZ_-)|- \bar H(\phi^{-1})(\bbZ_-)} \bigr]\le |\bar \gamma \setminus \bar \Gamma_1| 2u J^{2n+2}.
\end{equation*}
\item\label{dalo2}
For arbitrary $\gamma$ and $u$ as in \eqref{uvalu}
\begin{multline}\label{ckonveu}
\log \bE^{n,\sm, n}_{\gamma,\gG_1}\bigl[ e^{u |\phi^{-1}(\bbZ_-)|- \bar H(\phi^{-1})(\bbZ_-)} \bigr]\\
\ge |\bar \gamma \setminus \bar \Gamma_1|
\Bigl(\frac{1}{2}u J^{2n}- 40 J^{3n+3}\Bigr)-\frac{1}{4}\left(|\tilde \gamma|+L(\gG_1) \right).
\end{multline}
\item\label{dalo3}
For negative $\gamma$ and $u$ as in \eqref{uvalu}
\begin{multline*}%
\log \bE^{n-1,\sm, n}_{\gamma,\gG_1}\bigl[ e^{u |\phi^{-1}(\bbZ_-)|- \bar H(\phi^{-1})(\bbZ_-)} \bigr]\\[-5pt]
\le |\bar \gamma \setminus \bar \Gamma_1|
\Bigl(4u J^{2(n-1)}-\frac{1}{4}J^{3n}\Bigr)+\frac{1}{4}\left(|\tilde \gamma|+L(\gG_1) \right).
\end{multline*}
\end{enumeratei}
\end{lemma}

\begin{proof}[Proof of Proposition \ref{difficilissimo}].
We start with the case of positive contour with $n=0$.
Using Lemma \ref{dzip} we have
\begin{equation*}
\log \biggl(\frac{ Z^{u,\sm}_{1}[\gamma,\Gamma_1,1]}{ Z^{u,\sm}_{0}[\gamma,\Gamma_1,1]}\biggr)
\le |\bar \gamma \setminus \bar \gG_1|(J^2+2J^3+2CJ^4-h)+ \frac 1 2\left(|\tilde \gamma|+L(\gG_1)\right).
\end{equation*}
Recall now (\cf \eqref{uvalu}) that we are interested in the case
\begin{equation*}
h=h_w(\gb)+u^+_1= \log\Bigl(\frac{ e^{4\gb} }{e^{4\gb}-1}\Bigr)+200J^3,
\end{equation*}
hence we obtain
\begin{equation*}
\log \biggl(\frac{ Z^{u,\sm}_{1}[\gamma,\Gamma_1,1]}{ Z^{u,\sm}_{0}[\gamma,\Gamma_1,1]}\biggr)
\le -|\bar \gamma \setminus \bar \gG_1| J^3+\frac 1 2 \left(|\tilde \gamma|+L(\gG_1)\right).
\end{equation*}
We let the reader check that similarly for negative contours and $u=u^-_1$ we have
\begin{equation*}
\log \biggl(\frac{ Z^{u,\sm}_{0}[\gamma,\Gamma_1,1]}{ Z^{u,\sm}_{1}[\gamma,\Gamma_1,1]}\biggr)
\le -|\bar \gamma \setminus \bar \gG_1| J^3+\frac{1}{2}\left(|\tilde \gamma|+L(\gG_1)\right).
\end{equation*}
Let us now treat the case of a positive contour for $n\ge 1$.
Using \eqref{crounch} and Lemma~\ref{dalo}, we have for $u=u^+_{n+1}$ (recall\eqref{uvalu})\vspace*{-5pt}%
\begin{multline*}
\log \biggl(\frac{ Z^{u,\sm}_{n+1}[\gamma,\Gamma_1,n]}{ Z^{u,\sm}_{n}[\gamma,\Gamma_1,n]}\biggr)\\
\le |\bar \gamma \setminus \Gamma_1|\Bigl[ u\Bigl(2J^{2n+2}- \frac{1}{2}J^{2n}\Bigr)+40J^{3n+3} \Bigr]+ \frac{1}{2}(|\tilde \gamma|+L(\gG_1))\\
\le |\bar \gamma \setminus \Gamma_1|(40J^{3n+3}-\frac{1}{4}u J^{2n})+ \frac{1}{2}(|\tilde \gamma|+L(\gG_1)).
\end{multline*}
Recalling now \eqref{uvalu},
we obtain the result by observing that
\begin{equation*}
40J^{3n+3}-\frac{1}{4}u^+_{n+1} J^{2n}\le -J^{3n+3}.
\end{equation*}
In a similar manner in the case of negative contour and $n\ge 2$ we have as a consequence of Lemma \ref{dalo} \eqref{dalo2}--\eqref{dalo3}, for $\gb$ sufficiently large,
\begin{equation*}
\log \biggl(\frac{ Z^{u,\sm}_{n-1}[\gamma,\Gamma_1,n]}{ Z^{u,\sm}_{n}[\gamma,\Gamma_1,n]}\biggr)
\le
|\bar \gamma \setminus \Gamma_1| \Bigl(4u J^{2(n-1)}-\frac{1}{5}J^{3n}\Bigr)+ \frac{1}{2}(|\tilde \gamma|+L(\gG_1))
\end{equation*}
and we conclude by observing that
\begin{equation*}
4u^-_n J^{2(n-1)}-\frac{1}{5}J^{3n}\le -J^{3n+3},
\end{equation*}
so that \eqref{cookries} is satisfied in all cases.
\end{proof}

\subsection{Proof of Proposition \ref{dificilissimus}}\label{dcimus}

We can relax for this proof the notion of compatibility, meaning we consider the sum over a superset of $\cK^{\larg}_{\ext}(\gamma,n)$.
We consider, in this section only, that two contours are externally compatible if $\bar \gamma \cap \bar \gamma'=\emptyset$.
Adding a factor $2$ to take the sign into account (that is replacing each factor $e^{-(\gb-3)|\tilde \gamma_1|}$ by $2e^{-(\gb-3)|\tilde \gamma_1|}$ in \eqref{crookies}),
we choose to consider geometric contours instead of signed contours (and use $\bar \gamma_1$ to denote the interior of $\tilde \gamma_1$).

Our proof works by induction and leads us to consider sets of external contour in a general domain $\gL\subset \bbZ^2$
which are not necessarily simply connected.
We use, in this section only, the notation $\gG$ and $\gamma$ instead of $\gG_1$ and $\gamma_1$.
We maintain that all contours must satisfy $\bar \gamma\subset \gL$ and thus cannot surround holes.

We let $\tilde \cK^{\larg}_{\ext}(\gL,n)$ denote the set of collections of externally compatible $n$-large geometric contours
with the above mentioned notion of compatibility.
The result \eqref{crookies} will follow (provided that $e^{-4(\gb/2)}\ge 2e^{-4(\gb-3)}$) if we can prove that for every $\gL\subset \bbZ^2$
\begin{equation}\label{wookies}
\sum_{\tilde \gG\in \tilde \cK^{\larg}_{\ext}(\gL)} e^{-J^{3n+3} |\gL\setminus \bar \gG|-(\gb/2) L(\tilde \Gamma)}\le 1,
\end{equation}
where $L(\tilde \gG)$ and $\bar \gG$ are the length and area associated with $\tilde \gG$
defined in analogy with \eqref{interiors}.
We prove a more general version of the statement.
\begin{proposition}
For any finite domain in $\gL$ and any $\ell\ge 4$ we have
\begin{equation}\label{dunnies}
\sum_{\tilde \gG\in \tilde \cK^{\ell+}_{\ext}(\gL)} e^{-2^{-\ell} |\gL\setminus \bar \gG| }e^{-(\gb/2) L(\tilde\gG)}\le 1,
\end{equation}
where $\tilde \cK^{\ell+}_{\ext}(\gL)$ denote the set of collections of externally compatible geometric contours with length larger than $2\ell$.
\end{proposition}
If we apply this proposition for $\ell=n^2\gb^2$, \eqref{dunnies} implies \eqref{wookies} provided that
$2^{-n^2\gb^2}\le J^{3n+3}$, which is valid for every $n\ge 1$ provided that $\gb$ is sufficiently large (we have also $2^{-\gb^2}\le J^{3}$ to cover the case $n=0$).

\begin{proof}
We prove the result by induction on $\ell$. We let $\tilde \cK^{\ell}_{\ext}(\gL)$ be
the set of collection externally compatible of geometric contours with length equal to $2\ell$.
The key step consists in proving that
\begin{equation}\label{industaps}
\sum_{\tilde \gG\in \tilde \cK^{\ell}_{\ext}(\gL)} e^{-2^{-\ell} |\gL\setminus \bar \gG|-(\gb/2) L(\tilde\gG)}\le e^{-2^{-(\ell+1)} |\gL\setminus \bar \gG|}.
\end{equation}
Let us show how \eqref{dunnies} is deduced from \eqref{industaps}.
First let us observe that \eqref{dunnies} is obviously satisfied when $\ell$ is larger than the total number of edges in $\gL$. Hence we can proceed
by descending induction, assuming that the
statement is valid for $\ell+1$ and proving it for $\ell$.

Obviously $\tilde \gG \in \tilde \cK^{\ell+}_{\ext}(\gL)$ can be written in the form $\tilde \gG_1\cup \tilde \gG_2$ where $\tilde \gG_1\in \tilde \cK^{(\ell+1)+}_{\ext}(\gL)$
and $\tilde \gG_2\in \tilde\cK^{\ell}_{\ext}(\gL \setminus \bar \gG_1)$.
Hence\vspace*{-3pt}
\begin{multline*}
\sum_{\tilde \gG\in \tilde \cK^{\ell+}_{\ext}(\gL)} e^{-2^{-\ell} |\gL\setminus \bar \gG| }e^{-(\gb/2) L(\tilde\gG_1)}\\[-5pt]
=\sum_{\tilde \gG_1\in \tilde \cK^{(\ell+1)+}_{\ext}(\gL)} e^{-(\gb/2) L(\tilde\gG_1)}\sum_{\tilde \gG_2\in \tilde \cK^{\ell}_{\ext}(\gL\setminus \bar \gG_1)}
e^{-2^{-\ell} |(\gL\setminus \bar \gG_1)\setminus \bar \gG_2|} e^{-(\gb/2) L(\tilde\gG_2)}
\\ \le \sum_{\tilde \gG_1\in \tilde \cK^{(\ell+1)+}_{\ext}(\gL)} e^{-(\gb/2) |\tilde\gG_1|} e^{-2^{-(\ell+1)} |\gL\setminus \bar \gG_1|} \le 1,
\end{multline*}
where in the first inequality we use \eqref{industaps} for the domain $\gL\setminus \bar \gG_1$ and in the second one the induction hypothesis.

Let us now prove \eqref{industaps}.
We have to distinguish between two sorts of contributions, according to the number of contours.
Let us first consider the contribution where the number of contours is smaller than $m:=\lfloor |\gL|\ell^{-2} \rfloor$.
Keeping in mind that, from isoperimetric inequalities, a contour encloses at most $(\ell/2)^2$ sites, we have
\[
|\gL\setminus \bar \gG|\ge |\gL|-(\ell/2)^2m \ge \frac{3}{4}|\gL|
\]
and hence
\begin{equation*}
\sum_{\{ \tilde \gG\in \tilde \cK^{\ell}_{\ext}(\gL)\sep \# \tilde \gG \le m \}} e^{-2^{-\ell} |\gL\setminus \bar \gG|-(\gb/2) L(\tilde\gG)}
\le \exp\Bigl(-\frac{3}{2} 2^{-(\ell+1)}|\gL|\Bigr) \sum_{\tilde \gG\in \tilde \cK^{\ell}_{\ext}(\gL)}e^{-(\gb/2)L(\tilde\gG)}.
\end{equation*}
Now, for each site $x\in \bbZ^2$, we let $\tilde\cC^{\ell}(x,\gL)$ denote the set of geometric contour longer than $2\ell$ such that
$\bar \gamma \subset \gL$, for which $x$ is the
smallest vertex in $\bar \gamma$ for the lexicographical order. We write $\tilde\cC^{\ell}(x)$ for the set corresponding to $\gL=\bbZ^2$.
We have
\begin{equation*}
\sum_{\tilde \gG\in \tilde \cK^{\ell}_{\ext}(\gL)}e^{-(\gb/2) L(\tilde\gG)}
\le \prod_{x\in \gL} \biggl(1+\sum_{\tilde \gamma \in \tilde\cC^{\ell}(x,\gL)} e^{-\gb\ell}\biggr) \le
\bigl(1+\#\tilde\cC^{\ell}({\bf 0}) e^{-\gb\ell}\bigr)^{|\gL|}.
\end{equation*}
where the first inequality is obtained by summing over all collections of contours instead of externally compatible ones and
the second one by extending the sum to $\tilde\cC^{\ell}(x)$ and using translation invariance.
Using the fact that, by a standard counting argument, we have that, for $\gb$ sufficiently large,
\[
\#\tilde\cC^{\ell}({\bf 0}) e^{-\gb\ell}\le 9^{\ell} e^{-\gb\ell}\le e^{-\gb\ell /2},
\]
which implies in particular that, provided $\gb$ is sufficiently large,
\begin{multline}\label{zapi}
\sum_{\{ \tilde \gG\in \tilde \cK^{\ell}_{\ext}(\gL)\sep \# \tilde \gG \le m \}} e^{-2^{-\ell} |\gL\setminus \bar \gG|}e^{-(\gb/2) L(\tilde\gG)}
\\[-5pt] \le \exp\Bigl(|\gL|\Bigl(e^{-(\gb /2) \ell}-\frac{3}{2} 2^{-(\ell+1)}\Bigr) \Bigr)\le \exp\Bigl(-\frac{5}{4}2^{-(\ell+1)} |\gL|\Bigr).
\end{multline}
Now, concerning the contribution of collections of cardinality larger than $m$, we neglect the penalty for uncovered area
\begin{equation*}
\sum_{\{ \tilde \gG\in \tilde \cK^{\ell}_{\ext}(\gL)\sep \# \tilde \gG > m \}} e^{-2^{-\ell} |\gL\setminus \bar \gG|}e^{-(\gb/2) L(\tilde\gG)}
\le \sum_{\{ \tilde \gG\in \tilde \cK^{\ell}_{\ext}(\gL)\sep \# \tilde \gG > m \}} e^{-(\gb/2) L(\tilde\gG)}.
\end{equation*}

To estimate the sum in the r.h.s.\, we consider that if $\# \tilde \gG=k$ then to select $k$ contours, we must first choose
$k$ vertices to be the minimal (for the lexicographical order) vertices enclosed by each contour
(there are $\binom{|\gL|}{k}\le \left(\sfrac{e |\gL|}{k}\right)^k$ ways to do this)
and then ignoring further compatibility conditions, consider that for each vertex, there are at most $9^{\ell}$ eligible contours of length $\ell$.
Thus we obtain that for $\gb$ sufficiently large we have
\begin{multline} \label{zapa}
\sum_{\{ \tilde \gG\in \tilde \cK^{\ell}_{\ext}(\gL)\sep \# \tilde \gG > m \}} e^{-(\gb/2) L(\tilde\gG)}\le
\sum_{k> m} e^{-\gb k\ell} \binom{|\gL|}{k} 9^{k\ell}
\le \sum_{k> m} e^{-\sfrac{\gb k\ell}{2}} \left(\sfrac{e |\gL|}{k}\right)^k\\
\le \sum_{k> m}\bigl(e^{-\Psfrac{\gb \ell}{2}+1} \ell^2 \bigr)^k
\le \sum_{k> m} e^{-\sfrac{\gb k\ell}{4}}\le 2 e^{-\sfrac{\gb \ell(m+1)}{4}} \le 2e^{-\Psfrac{\gb}{4}\min \left(\ell, |\gL|/\ell \right)}.
\end{multline}
Overall, combining \eqref{zapi} and \eqref{zapa} we obtain that
\begin{multline*}
\sum_{\tilde \gG\in \tilde \cK^{\ell}_{\ext}(\gL)}e^{-(\gb/2) L(\tilde\gG)-2^{-\ell} |\gL\setminus \bar \gG|}\\[-10pt]
\le
\exp\Bigl(-\frac{5}{4}2^{-(\ell+1)} |\gL|\Bigr)\ind_{\{m\ge 1\}}
+ 2e^{-\Psfrac{\gb}{4}\min \left(\ell, |\gL|/\ell \right)} \le e^{-2^{-(\ell+1)} |\gL|},
\end{multline*}
where to check the last inequality we have to check separately the cases $m\ge 1$ and $m=0$.
\end{proof}

\subsection{The large contour case for Proposition \ref{mejor}}\label{mejorplus}

We conclude this section by extending the stability result on a larger interval.
This can be treated in a relatively simple fashion by induction if we rely on cluster expansion estimates (Lemma \ref{finalfrontier}).

We only have to check the stability for large contours since that of small contours has been checked in Proposition \ref{lesmall}.
The proof works using the same induction as for Proposition \ref{truta}.
We define the property
\begin{equation*}
\cP(n,\gamma):=\begin{cases}
\big[ \text{the contour } \gamma \text{ is stable at level } n \text{ for } u\ge u^*_{n+1} \big] &\text{if } \gep(\gamma)=+,\\
\big[ \text{the contour } \gamma \text{ is stable at level } n \text{ for } u\le u^*_{n} \big] &\text{if } \gep(\gamma)=-.\\
\end{cases}
\end{equation*}
and
\begin{equation*}
\cP(\gamma):=\big[\cP(n,\gamma) \text{ is satisfied for all } n\ge 0 \big].
\end{equation*}
As for Proposition \ref{truta},
we need to prove $\cP(n,\gamma)$ assuming $\cP(n+\gep(\gamma),\gamma)$ and $\cP(\gamma')$ for $\bar \gamma'\subset \bar\gamma$, $\bar \gamma'\ne \bar \gamma$.
After fixing $n$ and $\gamma$ we assume below that
\begin{equation}\label{redevu}
u=u(n,\gamma):= \begin{cases} u^*_{n+1} &\text{if } \gep(\gamma)=+,\\
u^*_{n} &\text{if } \gep(\gamma)=-.
\end{cases}
\end{equation}
Using Lemma \ref{reduc} (or rather its proof) we can reduce ourselves to proving
\begin{equation*}
\frac{\bar z^u_{n+\gep(\gamma)}}{\bar z^u_n}(\gamma)\le \frac{1}{2}\,e^{|\tilde\gamma|}.
\end{equation*}
By induction hypothesis, all contours involved in the partition function $\bar z^u_{n+\gep(\gamma)}(\gamma)$, $\bar z^u_n(\gamma)$
are stable for $u$ as in \eqref{redevu}.
Recalling the definition for the truncated potential \eqref{trunk} and the definitions of Section \ref{fffa},
this stability implies that
\begin{equation*}
\bar z^u_{n+\gep(\gamma)}(\gamma)=Z\bigl[\gamma, w^{u,\tr}_{n+\gep(\gamma)}\bigr] \quad\text{and}\quad \bar z^u_{n}(\gamma)=Z\bigl[\gamma, w^{u,\tr}_{n}\bigr].
\end{equation*}
As the truncated potentials satisfy \eqref{usualchoice}, Lemma \ref{finalfrontier} enables us to deduce that
\begin{equation*}
\begin{split}
|\log \bar z^u_{n+\gep(\gamma)}(\gamma)- |\bar \gamma| \tf^{\tr}_{n+\gep(\gamma)}(\gb,u) |&\le |\tilde \gamma|/4,\\
|\log \bar z^u_{n}(\gamma)- |\bar \gamma| \tf^{\tr}_{n}(\gb,u) |&\le |\tilde \gamma|/4.
\end{split}
\end{equation*}
Using the definition of $u^*_{n+1}$ or $u^*_n$ (depending on the sign of the contour)
we have $\tf^{\tr}_{n+\gep(\gamma)}(\gb,u)=\tf^{\tr}_{n}(\gb,u)$ and thus we can deduce that
\begin{equation*}
\frac{ \bar z^u_{n+\gep(\gamma)}}{\bar z^u_{n}}(\gamma)\le \exp\left(|\tilde \gamma|/2\right).
\end{equation*}

\section{Estimates for restricted partition functions: Proof of Lemma \ref{dzip} and~\ref{dalo}}\label{secrestrict}

In this section we prove the two remaining technical lemmas from Section \ref{dcimo}, and fully complete the proof of Theorem \ref{converteo}.

\subsection{Proof of Lemma \ref{dzip}}
We first prove upper bound results which are easier. The idea is to write a contour decomposition and to relax the compatibility assumption in the sum to obtain an upper bound.
Let us start with the case of zero boundary condition.
Recall that
\begin{equation*}%\label{sumi}
Z^{u,\sm}_{0}[\gamma,\gG_1,1]:= \sum_{\{ \gG_2 \in \cK^{\sm,+}_{\ext}(\gamma,1) \sep \gG_2 \parallel \gG_1\}} e^{h|\bar \gamma \setminus \bar \gG_1\cup \bar \gG_2|} \prod_{\gamma_2\in \gG_2 }
e^{-\gb |\tilde\gamma_2|}
z^u_{1}(\gamma_2),
\end{equation*}
where (recall \eqref{smallandbig})
\[
\cK^{\sm,+}_{\ext}(\gamma,1)=\{ \gG_2\in \cK^{\sm}_{\ext}(\gamma,1) \sep \forall \gamma\in \gG_2, \gep(\gamma_2)=+\}.
\]

To obtain an upper bound, we replace the first exponential term by $e^{h|\bar \gamma \setminus \bar \gG_1|}$, and we let the sum range over all collections of small
positive $1$-small contours
$\cC^{\sm,+,1}_{\bar \gamma\setminus \bar \gG_1}$ without imposing any compatibility condition.
The sum factorizes and we obtain
\begin{equation*}
Z^{u,\sm}_{0}[\gamma,\gG_1,1]\le
e^{h|\bar \gamma \setminus \bar \gG_1|} \prod_{\gamma_2 \in \cC^{\sm,+,1}_{\bar \gamma\setminus \bar \gG_1}}
\bigl(1+ e^{-\gb |\tilde\gamma_2|}
z^u_{1}(\gamma_2) \bigr).
\end{equation*}
Now, given $x\in \bbZ^2$, we let $\cC^{\sm,+,1}_x$ be the set of positive small contours $\gamma_2$ for which~$x$ is the minimal point in $\bar \gamma_2$ for the lexicographical order.
We have
\begin{multline*}
\sum_{ \gamma_2 \in \cC^{\sm,+,1}_{\bar \gamma\setminus \bar \gG_1}}\hspace*{-2mm} \log \bigl(1+ e^{-\gb |\tilde\gamma_2|}
z^u_{1}(\gamma_2)\bigr)\le \sum_{x\in |\bar \gamma\setminus \bar \gG_1|} \sum_{\gamma_2 \in \cC^{\sm,+,1}_x} \hspace*{-2mm}\log \bigl(1+ e^{-\gb |\tilde\gamma_2|}
z^u_{1}(\gamma_2)\bigr)\\
=|\bar \gamma\setminus \bar \gG_1| \sum_{\gamma_2 \in \cC^{\sm,+,1}_{\bf 0}} \hspace*{-2mm}\log \bigl(1+ e^{-\gb |\tilde\gamma_2|}
z^u_{1}(\gamma_2)\bigr),
\end{multline*}
because the right hand side includes all the terms of the left hand side, plus a few extra contours that are not contained in $\bar \gamma\setminus \bar \gG_1$.
We observe that
\begin{equation}\label{night}
\begin{cases}
z^u_{1}(\gamma_2)=(1-J^2)^{-1} & \text{when } \bar \gamma_2=\{x\},\\
z^u_{1}(\gamma_2)=(1-J^3)^{-1}(1+J^2)(1-J^2)& \text{when } \bar \gamma_2=\{x,y\} \text{ with } x\sim y,\\
z^u_{1}(\gamma_2)\le e & \text{for other small contours}.
\end{cases}
\end{equation}
For the first two lines, the full computation is performed in \cite[Proof of Lem.\,3.2]{cf:part1} and the last one can be derived like \eqref{truxc}.
Noting that in $\cC^{\sm,+,1}_{\bf 0}$ there is one contour of length $4$ and two of length $6$, we obtained that
\begin{equation*}
\sum_{\gamma_2 \in \cC^{\sm,+,1}_x}\log \bigl(1+ e^{-\gb |\tilde\gamma_2|}
z^u_{1}(\gamma_2)\bigr)\le J^2+2J^3+ C J^4,
\end{equation*}
where the term $CJ^4$ includes the contribution of all contours of length $8$ or more.

For the case of boundary condition equal to one we observe similarly that
\begin{equation*}
Z^{u,\sm}_{1}[\gamma,\gG_1,1]\le
\prod_{\gamma_2 \in \cC^{\sm,1}_{\bar \gamma\setminus \bar \gG_1}}
\bigl(1+ e^{-\gb |\tilde\gamma_2|}
z^u_{1+\gep(\gamma_2)}(\gamma_2) \bigr),
\end{equation*}
where $\cC^{\sm,1}_{\bar \gamma\setminus \bar \gG_1}$ denote the set of $1$-small contours in $\bar \gamma\setminus \bar \gG_1$.
Similarly to \eqref{night}, one can check that (recall $h\le 2J^2$)
\begin{equation*}%\label{night2}
\begin{cases}
z^u_{1+\gep(\gamma_2)}(\gamma_2)\le 1+CJ^2 &\text{when } L(\tilde\gamma_2)\le 6,\\
z^u_{1}(\gamma_2)\le e &\text{for other small contours },
\end{cases}
\end{equation*}
and the results follows as for zero boundary condition case.

We obtain the lower bound results by restricting our sums to contours of length $4$ and $6$.
We set
\begin{equation}\label{csta}
\cC^{*,+}(\gamma,\gG_1):=\{ \gamma_2 \sep \gep(\gamma_2)=+, \ |\tilde \gamma_2|\le 6,\ \bar \gamma_2\subset \bar \gamma \setminus \bar \gG_1,\
d(\gamma_2, \{\gamma\}\cup \gG_1)>0 \},
\end{equation}
where $d(\gamma_2, \{\gamma\}\cap \gG_1)$ denotes the minimal distance between the geometric contours~$\tilde \gamma_2$ and those in the set $\{\gamma\}\cup \gG_1$. This condition ensures in particular compatibility. We~set
\[
v(\gamma_2):=e^{-\gb |\tilde\gamma_2|-h|\bar \gamma_2|}
z^u_{1}(\gamma_2).
\]
We obtain
\begin{multline}\label{drac}
Z^{u,\sm}_{0}[\gamma,\gG_1,1]
\ge e^{h|\bar \gamma \setminus \bar \gG_1|} \sum_{\{ \gG_2 \in \cK^{\sm,+,1}_{\ext}(\gamma) \sep \gG_2 \parallel \gG_1\}} \prod_{\gamma_2\in \gG_2 }
v(\gamma_2)\ind_{\{\gamma_2\in \cC^{*,+}(\gamma,\gG_1)\}}\\
=:e^{h|\bar \gamma \setminus \bar \gG_1|}
Z[v,\cC^{*,+}(\gamma,\gG_1)].
\end{multline}
We can apply cluster expansion results for the relation of external compatibility with the weights being given by $v$ (see Remark \ref{laremarc}).
We set for $\bC\subset \cC^{*,+}(\gamma,\gG_1)$
\begin{equation}\label{modpot}
v^{R}(\bC):=\sum_{\bB\subset \bC} (-1)^{|\bB|+|\bC|} \log Z[v,\bB],
\end{equation}
where $Z[v,\bB]$ is defined as $Z[v,\cC^{*,+}(\gamma,\gG_1)]$ in \eqref{drac} but with the indicator $\ind_{\{\gamma_2\in \bB\}}$.
We let $\cR^+(\gamma,\gG_1)$ denote the set of clusters associated with external compatibility in $\cC^{*,+}(\gamma,\gG_1)$
that is $\bC\subset \cC^{*,+}(\gamma,\gG_1)$ is in $\cR^+(\bar \gamma\setminus\bar \gG_1)$ if (recall \eqref{noncluster})
\begin{equation*}
\forall \bB \subset \bC, \quad \bB\nparallel \bC\setminus \bB,
\end{equation*}
where $\bB\nparallel \bC$ is the negation of $\bB\parallel \bC$.
After observing that $v^{R}(\bC)=0$ for $\bC\notin \cR^+(\bar \gamma\setminus\bar \gG_1)$
we obtain
\begin{multline}\label{dox}
\log Z[v,\cC^{*,+}\gamma,\gG_1)]= \sum_{\{ \bC\in \cR^+(\bar \gamma\setminus\bar \gG_1)\}} v^R(\bC)
\\
\ge \sum_{\{ \bC\in \cR^+(\gamma, \gG_1) \sep L(\bC)\le 6\}} v^R(\bC)
-\sum_{\{ \bC\in \cR^+(\gamma, \gG_1) \sep L(\bC)\ge 8\}} \left|v^R(\bC)\right|.
\end{multline}
Note that the clusters in the first sum in \eqref{dox} consist in only one contour. Using \eqref{night} again, and the assumption $h\le J^2$, the reader can check that for some appropriate constant $C$ we have
\begin{equation}\label{speel}
v^R(\{\gamma_2\})=\log\left(1+ v(\gamma_2)\right)\ge \begin{cases} J^2-C J^4, \quad \text{ if } \bar \gamma_2=\{x\},\\
J^3-C J^4, \quad \text{ if } \bar \gamma_2=\{x,y\} \text{ with } x\sim y.
\end{cases}
\end{equation}
We let the reader check that the number of contours of length $4$ and $6$ in $\cC^{*,+}(\gamma,\gG_1)$ satisfies
\begin{equation*}
\begin{split}
\#\{ \gamma_2\in \cC^{*,+}(\gamma,\gG_1) \sep |\tilde \gamma_2|=4\}&\ge |\bar \gamma\setminus\bar \gG_1|-(|\tilde \gamma|+L(\gG_1)),\\
\#\{ \gamma_2\in \cC^{*,+}(\gamma,\gG_1) \sep |\tilde \gamma_2|=6\}&\ge 2|\bar \gamma\setminus\bar \gG_1|-6(|\tilde \gamma|+L(\gG_1)),
\end{split}
\end{equation*}
the second term being caused by boundary effects.
Hence we have for $\gb$ sufficiently large
\begin{equation}\label{wick}
\sum_{\{ \bC\in \cR^+(\bar \gamma,\bar \gG_1) \sep L(\bC)\le 6\}} v^R(\bC)
\ge |\bar \gamma\setminus\bar \gG_1|(J^2+2J^3-CJ^4)- \frac{1}{2}(L(\gG_1)+|\bar \gamma|),
\end{equation}
where the second term is present to account for the fact that the number of contours is not exactly proportional to the volume $|\bar \gamma\setminus\bar \gG_1|$.

To control the second term in \eqref{dox}, we use Theorem \ref{superclust} for external compatibility with
\[
a(\gamma_2)=|\tilde \gamma_2| \text{ and } d(\gamma_2)=(\gb-5)|\tilde \gamma_2|.
\]
In that case, \eqref{condit} holds and one deduces from \eqref{dabound} that
\begin{equation*}
\sum_{\{ \bC\in \cR^+(\bar \gamma,\bar \gG_1) \sep L(\bC)\ge 8\}} \left|v^R(\bC)\right|\le CJ^4|\bar \gamma\setminus\bar \gG_1|,
\end{equation*}
which, together with \eqref{wick}, yields
\begin{equation*}
\log Z[v,\cC^{*,+}(\gamma,\gG_1)]\ge |\bar \gamma\setminus\bar \gG_1|(J^2+2J^3-C'J^4)- \frac{1}{2}(L(\gG_1)+|\bar \gamma|),
\end{equation*}
and leads to the conclusion (recall \eqref{drac}).

For the case with boundary condition equal to one, we have
\begin{equation}\label{sumikr}
Z^{u,\sm}_{1}[\gamma,\gG_1,1]:= \sum_{\{ \gG_2 \in \cK^{\sm}_{\ext}(\gamma,1) \sep \gG_1 \parallel \gG_2\}}
\prod_{\gamma_2\in \gG_2 }
e^{-\gb |\tilde\gamma_2|}
z^u_{1+\gep(\gamma_2)}(\gamma_2).
\end{equation}
The same argument as the one used at level $0$ makes it possible to obtain a lower bound.
Restric\-ting the sum to $\gamma_2\in \cC^{*}(\gamma,\gG_1)$, where
\[
\cC^{*}(\gamma,\gG_1):=\{ \gamma_2 \sep \ |\tilde \gamma_2|\le 6,\ \bar \gamma_2\subset \bar \gamma \setminus \bar \gG_1,\
d(\gamma_2, \{\gamma\}\cup \gG_1)>0 \},
\]
and setting
\[
v'(\gamma_2):=e^{-\gb |\tilde\gamma_2|}
z^u_{1+\gep(\gamma_2)}(\gamma_2),
\]
we have
\begin{equation*}
Z^{u,\sm}_{1}[\gamma,\gG_1,1]
\ge \sum_{\{ \gG_2 \in \cK^{\sm,+,1}_{\ext}(\gamma) \sep \gG_2 \parallel \gG_1\}} \prod_{\gamma_2\in \gG_2 }
v'(\gamma_2)\ind_{\{\gamma_2\in \cC^*(\gamma,\gG_1)\}}.
\end{equation*}
We can then check (using the fact that $h\le 2J^2$) that \eqref{speel} is satisfied for $(v')^R$ defined as in Equation \eqref{modpot}
and conclude in a similar manner (the coefficient are multiplied by two because $\cC^{*}(\gamma,\gG_1)$ contains twice as many contours).
\qed

\subsection{Proof of Lemma \ref{dalo}}

We start our proof with the lower bounds which are easier to achieve since they are a consequence of Jensen's inequality.
The upper bound results require a more delicate analysis and are treated afterward.

\subsubsection{Lower bounds, the proof of \eqref{dalo2}}

From Jensen's inequality, we have
\begin{equation*}
\log \bE^{n,\sm,n}_{\gamma,\gG_1}\bigl[e^{u|\phi^{-1}(\bbZ_-)|-\bar H(\phi^{-1}(\bbZ_-))}\bigr]
\ge \bE^{n,\sm,n}_{\gamma,\gG_1}\left[u|\phi^{-1}(\bbZ_-)|-\bar H(\phi^{-1}(\bbZ_-))\right],
\end{equation*}
and thus Equation \eqref{ckonveu}
can be deduced from the following inequality:
\begin{multline*}
\bE^{n,\sm,n}_{\gamma,\gG_1}\left[u|\phi^{-1}(\bbZ_-)|-\bar H(\phi^{-1}(\bbZ_-))\right]\\
\ge
|\bar \gamma \setminus \bar \gG_1|\Bigl(\frac{1}{2}u J^{2n}-40J^{3n+3}\Bigr)-3u(|\tilde \gamma|+L(\gG_1)),
\end{multline*}
provided that $u\le 1/12$.
We are going to bound terms separately and to show that%
\begin{equation}\label{lixiout}\begin{split}
\bE^{n,\sm,n}_{\gamma,\gG_1}\left[|\phi^{-1}(\bbZ_-)| \right]&\ge
\frac{1}{2}J^{2n}\bigl(|\bar \gamma\setminus\bar \gG_1|-3(|\tilde \gamma|+L(\gG_1))\bigr),\\
\bE^{n,\sm,n}_{\gamma,\gG_1}\left[\bar H(\phi^{-1}(\bbZ_-)) \right]&\le 40 J^{3n+3}|\bar \gamma\setminus\bar \gG_1|.
\end{split}\end{equation}
The first inequality can directly be deduced from \eqref{infin}: the equality is valid as soon as $x$ is not constrained by the boundary condition,
which might happen only if $x$ lies in $\Delta^-_{\gamma}$ or in $\Delta^+_{\gamma_1}$ for $\gamma_1\in \gG_1$,
hence the number of $x$ such that \eqref{infin} does not apply is proportional to $|\tilde \gamma|+L(\gG_1)$.

Concerning the second line in \eqref{lixiout}, we let $f_{i}(\phi)$ denote the number
of points which lie in a connected component of $\phi^{-1}(0)$ of size $i$ or larger:\vspace*{-3pt}
\begin{equation*}
f_{i}(\phi):=\{ x\in \bar\gamma \setminus \bar \gG_1 \sep \exists \gL\subset \bar\gamma \setminus \bar \gG_1,\ \gL \text{ connected };\ x\in \gL \forall z\in \gL, \ \phi(z)\le 0 \}.
\end{equation*}
Using Lemma \ref{zoomats} (we observe \eqref{sitmats} implies that $\bar H \{ x,y\} \le J^3$ for $\gb$ large), we obtain that\vspace*{-3pt}
\begin{equation}\label{lamine}
\bar H(\phi^{-1}(\bbZ_-))\le \frac{J^{3}}{2}\,f_2(\phi)+ 2 J^{2}f_{3}(\phi).
\end{equation}
From Proposition \ref{rouxrou}, we have
\begin{equation*}
\bE^{n,\sm,n}_{\gamma,\gG_1}[f_2(\phi)]
\le |\bar \gamma \setminus \bar \gG_1|4\max_{\{x,y \in \bar \gamma \setminus \bar \gG_1 \sep x\sim y\}}\bE^{n,\sm,n}_{\gamma,\gG_1}[ \max(\phi(x),\phi(y))\le 0]
\le 8 J^{3n}.
\end{equation*}
Similarly we obtain also using Proposition \ref{rouxrou}\vspace*{-3pt}\enlargethispage{.5\baselineskip}%
\begin{multline}\label{diatta}
\bE^{n,\sm,n}_{\gamma,\gG_1}[f_3(\phi)]\le 18
|\bar \gamma \setminus \bar\gG_1| \max_{\{x\sim y \sim z\}} \bE^{n,\sm,n}_{\gamma,\gG_1}[ \max(\phi(x),\phi(y),\phi(z))\le 0]\\[-3pt]
\le 36n J^{4n}\le 36 J^{3n+1}.
\end{multline}
The coefficient $18$ corresponds to the number of ways of choosing a connected set of size three which contains a given $x$.
Combining \eqref{lamine}--\eqref{diatta} we conclude that the second inequality in \eqref{lixiout} holds.
\qed

\subsubsection{Upper bounds-Proof of \eqref{dalo1} and \eqref{dalo3}}

The upper bound is a bit more delicate since the proof relies on some decorrelation property of the measure $\bP^{m,\sm,n}_{\gamma,\gG_1}$.
We are going to use the following technical statement, which ensures that the bounds from
Proposition \ref{rouxrou} remain valid after conditioning to the realization of the field outside a large ball.
For the rest of the proof one sets $r=3n^2\gb^2$.

\begin{lemma}\label{locall}
We have for any $m\ge 1$, for all of $\phi$, and any $x\in\bar \gamma \setminus \bar \gG_1$, we have\vspace*{-3pt}
\begin{equation*}%\label{leff}
\bP^{m,\sm,n}_{\gamma,\gG_1}[\phi(x)\le 0\ | \ \phi(z),\ |z-x|\ge r]\le 2 J^{2m}.
\end{equation*}
If moreover\vspace*{-3pt}
\begin{equation}\label{badcond}
\{x,y\}\cap \Bigl[\Delta^-_{\gamma}\cup \Bigl(\bigcup_{\gamma_1\in \gG_1} \Delta^+_{\gamma_1}\Bigr) \Bigr]=\emptyset,\vspace*{-3pt}
\end{equation}
then\vspace*{-3pt}
\begin{multline*}%\label{leffe}
\bP^{m,\sm,n}_{\gamma,\gG_1}\big[\min(\phi(x),\phi(y))\le 0; \ \forall w\in \partial \{x,y\}, \ \phi(w)>0 \\[-3pt]
| \ \phi(z),\ |z-x|\ge r\big]\ge \frac{1}{2} J^{3m}.
\end{multline*}
\end{lemma}
The condition \eqref{badcond} is present to ensure that both $\phi(x)$ and $\phi(y)$ are allowed to take negative value under
$\bP^{m,\sm,n}_{\gamma,\gG_1}$. The condition is sufficient but not always necessary since only positive contours might prevent to have $\phi(x)\le m$.

Let us now prove \eqref{dalo1}, which we choose to replace by a slightly more general statement. For $\gb$ sufficiently large one, for every $n\ge 1$,
$m\ge \max(n-1,1)$ and $v\le J^n$, we show that%
\begin{equation}\label{lecarlat}
\log \bE^{m,\sm,n}_{\gamma,\gG_1}\bigl[ e^{v |\phi^{-1}(\bbZ_-)|} \bigr] \le |\bar \gamma\setminus \bar \gG_1|3 v J^{2m}.
\end{equation}
Recall that $\delta_x:=\ind_{\{\phi_x\le 0\}}$.
In order to control the effect of correlation we choose to split $\bar \gamma \setminus \bar \gG_1$
according to the value of $\lfloor x/r \rfloor$ modulo $2$.
More precisely we split $\bar \gamma \setminus \bar \gG_1$ into squares of side-length $r$ and regroup these squares into four collections so that two squares in the same collection are never adjacent (see Figure \ref{fig:amma}).
For $i \in \lint 1, 4\rint$, $z\in\bbZ^2$ we set
\begin{equation}
\label{ztook}
\cB_i(z):=\{ x\in \bar \gamma \setminus \bar \gG_1 \sep \lfloor x_j/r \rfloor=\alpha_j(i)+2z_j,\ j=1,2\ \},
\end{equation}
where $\alpha_j(i)$ is the $j$-th digit in the dyadic development if $i-1$.

\begin{figure}[ht]
\begin{center}
\includegraphics[scale=1]{figures/amma}
\end{center}
\caption{\label{fig:amma}\smaller
A contour $\gamma$ is represented together with the collection $\Gamma_1$ of large external contours inside it (with hatched interior). We split the set $\bar \gamma \setminus \bar \gG_1$ into squares of sidelength $r$, and attribute to each square a label between $1$ and $4$ such that two squares with the same label are never adjacent. }
\end{figure}

We have from Hölder's inequality
\begin{equation}\label{zolder}
\bigl(\bE^{m,\sm,n}_{\gamma,\gG_1}\bigl[ e^{v |\phi^{-1}(\bbZ_-)|} \bigl]\bigr)^4\le \sum_{i=1}^4
\bE^{m,\sm,n}_{\gamma,\gG_1}\bigl[ e^{4v \sum_{z\in \bbZ^2}\sum_{x\in \cB_i(z)} \delta_x} \bigr].
\end{equation}
Using the fact that $e^x\le 1+e^K x$ for $x\in[0,K]$ and Lemma \ref{locall}, we obtain that for each $z\in \bbZ$
\begin{multline*}
\bE^{m,\sm,n}_{\gamma,\gG_1}\Bigl[ e^{4v \sum_{x\in \cB_i(z)} \delta_x} \ | \ \phi(y), y\in \bigcup_{z'\ne z}\cB_i(z') \Bigr]\\
\le 1+ 4ve^{4vr^2}\bE^{m,\sm,n}_{\gamma,\gG_1}\biggl[ \sum_{x\in \cB_i(z)} \delta_x \ | \ \phi(y), y\in \bigcup_{z'\ne z}\cB_i(z') \biggr]\\
\le 1+ 8ve^{4vr^2} J^{2m}|\cB_i(z)|.
\end{multline*}
Using this inequality iteratively and combining it with \eqref{zolder} we obtain that
\begin{equation*}
\log \bE^{m,\sm,n}_{\gamma,\gG_1}\bigl[ e^{v |\phi^{-1}(\bbZ_-)|} \bigr]\le 2ve^{4vr^2} J^{2m}|\bar \gamma\setminus \bar \gG_1|.
\end{equation*}
The inequality \eqref{lecarlat} follows provided $e^{4vr^2}\le 3/2$, which is the case under our assumption provided $\gb$ is sufficiently large.

Let us now turn to the more delicate case \eqref{dalo3}.
Using Cauchy-Schwartz inequality we have
\begin{multline*}
\bE^{n-1,\sm}_{\gamma,\gG_1}\bigl[ e^{u |\phi^{-1}(\bbZ_-)|-\bar H(\phi^{-1}(\bbZ_-))} \bigr]^2\\
\le
\bE^{n-1,\sm}_{\gamma,\gG_1}\bigl[ e^{2u |\phi^{-1}(\bbZ_-)|} \bigr] \bE^{n-1,\sm}_{\gamma,\gG_1}\bigl[ e^{-2\bar H(\phi^{-1}(\bbZ_-))} \bigr].
\end{multline*}
To evaluate the first term we can rely on \eqref{lecarlat} and conclude that
\begin{equation*}
\log \bE^{n-1,\sm}_{\gamma,\gG_1}\bigl[ e^{2u |\phi^{-1}(\bbZ_-)|} \bigr] \le 6 u J^{2(n-1)}.
\end{equation*}
The final step is to prove that
\begin{equation*}
\log \bE^{n-1,\sm}_{\gamma,\gG_1}v[ e^{-\bar H(\phi^{-1}(\bbZ_-))} \bigr] \le -
\frac{1}{4} J^{3n}\left(|\bar \gamma\setminus\bar \gG_1|-10\left(|\tilde \gamma|+L(\gG_1)\right)\right).
\end{equation*}

As $\bar H$ is positive, this yields the same upper bound for the expectation of $e^{-2\bar H(\phi^{-1}(\bbZ_-))}$.
We set for $\gL\subset\bbZ^2$ finite
\begin{equation*}
G^{k}(\gL):= \log \bE^+_{\gL,\gb} \bigl[e^{-\bar H(\phi^{-1}[k,\infty))}\bigr].
\end{equation*}
Conditioning to the level set at level $n-1$ we obtain
\begin{multline*}%\label{rouge}
\bE^{n-1,\sm,n}_{\gamma,\gG_1}\bigl[ e^{-\bar H(\phi^{-1}(\bbZ_-))} \bigr]
\\=\bE^{n-1,\sm,n}_{\gamma,\gG_1}\bigl[ e^{-G^{n-2}(\phi^{-1}(-\infty,n-2])} \bigr]=
\bE^{1,\sm,n}_{\gamma,\gG_1}\bigl[ e^{G^{n-2}(\phi^{-1}(\bbZ_-))} \bigr].
\end{multline*}
We refer to \cite[Eq.\,(6.7)]{cf:part1} for details concerning the first equality, the second one simply uses vertical translation invariance.
Note that $G^k(\gL)=\sum_I G^k(\gL_i)$, as this is the case for $\bar H$ (recall \eqref{maxdecomp}), where $(\gL_i)_{i\in I}$ is the decomposition
into maximal connected components of $\gL$.
Note also that we have, according to \cite[Lem.\,6.2]{cf:part1}, for $\gb$ sufficiently large, for any pair of neighbors $\{x,y\}$ in $\bbZ^2$,
\begin{equation*}
G^{n-2}(\{x,y\})=\log \Bigl(1- \frac{J^3-J^4}{1-J^4}J^{3(n-2)} \Bigr)\le -\frac{1}{2}J^{3n-3}.
\end{equation*}
Hence, setting
\begin{equation*}
\delta_2(x):=\ind_{\{\phi(x)\le 0 \text{ and the connected component of $x$ in $\phi^{-1}(\bbZ_-)$ has size two}\}},
\end{equation*}
and ignoring the contribution to $G^{n-2}$ of connected components of larger size (for singletons, note that $G^{k}\{x\}=0$), we have
\begin{equation*}
G^{n-2}(\phi^{-1}(\bbZ_-))\le -\frac{1}{8}J^{3n-3} \sum_{x\in \bar \gamma \setminus \bar \gG_1}\delta_2(x).
\end{equation*}
As we are going to use Lemma \ref{locall}, we are going to consider the sum restricted to
\begin{equation*}
\gL(\gamma,\gL_1):=\Bigl\{ x\in \bar \gamma \setminus \bar \gG_1 \sep d\Bigl(x, \Delta^-_{\gamma}\cup \Big(\bigcup_{\gamma_1\in \gG_1} \Delta^+_{\gamma_1}\Big)\Bigr)\ge 2\Bigr\}.
\end{equation*}
Note that
\begin{equation*}
\gL(\gamma,\gL_1)\ge |\bar \gamma \setminus \bar \gG_1|- 10 \left(|\tilde \gamma| + L(\gG_1) \right).
\end{equation*}
To conclude we thus need to prove that
\begin{equation}\label{zoyer}
\log \bE^{1,\sm,n}_{\gamma,\gG_1}\bigl[ e^{-\Psfrac{1}{2}J^{3n-3} \sum_{x\in \gL}\delta_2(x)} \bigr]\le \frac{1}{4}J^{3n}|\gL|.
\end{equation}

We now choose to proceed as in the proof of \eqref{lecarlat} to deal with the correlation
between the variables $\delta_2(x)$.
We set with the same notation as in \eqref{ztook} for $i\in \lint 1, 4 \rint$, $z\in \bbZ^2$,
\begin{equation*}%\label{ztookit}
\cB_i(z):=\{ x\in \gL \sep \lfloor x_j/r \rfloor=\alpha_j(i)+2z_j,\ j=1,2\ \}.
\end{equation*}
Note that from Lemma \ref{locall} (observe that, to have $\delta_2(x)=1$, there are four ways to chose a neighbor of $x$ for which the field is negative) we have
\begin{equation}\label{wiou}
\bE^{1,\sm,n}_{\gamma,\gG_1}\biggl[ \sum_{x\in \cB_i(z)} \delta_2(x) \ | \ \phi(y), \ y \in \bigcup_{z'\ne z}\cB_i(z') \biggr]\ge 2 J^3.
\end{equation}
Like for \eqref{zolder} we have from Hölder's inequalities\vspace*{-3pt}
\begin{multline}\label{miax}
\Bigl(\bE^{1,\sm,n}_{\gamma,\gG_1}\bigl[ e^{ -\Psfrac{1}{4}J^{3n-3} \sum_{x\in \gL}\delta_2(x)} \bigr]\Bigr)^4\\[-5pt]
\le \prod_{i=1}^4 \bE^{1,\sm,n}_{\gamma,\gG_1}\bigl[ e^{-J^{3n-3} \sum_{z\in \bbZ}\sum_{x\in \cB_i(z)} \delta_2(x)} \bigr].
\end{multline}
Using $e^{-x}\le 1-e^{-K}x$ for $x\in[0,K]$ and \eqref{wiou}, we obtain that for any $z$ and $i$\vspace*{-3pt}
\begin{multline*}
\bE^{1,\sm,n}_{\gamma,\gG_1}\Bigl[ e^{-J^{3n-3} \sum_{x\in \cB_i(z)} \delta_2(x)} \ | \ \phi(y), \ y \in \bigcup_{z'\ne z}\cB_i(z') \Bigr]\\
\le 1-e^{-J^{3n-3}r^2}J^{3n-3}\bE^{1,\sm,n}_{\gamma,\gG_1}\biggl[ \sum_{x\in \cB_i(z)} \delta_2(x) \ | \ \phi(y), \ y \in \bigcup_{z'\ne z}\cB_i(z') \biggr]
\\ \le 1- J^{3n}|\cB_i(z)|.
\end{multline*}
Using this inequality to evaluate each term factor in the l.h.s of \eqref{miax}\vspace*{-3pt}
\begin{equation*}
\prod_{i=1}^4 \bE^{1,\sm,n}_{\gamma,\gG_1}\left[ e^{-J^{3n-3} \sum_{z\in \bbZ}\sum_{x\in \cB_i(z)} \delta_2(x)} \right]\le
\left(1-J^{3n}\right)^{|\gL|}.
\end{equation*}
which in view of \eqref{miax} concludes the proof of \eqref{zoyer}. \qed

\begin{proof}[Proof of Lemma \ref{locall}]
We are going to show that the inequalities are valid when a stronger conditioning is considered.
We set\vspace*{-3pt}
\[
\cC[r,x]:=\left\{ \gamma \in \cC[\gamma,\gG_1,m] \sep \bar \gamma\cap B(x,r/2)\ne \emptyset \right\}
\]
where $B(z,\ell)$
the Euclidean ball of center $z$ radius $\ell$. We let $\hat \cC[r,x]$ be the corresponding set of cylinders.

Conditioned to $\hat\Upsilon_{m}(\phi)\cap (\hat \cC[r,x])^{\cc}$ (the set of cylinders of $\phi$ which do not intersect $B(x,r/2)$), which, due to smallness of contours,
is a stronger conditioning than
$(\phi(z),\ |z-x|\ge r)$,
the distribution of
$\phi$ restricted to $B(x,r/2)$
is of the type $\bP_{\bL,\gL,\gb}$ described at the beginning of Section \ref{peaksec}, where
\[
\bL:=\{\gamma_3 \in \cC[r,x] \sep \forall \gamma_2\in \hat\Upsilon_{m}(\phi)\cap (\hat \cC[r,x])^{\cc},\ \gamma_3 \mid \gamma_2 \}.
\]
The results follows then by applying Proposition \ref{rouxrou}.
\end{proof}

\section{First consequences of the cluster expansion convergence}\label{dalast}

In this section, we exploit contour stability to obtain
the convergence of Gibbs measures $\bP^{n,h}_{\gL,\gb}$ for appropriate values of $h$ in Section \ref{moots}.
We also exploit the decay of correlation to prove the regularity of the free energy in Section \ref{laprov}.
Finally we close the proof of Theorem \ref{main} in Section \ref{lapriv} where we prove Proposition \ref{convirj}.

\subsection{Convergence of the Gibbs measure}\label{moots}

In view of the results of Section \ref{clusexp}, the stability of the contour implies the convergence of some measure when the size of the box grows.
We need however some work to convert this result into a convergence result for the SOS measure.
The proof relies on the following observation: the distribution of the set of external contour $\Upsilon^{\ext}_n(\phi)$
under $\bP^{n,h}_{\gL,\gb}$ is the same that the distribution of external contour for the measure $\bbP^{w^{h}_{n}}_{\gL}$
(this is simply a by-product of the proof of Proposition \ref{trax}, recall Equation \eqref{tapouz}).
The second observation is that, conditioned to the set of external contours, the distribution of the field inside each contour is independent from the
rest and with an explicit distribution.

\begin{proposition}\label{conviark}
Let $f$ and $g$ be local functions $\gO_{\infty}\to [0,1]$ with respective supports $A$ and $B$.
For $n\ge 0$ and $h\in [h^*_{n+1},h^*_n]$, we have for any simply connected $\gL,\gL'\subset \bbZ^2$ which satisfies $A\subset \gL\cap \gL'$
and $d(A,\gL)\le d(A,\gL')$, there exists a positive constant
$C$ such that, for all $\gb$ sufficiently large,
\begin{equation}\label{clox}
\bigl| \bP^{n,h}_{\gL,\gb}[f(\phi) ]-\bP^{n,h}_{\gL',\gb}[f(\phi)]\bigr|\le C|A|e^{\Psfrac{\gb}{100}d(A,\gL)}.
\end{equation}
As a consequence the sequence of measure $\bP^{n,h}_{\gL,\gb}$ converges, when $\gL$ exhausts $\bbZ^2$,
to a limit $\bP^{n,h}_{\gb}$ which satisfies
\begin{equation}\label{cloxf}
\bigl| \bE^{n,h}_{\gL,\gb}[f(\phi) ]-\bE^{n,h}_{\gb}[f(\phi)]\bigr|\le C|A|e^{\Psfrac{\gb}{100}d(A,\gL)}.
\end{equation}
Moreover for we have for all $\gL$ and $h\in [h^*_{n+1},h^*_n]$,
\begin{equation}\label{clax}
\bigl| \bE^{n,h}_{\gL,\gb}[ f(\phi)g(\phi) ]- \bE^{n,h}_{\gL,\gb}[ f(\phi)] \bE^{n,h}_{\gL,\gb}[g(\phi) ] \bigr| \le C|A|e^{\Psfrac{\gb}{100}d(A,B)}.
\end{equation}
\end{proposition}

A consequence of \eqref{cloxf} is that the measure $\bE^{n,h}_{\gb}$ is translation invariant.
Note that \eqref{clox} ensures that the convergence of $\bP^{n,h}_{\gL,\gb}[f(\phi)]$ (which is a sequence of continuous functions in $h$)
is uniform in $h\in [h^*_{n+1},h^*_n]$ and thus it implies that $\bP^{n,h}_{\gb}[f(\phi)]$ is continuous on that interval.

\begin{proof}
Similarly to Proposition \ref{propinfi}, the proof of \eqref{cloxf} amounts to evaluating total variation distances for the distribution of $\phi\restrict_A$.
We are going to show that in fact the result can be obtained as a consequence of Proposition \ref{propinfi} and Theorem \ref{converteo}.

Firstly, we notice that conditionally on the set of external contour $\Upsilon^{\ext}_n(\phi)$,
the conditional distribution of the field in $\bar \gamma$ for $\gamma \in \Upsilon^{\ext}_n(\phi)$ is independent of that in $\gL\setminus \bar \gamma$
and given by $\bP_{\gamma}^{n,h}$ (recall \eqref{defpnu}).
Hence the distribution of $\phi\restrict_{A}$ is completely determined by the set of external contours which intersect $A$ (recall the definition of~$\cC'_A$ below
\eqref{locallyfinite}). For this reason one has
\begin{multline*}%\label{dafa}
\| \bP^{n,h}_{\gL,\gb}(\phi\restrict_A\in\csbullet)- \bP^{n,h}_{\gL',\gb}(\phi\restrict_A\in\csbullet)\|_{TV}\\\le
\| \bP^{n,h}_{\gL,\gb}(\Upsilon^{\ext}_n(\phi) \cap \cC'_{A}\in\csbullet)- \bP^{n,h}_{\gL',\gb}(\Upsilon^{\ext}_n(\phi) \cap \cC'_{A}\in\csbullet)\|_{TV}\\
\le \bbP^{w^{h}_{n}}_{\gL}(\Upsilon \cap \cC'_{A}\in\csbullet)- \bbP^{w^{h}_{n}}_{\gL'}(\Upsilon \cap \cC'_{A}\in\csbullet)\|_{TV}.
\end{multline*}
The second inequality is due to the fact that the $\Upsilon^{\ext}_n$ is the same as that of $\Upsilon^{\ext}$
the set of external contours associated with $\Upsilon$ under $\bbP^{w^{h}_{n}}_{\gL}$.
We can conclude using Proposition~\ref{propinfi} and Theorem \ref{converteo}.
For the proof of \eqref{clax}, we condition the expectation according to the realization of the set of external contours.
Let us start by observing that when $\Upsilon^{\ext}_n(\phi)\cap \cC'_A \cap \cC'_B=\emptyset$, we have
\begin{equation}\label{troum}
\bE^{n,h}_{\gL,\gb}[f(\phi)g(\phi) \ | \ \Upsilon^{\ext}_n(\phi)]= \bE^{n,h}_{\gL,\gb}[f(\phi) \, | \, \Upsilon^{\ext}_n(\phi)] \bE^{n,h}_{\gL,\gb}[g(\phi) \, | \,
\Upsilon^{\ext}_n(\phi)].
\end{equation}
Now, applying Proposition \ref{propinfi} for the distribution of $\Upsilon^{\ext}_n(\phi)$ (which is distributed like the set of external contours under
$\bbP^{w^{h}_{n}}_{\gL}$) and for the functions $\tilde f$ and $\tilde g$ which are the conditional expectations of $f$ and $g$ given the external set of contours,
we obtain that
\begin{multline*}
\left| \bE^{n,h}_{\gL,\gb}\left[ \bE^{n,h}_{\gL,\gb}[f(\phi) \, | \, \Upsilon^{\ext}_n(\phi)] \bE^{n,h}_{\gL,\gb}[g(\phi) \ | \
\Upsilon^{\ext}_n(\phi)]\right] - \bE^{n,h}_{\gL,\gb}[f(\phi)] \bE^{n,h}_{\gL,\gb}[g(\phi)] \right|
\\ \le C|A|e^{-(\gb/100) d(A,B)}.
\end{multline*}
And we conclude by noticing that from \eqref{troum} we have
\begin{multline*}
\left| \bE^{n,h}_{\gL,\gb}[f(\phi)g(\phi)]- \bE^{n,h}_{\gL,\gb}\left[ \bE^{n,h}_{\gL,\gb}[f(\phi) \, | \, \Upsilon^{\ext}_n(\phi)] \bE^{n,h}_{\gL,\gb}[g(\phi) \ | \
\Upsilon^{\ext}_n(\phi)]\right]\right|\\ \le
\bP^{n,h}_{\gL,\gb}[\Upsilon^{\ext}_n(\phi)\cap \cC'_A \cap \cC'_B=\emptyset]\le |A|e^{-(\gb/2) d(A,B)}.
\end{multline*}
where in the last inequality we used \eqref{pignouf}.
\end{proof}

\Subsection{Exponential decay for Ursell function and proof of Proposition \ref{infinitz}}\label{laprov}

The fact that exponential decay of correlation implies differentiability relies
on a well established theory exposed \eg in \cite{cf:vDK} in the case of the Ising model.
The argument displayed in \cite{cf:vDK} adapts verbatim to our problem.
For the sake of completeness, we provide here the main steps.

In order to prove Proposition \ref{infinitz} we are going to show that
for every $h\in [h^*_{n+1},h^*_n]$, $\gL\subset \bbZ^2$, $k\ge 1$, there exists a constant $C_k$ such that
\begin{equation}\label{bootk}
\left|\partial^k_h\bigl(\log \cZ^{n,h}_{\gL,\gb}\bigr)\right|\le C_k |\gL|
\end{equation}
Then using Arzelà-Ascoli's theorem,
we can deduce from \eqref{bootk} that $\tf(\gb,h)$ is infinitely differentiable on $(h^*_{n+1},h^*_n)$ and that
\begin{equation}\label{koomits}
\lim_{N\to \infty}\frac{1}{N^2}\,\partial^k_h\bigl(\log \cZ^{n,h}_{N,\gb}\bigr)= \partial^k_h \tf(\gb,h).
\end{equation}
To prove \eqref{bootk}, setting $\delta_x:=\ind_{\{\phi(x)=0\}}$, we use the fact that
\begin{equation}\label{mood}
\partial^k_h\bigl(\log \cZ^{n,h}_{\gL,\gb}\bigr)= \sum_{(x_1,\dots,x_k)\in \gL^k} \bE^{n,h}_{\gL,\gb}\left[ \delta_{x_1} ; \dots ; \delta_{x_k} \right],
\end{equation}
where $\bE^{n,h}_{\gL,\gb}\left[ \delta_{x_1} ; \dots ; \delta_{x_k} \right]$ denotes the $k$ point truncated correlation function (or~Ursell function)
defined by
\begin{equation*}
\bE^{n,h}_{\gL,\gb}\left[ \delta_{x_1} ; \dots ; \delta_{x_k} \right]:=
\sum_{\cP} (-1)^{|\cP|+1}(|\cP|+1)! \prod_{P\in \cP}\bE^{n,h}_{\gL,\gb}\biggl[\prod_{i \in P} \delta_{x_i}\biggr],
\end{equation*}
where the sum in $\cP$ ranges over the set of partitions of $\lint 0,k \rint$, and $|\cP|$ denotes the cardinal of the partition.

The identity \eqref{mood} can be obtained by induction on $k$ and is a particular case of \cite[Eq.\,(1.9)]{cf:vDK}.
We can conclude by using some decay estimates for these correlation functions.

\begin{proposition}\label{mickey}
For any $k$ there exists a positive constant $m_k$ such that for every $h\in (h^*_{n+1},h^*_n)$ and $x_1,\dots,x_k\in \gL$
\begin{equation*}
\bE^{n,h}_{\gL,\gb}\left[ \delta_{x_1} ;\dots ;\delta_{x_k} \right]\le C_k e^{-m_k d(x_1,\dots,x_k)},
\end{equation*}
where
\[
d(x_1,\dots,x_k)\text{ is the smallest cardinal of a connected $\bbZ^2$ set containing $\{x_1,\dots,x_k\}$.}
\]
\end{proposition}
To deduce that \eqref{bootk} holds, one only needs to check that we have
\begin{equation*}
\sum_{x_2,\dots,x_k\in \bbZ^2} e^{-m_k d(x_1,\dots,x_k)}<\infty.
\end{equation*}
This is an obvious consequence of
\begin{equation*}
d(x_1,\dots,x_k)\ge \max_{i\in \lint 2,k\rint} |x_1-x_i|\ge \frac{1}{k-1} \sum_{i=2}^k |x_1-x_i|.
\end{equation*}
\begin{proof}[Proof of Proposition \ref{mickey}]
We can follow line by line the proof of \cite[Lem.\,3.1]{cf:vDK} where we replace \cite[Ass.\,(3.2)]{cf:vDK} by \eqref{clax}.
\end{proof}

\subsection{Proof of Proposition \ref{convirj}}\label{lapriv}

Recalling Equations \eqref{koomits}--\eqref{mood} for $k=1$, we only need to prove that
for $h\in [h^*_{n+1},h^*_n]$ we have
\begin{equation*}
\frac{1}{2}\, J^{2n}\le \frac{1}{N}\bE^{n,h}_{N,\gb}\biggl[ \sum_{x\in \lint 1, N\rint}\delta_x\biggr]\le 2J^{2n}.
\end{equation*}
We can in fact prove the inequality holds for $\bP^{n,h}_{\gL,\gb}\left[ \phi(x)=0 \right]$, uniformly in $\gL$ and $x$.
Let us start with the special case $n=0$, for which the upper bound is trivial.
For the lower bound we apply \eqref{godstim} to $w^h_n$ and $A=\{x\}$,
we obtain that
\begin{multline*}
\bP^{n,h}_{\gL,\gb}[\phi(x)=0]\ge \bP^{n,h}_{\gL,\gb}[ \text{No contour in } \Upsilon_n(\phi) \text{ encloses } x]\\
=\bbP^{w^h_n}_{\gL}\left[ \Upsilon \cap \cC'_{\{x\}}=\emptyset \right]
=\exp\biggl(-\sum_{\bC\in \cQ(\gL,x,\emptyset)} (w^h_n)^T(\bC)\biggr),
\end{multline*}
where
\begin{equation*}
\cQ(\gL,\{x\},\emptyset):=\{ \bC \in \cQ(\cC_{\gL}) \sep \exists \gamma \in \bC, x\in \bar \gamma\}.
\end{equation*}
The estimate \eqref{dabound2} is then sufficient to conclude that
\begin{equation*}
\sum_{\bC\in \cQ(\gL,x,\emptyset)} |(w^h_n)^T(\bC)|\le 1/2,
\end{equation*}
which is sufficient for our purpose.

When $n\ge 1$, we let $\bg_x$ be the unique external contour enclosing $x$ whenever it exists.
As a vertex not enclose in any contour is by definition at level $n$ we have
\begin{multline*}
\bE^{n,u}_{N,\gb}[ \phi(x)=0]= \sum_{\gamma\in \cC_N: x\in \bar \gamma} \bP^{n,u}_{N,\gb}\left[ \phi(x)=0\, ; \, \bg_x=\gamma \right]\\[-5pt]
=\sum_{\gamma\in \cC_N: x\in \bar \gamma} \bP^{n,u}_{N,\gb}\left[ \bg_x=\gamma \right]\bP^{n,u}_{\gamma}\left[ \phi(x)=0\right],
\end{multline*}
where in the last equality, we used that the conditional distribution of $\phi$ restricted to~$\bar \gamma$ is given by
$\bP^{n,u}_{\gamma}$ defined in \eqref{defpnu}. We are going to show that most of the contribution to the sum
is given by the negative contour $\gamma^{x,-}$ which satisfies $\bar \gamma^{x,-}=\{x\}$ (there are two contours of length $4$ such that $x\in \bar \gamma$ but
$\gamma^{x,-}$ is the only one which contributes to the sum).
More precisely we are going to show that for $\gb$ sufficiently large and $u$ satisfying the assumption we have
\begin{equation}\label{precizion}
\left| \bP^{n,u}_{N,\gb}\left[ \phi(x)=0\, ; \, \bg_x=\gamma^{x,-} \right]- e^{-4\gb n} \right|\le C e^{-4\gb n}\left(u+ e^{-4\gb}\right).
\end{equation}
and
\begin{equation}\label{pro}
\sum_{\gamma\in \cC_N: x\in \bar \gamma,\ |\tilde \gamma|\ge 6} \bP^{n,u}_{N,\gb}\left[ \phi(x)=0\, ; \, \bg_x=\gamma \right] \le C e^{-(4n+2)\gb},
\end{equation}
which yields a sharper result than required.
For \eqref{precizion}
using \eqref{godstim} for $A=\{x\}$ again, we obtain that
\begin{multline*}
\bP^{n,u}_{N,\gb}\left[ \bg_x=\gamma^{x,-}\right]= \bbP^{w^h_n}_{\gL}[\Upsilon \cap \cC'_{\{x\}}=\gamma^{x,-}]\\
=w^h_n(\gamma^{x,-})\exp\biggl(-\sum_{\bC\in \cQ(\gL,x,\{\gamma^{x,-}\})} (w^h_n)^T(\bC)\biggr),
\end{multline*}
where
\begin{equation*}
\cQ(\gL,\{x\},\{\gamma^{x,-}\}):=\{ \bC \in \cQ(\cC_{\gL}\setminus \gamma^{x,-}) \sep \exists \gamma \in \bC,\ x\in \bar \gamma \text{ or } \bC \perp \gamma^{x,-} \}.
\end{equation*}
We have\vspace*{-3pt}
\begin{equation*}
w^h_n(\gamma^{x,-})=e^{-4\gb}\,\frac{1+(e^{u-1})e^{-4\gb n}}{1+(e^{u-1})e^{-4\gb{n-1}}},
\end{equation*}
and as a consequence of \eqref{dabound2},\vspace*{-3pt}
\begin{equation*}
\biggl| \sum_{\bC\in \cQ(\gL,x,\{\gamma^{x,-}\})} (w^h_n)^T(\bC) \biggr| \le C e^{-4\gb},
\end{equation*}
which is sufficient to yield \eqref{precizion}.

For \eqref{pro}
first, let us notice that we can discard the contribution of contours longer than $100 n$, because using a union bound and \eqref{pignouf} we have
\begin{equation*}
\bP^{n,u}_{N,\gb}\left[ \Diam(\bg_x)\ge 100 n \right]\le e^{-(4n+2)\gb}.
\end{equation*}
For smaller contours,
we need an estimate $\bP^{n,u}_{\gamma}\left[ \phi(x)=0\right]$.
Recalling the definition of $\bP^{n}_{\gamma}$ below \eqref{tambem} we have\vspace*{-3pt}\enlargethispage{.5\baselineskip}%
\begin{multline}\label{zipa}
\bP^{n,u}_{\gamma}\left[ \phi(x)=0\right]=
\frac{\bE^{n+\gep(\gamma)}_{\gamma}\bigl[\delta_x e^{u|\phi^{-1}(0)|-H(\phi^{-1}(\bbZ_-))} \bigr]}{\bE^{n+\gep(\gamma)}_{\gamma}\bigl[e^{u|\phi^{-1}(0)|-H(\phi^{-1}(\bbZ_-))} \bigr]}
\\
\le \frac{e^{|\bar \gamma| u}\bE^{n+\gep(\gamma)}_{\gamma}\left[\delta_x \right]}{\bP^{n+\gep(\gamma)}_{\gamma}[\forall x\in \bar \gamma, \phi(x)\ge 1]}
\le 4e^{- 4\gb(n+\gep(\gamma))}.
\end{multline}
where in the last inequality, we used Proposition \ref{rouxrou} to estimate both the numerator and the denominator.
Now, using \eqref{pignouf}, one has\vspace*{-3pt}
\begin{equation}\label{zipi}
\bP^{n,u}_{N,\gb}\left[ \bg_x=\gamma \right]\le w^h_n(\gamma) \le e^{-(\gb-1)|\tilde \gamma|},
\end{equation}
where the last inequality is a consequence of Proposition \ref{lesmall}.
Combining \eqref{zipa} and \eqref{zipi} as well as a standard bound for the number of contours of a given length enclosing $x$ we find that\vspace*{-3pt}
\begin{multline*}
\sum_{\gamma\in \cC_N: \ x\in \bar \gamma,\ |\tilde \gamma|\ge 6,\ \Diam(\gamma)\le 100 n}
\bP^{n,u}_{N,\gb}\left[ \bg_x=\gamma \right]\bP^{n,u}_{\gamma}\left[ \phi(x)=0\right]\\\le 4e^{- 4\gb(n-1)}
\sum_{\gamma\in \cC_N: x\in \bar \gamma,\ |\tilde \gamma|\ge 6}e^{-(\gb-1)|\tilde \gamma|}\le C e^{- 4\gb(n+2)},
\end{multline*}
which is sufficient to conclude the proof of \eqref{pro}.
\qed

\section{Properties of Gibbs measure: the proof for Theorem \ref{Gibbs}}\label{daverylast}

In this final section, we prove the remaining unproved statements from Theorem~\ref{Gibbs}.
First in Section \ref{start}, we prove our statement concerning $\star$-connectivity of the level sets.
Then in Section \ref{noon} we prove that there exists no Gibbs states for $h\le h_w(\gb)$.
In Section \ref{cf}, we identify the contact fraction for each Gibbs states which has been obtained in Proposition \ref{conviark}.
This yields in particular \eqref{variousgibbs}.
In Section \ref{zerob} we identify the minimal Gibbs states, which is the one obtained by taking the limit of zero boundary condition.
In Section \ref{unik} we prove uniqueness of Gibbs states at differentiability points, and in Section \ref{sanduba}, we prove that at
$\bP^{n,h^*_n}_{\gb}$ and $\bP^{n-1,h^*_n}_{\gb}$ are respectively the maximal and minimal Gibbs states corresponding to $h^*_n$ proving \eqref{ssanduba}.

\subsection{Percolative properties of level sets}\label{start}

Let us check that for all $h\in [h^*_n,h^*_{n+1}]$, the random field $\phi$ percolates at level $n$ under $\bP^{n,h}_{\gb}$.
The external contour lines of $\phi$ under $\bP^{n,h}_{\gb}$ can be obtained by considering the set of external contour of a sample of
$\bbP^{\go^h_n}$. In particular, this implies that almost surely there are no infinite contour lines.

As each maximal connected components of $\phi^{-1}[n+1,\infty)$ \resp $\phi^{-1}(-\infty,n-1]$ is enclosed in a positive contour \resp negative contour,
this implies that they are all finite. We can even prove using a union bound argument and Theorem \ref{converteo} that the diameter of the
largest such component in a box of side-length $N$ is of order $\log N$).

Proving the existence of an infinite component for $\phi^{-1}(n)$ is more tricky, as some points which are not enclosed by any contour can belong to finite clusters of
$\phi^{-1}(n)$. Now, our result will hold if we can prove that
\begin{equation}\label{thatsallfolks}
\bbR^2\setminus \bigcup_{\gamma\in \Upsilon^{\ext}_n(\phi)} \bar \gamma \quad \text{ has a unique unbounded connected component},
\end{equation}
where in the equation above, with a small abuse of notation, $\bar \gamma$ denotes the closed subset of $\bbR^2$ enclosed by $\gamma$.

Our idea is to compare the set $\bigcup_{\gamma\in \Upsilon^{\ext}_n(\phi)} \bar \gamma$ with the occupied set of a Poisson Boolean percolation process \cite{cf:ATT, cf:Hall, cf:G}.
We know that the set of external contours under $\bP^{n,h}_{\gb}$ can be obtained as a subset of a sample of $\bbP^{\go^h_n}$
which itself is dominated (\eg by \cite[Lem.\,4.4]{cf:part1}) by
by a random collection of contours $\chi$ where each contour $\gamma$ is present independently with probability
$\spfrac{\go^h_n(\gamma)}{1+\go^h_n(\gamma)}\le e^{(1-\gb) |\tilde \gamma|}$ (the inequality being a consequence of Theorem \ref{converteo}).

We let the reader check that, for $\gb$ sufficiently small, $\bigcup_{\gamma\in \chi} \bar \gamma$ is stochastically dominated by a continuum percolation process,
where obstacles are balls whose centers are distributed according to a Poisson point process with intensity $\gl(\gb)=e^{-\gb}$ and
whose radius are IID with standard exponential distribution (see \eg \cite{cf:ATT, cf:Hall, cf:G} for a more formal definition).

It has been proved that for $\gl$ sufficiently small the vacant set for such a Boolean percolation process percolates, and that the occupied set is
only composed of bounded connected components (a much stronger result is displayed in \cite[Th.\,1]{cf:ATT} with optimal assumptions,
but the statement we need can also be extracted from earlier works
\eg \cite{cf:G, cf:Hall}). This proves \eqref{thatsallfolks} and concludes our reasoning.
\qed

\subsection{Absence of Gibbs state for $h\le h_w(\gb)$ }\label{noon}

As by the DLR relation a Gibbs states can always be obtained as a limit of finite volume measures with random boundary conditions,
we know that the limit obtained with zero boundary condition is, if finite, the minimal Gibbs state.
We are going to prove the following result which implies divergence of the distribution of $\phi$.

\begin{proposition}
For $\gb$ sufficiently large, for $h\le h_w(\gb)$ we have for any $x\in \bbZ^2$ and any $K>0$
\begin{equation*}
\lim_{\gL\to \bbZ^2} \bP^{h}_{\gL,\gb}(\phi(x)\le K)=0.
\end{equation*}
\end{proposition}

Note that by monotonicity in $h$ (Corollary \ref{FKalt}), it is sufficient to check the statement for $h=h_w(\gb)$.

\begin{proof}
Now, using the DLR relation for the neighborhood of $x$ and the definition of the measure, one obtains that for any $k\ge 0$ one has
\begin{equation*}
\bP^{h_w(\gb)}_{N,\gb}(\phi(x)=k+1 \ | \ \phi(y),\ y\sim x)\le e^{4\gb} \bP^{h_w(\gb)}_{N,\gb}(\phi(x)=k+1 \ | \ \phi(y),\ y\sim x).
\end{equation*}
This readily implies that for an explicit constant $C(\gb,K)$ one has
\begin{equation*}
\bP^{h_w(\gb)}_{\gL,\gb}(\phi(x)\le K)\le C(\gb,K) \bP^{h_w(\gb)}_{\gL,\gb}(\phi(x)=0).
\end{equation*}
To conclude we just need to show that
\begin{equation}\label{suit}
\lim_{\gL \to \bbZ^2}\bP^{h_w(\gb)}_{\gL,\gb}(\phi(x)=0)=0.
\end{equation}
As a consequence of Theorem \ref{oldmain} and \eqref{difencial},
we have
\begin{equation*}
\lim_{N\to \infty} \frac{1}{N^2}\sum_{x\in \lint 1,N\rint^2} \bP^{h_w(\gb)}_{N,\gb}(\phi(x)=0)=\partial_h \tf(\gb,h_w(\gb))=0.
\end{equation*}
As, by monotonicity (Corollary \ref{FKalt}), each term in the sum is larger than the limit one wants to compute, we obtain \eqref{suit}.
\end{proof}

\subsection{Identifying the contact fraction: the proof of \eqref{variousgibbs}}\label{cf}

Let us prove that for any $h\in [h^*_{n+1},h^*_n]$ we have
\begin{equation}\label{wootz}
\bP^{n,h}_{\gb}[\phi(x)=0]=\partial_h \tf(\gb,h),
\end{equation}
where the derivative as to be understood as the derivative on the right for $h=h^*_{n+1}$ and on the left for
$h=h^*_{n}$.
We already know as a consequence of \eqref{koomits} that for $h\in(h^*_{n+1},h^*_n)$ we have
\begin{equation}\label{zeone}
\lim_{N\to \infty} \frac{1}{N^2}\sum_{x\in \lint 1,N \rint^2 }\bP^{n,h}_{N,\gb}[\phi(x)=0]=\partial_h \tf(\gb,h).
\end{equation}

Note that the statement can be extended to the boundary of the interval $h\in\{h^*_{n+1},h^*_n\}$ using that
the second derivative $\partial^2_h \log Z^{n,h}_{N,\gb}$ is uniformly bounded on the interval $[h^*_{n+1},h^*_n]$ (recall \eqref{koomits}).

A consequence of the exponential decay of correlation \eqref{cloxf} is that
\begin{equation}\label{zetwo}
\biggl|\sum_{x\in \lint 0,N \rint^2 }\bP^{n,h}_{\gb,N}[\phi(x)=0]- N^2\bP^{n,h}_{\gb}[\phi(x)=0]\biggr|\le C N,
\end{equation}
and thus \eqref{wootz} is deduced from the combination of \eqref{zeone} and \eqref{zetwo}.

\subsection{Limit with zero boundary condition}\label{zerob}

So far we have only shown the existence of translation invariant Gibbs measure and non-uniqueness at the phase transition points.
To conclude we need an argument to show that they are the only one.

A first step is to show convergence of the measure when having zero boundary condition to a translation invariant limit which has the right contact fraction.
The proof uses essentially the same idea as those to prove a similar result for wetting of the harmonic crystal
\cite[\S 5]{cf:GL}.

\begin{proposition}
For any $h>h_w(\gb)$
the sequence of measure $\bP^{h,0}_{\gL,\gb}$ converge to an infinite volume limit which we call $\tilde \bP^{h}_{\gb}$.
We have for every $x\in 0$
\begin{equation}\label{zioup}
\tilde \bP^{h}_{\gb}[\phi(x)=0]=\partial^+_h \tf(\gb,h).
\end{equation}
\end{proposition}

\begin{proof}
A first observation is that $\bP^{h,0}_{\gL,\gb}$ restricted to $A\subset \gL$ increases stochastically when $\gL$ increases.
This is a consequence of the DLR property (recall \eqref{DLR}) and Corollary \ref{FKalt}: for $\gL'\supset \gL$ the restriction
$\bP^{h,0}_{\gL',\gb}$ to $\gL$ corresponds to $\bP^{h,\Psi}_{\gL,\gb}$ with a random boundary condition $\Psi\ge 0$ which thus dominates $\bP^{h,0}_{\gL,\gb}$.

The sequence of measure is tight because, from Corollary \ref{FKalt}, $\bP^{h,0}_{\gL,\gb}$ is dominated by $\bP^{n,h}_{\gL,\gb}$ which converges.
Hence we obtain the existence of the limit.

To prove the statement about the contact fraction, let us set $\eta:= \tilde \bP^{h}_{\gb}[\phi({\bf 0})=0]$.
Note that, by monotonicity for any $\gL$ and $x\in \gL$, we have
\begin{equation*}
\bP^{h,0}_{\gL,\gb}[\phi(x)=0] \ge \eta,
\end{equation*}
and hence \eqref{difencial} implies that, when the derivative exists,
\begin{equation*}
\partial_h\tf(\gb,h)\ge \eta.
\end{equation*}
Given $\gep>0$, using the definition of $\eta$ we can choose $K=K_\gep$ sufficiently large which satisfies
\[
\bP^{h}_{\lint -K,K\rint^2,\gb}[\phi(x)=0]\le \eta+\gep.
\]
By monotonicity we also have
\[
\bP^{h,0}_{N,\gb}[\phi(x)=0]\le \eta+\gep
\]
for all $x$ such that $d(x,\partial \gL_N)\ge K+1$.
Hence we have
\begin{equation*}
\sum_{x\in \gL_N} \tilde \bP^{h}_{N,\gb}[\phi(x)=0]\le 4 (K+1) N+ N^2(\eta+\gep).
\end{equation*}
Using \eqref{difencial} again, we conclude that $\partial_h\tf(\gb,h)\le \eta+\gep$ and thus obtain that \eqref{zioup} holds at all differentiability points.
To conclude, we remark that $\tilde \bP^{h}_{\gb}[\phi({\bf 0})=0]$ being the infimum (in $N$) of continuous non-decreasing function, it has to be right-continuous 
at every point, from which we deduce that the results also hold where $\tf$ is not differentiable.
\end{proof}

The next step is to show that the limit found above coincides with $\bP^{n,h}_{\gb}$.

\begin{proposition}\label{dabot}
For any $h\in [h^*_{n+1},h^*_n)$, we have
\begin{equation}
\tilde \bP^{h}_{\gb}= \bP^{n,h}_{\gb}.
\end{equation}
\end{proposition}

\begin{proof}
Because of the stochastic ordering induced by the boundary condition before taking the limit (recall Corollary \ref{FKalt}),
$\bP^{n,h}_{\gb}$ stochastically dominates $\tilde \bP^{h}_{\gb}$.
We can thus consider a coupling $\bQ$ of the two measures (the first marginal being distributed like $\bP^{n,h}_{\gb}$ and the second one like $\tilde \bP^{h}_{\gb}$) such that
\begin{equation*}
\forall y\in \bbZ^2, \phi_1(y)\ge \phi_2(y), \quad \bQ\ \text{a.s.}
\end{equation*}
We are going to show that we have in fact $\phi_1(y)= \phi_2(y)$ with probability one, which shows that the two marginal probabilities are equal.
Now, by \eqref{wootz} and \eqref{zioup} the two measures have the same contact fractions.
Given $x\in \bbZ^2$, using the DLR equation for the neighborhood of $x$ we have
\begin{equation}\label{superdlr}
\bQ[\phi_i(x)=0]=
\bQ\biggl[ \frac{e^{h-\gb\sum_{y\sim x}\phi_i(y)}}{e^{h-\gb\sum_{y\sim x}\phi_i(y)}+ \sum_{k\ge 1} e^{-\gb\sum_{y\sim x}|\phi_i(y)-k| }} \biggr].
\end{equation}
The reader can now check that the function on the right-hand side (let us denote it by $\chi(\phi)$) is strictly decreasing in $\phi(y)$ for all $y$ neighboring $x$. In particular our monotone coupling implies that $\chi(\phi_2)-\chi(\phi_1)\ge 0$.
As we have
\begin{equation}
0= \bQ[\phi_2(x)=0]-\bQ[\phi_1(x)=0]
=\bQ\left[ \chi(\phi_2)-\chi(\phi_1)\right],
\end{equation}
which implies $\chi(\phi_2)=\chi(\phi_1)$ almost surely, and strict monotonicity implies in turn that
$ \phi_2(y)=\phi_1(y)$ with probability one for all neighbors of $x$. As $x$ is arbitrary, the statement is valid for all $y \in \bbZ^2$, and this concludes the proof.
\end{proof}

\subsection{Uniqueness of Gibbs states}\label{unik}

The key point is to prove the following result.
\begin{lemma}\label{labon}
If $h$ is a differentiability point of $\tf(\gb,h)$, then there is only one translation invariant, finite mean Gibbs state.
\end{lemma}

\begin{proof}
Let $\nu^h$ be a translation invariant, finite mean Gibbs state.
Let $\hat \phi$ be a boundary condition sampled according to $\nu^h$ (but for practical reason we write $\hat \nu^h$ for the distribution).
Then the DLR equation implies that the law of $\phi\restrict_{\gL}$ under $\bP^{h,\hat \phi}_{\gL,\gb}$ corresponds to the restriction of $\nu^h$.
In particular, this implies that for any increasing local function
\begin{equation*}
\nu^h(f(\phi))=\hat \nu^h\bigl(\bE^{h,\hat \phi}_{\gL,\gb}f(\phi)\bigr)\ge \bE^{h,0}_{\gL,\gb}[f(\phi)].
\end{equation*}
Passing to the limit we conclude that $\nu^h$ dominates the limit obtained with zero boundary condition which has been identified in Lemma \ref{dabot}.
Now, from translation invariance we have
\begin{equation}\label{contrex}
\nu^h(\phi(x)=0)=\frac{1}{N^2}\,\hat \nu^{h}\bE^{h,\hat \phi}_{N,\gb}\biggl[\sum_{x\in \gL_N} \phi(x)=0\biggr].
\end{equation}
We are now going to prove that for every $h'\in \bbR$ we have
\begin{equation}\label{evian}
\lim_{N\to \infty}\frac{1}{N^2}\,\hat \nu^{h}\bigl(\log \cZ^{h',\hat \phi}_{N,\gb}\bigr)=\tf(\gb,h').
\end{equation}
This follows from the fact that convergence holds with zero boundary condition
and that the effect boundary conditions can be controlled via the following observation
\begin{align*}
|\log \cZ^{h,\hat \phi}_{N,\gb}-\log \cZ^{h}_{N,\gb}|&\le \gb\max_{\phi}|\cH^{\hat \phi}_{N}(\phi)- \cH_{N}(\phi)|,\notag
\\
|\cH^{\hat \phi}_{N}(\phi)- \cH_{N}(\phi)|&\le \sum_{x\in \partial \gL_N} \hat \phi(x).
\end{align*}
We can conclude using the fact that from our assumptions $\hat \nu^h\bigl(\sum_{x\in \partial \gL_N} \hat \phi(x)\bigr)\le CN$.

We now observe that \eqref{evian} and convexity imply that the right hand side of \eqref{contrex} converges to $\partial_h\tf(\gb,h)$,
and thus that $\hat \nu^h$ and $\tilde \bP^h_{\gb}$ have the same contact fraction.
Using the same trick as in the proof of Lemma \ref{dabot} (recall \eqref{superdlr}) we prove that the two measures coincide.
\end{proof}

\subsection{Stochastic sandwich at angular points: the proof of \eqref{ssanduba}}\label{sanduba}

Assume now that $h=h^*_n$.
While the derivative of the free energy does not exist, the proof of Lemma~\ref{labon} still implies that
\[
\nu^h(\phi(x)=0)\in [\partial^-_h \tf(\gb,h),\partial^+_h\tf(\gb,h)],
\]
and that $\nu^h$ stochastically
dominates the Gibbs states obtained in the limit with $0$ boundary condition which we know to be $\bP^{n-1,h}_\gb$.

A last thing to prove is that $\nu^h$ is dominated by $\bP^{n,h}_{\gb}$.
Consider $\hat \phi$ being distributed according to the measure $\nu^h$ (we write $\hat \nu^h$)
and consider the finite volume measure corresponding to boundary condition $n\vee\hat \phi=\max(n,\hat \phi)$.

If we allow $\phi(x)=\infty$ and make the set $\bbZ_+\cup\{\infty\}$ compact, then any sequence of measures is tight and thus admits a limit point.
We consider $\nu'$ a limit point of the following sequence of probability on $\gO_N$:
\begin{equation*}
\nu'_N(\cdot):=\frac{1}{(2N^2-2N+1)^2}\sum_{y\in \lint N-N^2,N^2-N\rint^2}\hat \nu^h \bP^{h,\hat \phi\vee n}_{\lint -N^2,N^2\rint^2,\gb}\bigl[ \left(\phi(x+y)\right)_{x\in \lint N,N\rint} \in \csbullet \bigr].
\end{equation*}
Note that, by construction, $\nu'$ is translation invariant and dominates both $\nu^h$ and $\bP^{n,h}_{\gb}$.
Moreover $\nu'$ also satisfies the following version of the DLR equation (recall~\eqref{DLR}):
For~every finite subset $\gL$ of $\bbZ^2$ and for every local bounded continuous \hbox{$g: (\bbZ_+\cup\{\infty\})^{\bbZ^2} \to \bbR$} --- in
particular, the limit of $g(\phi)$, when $\min_{x\in \gL}\phi(x)\to \infty$, exists and we call it $g(\infty)$ --- then we have $\nu'$ almost surely
\begin{multline*}%\label{superdlr2}
\nu'\left[ g(\phi) \ | \ \phi(x)=\psi(x), \forall x\in \partial \gL\right]\\=
\begin{cases}
\frac{1}{\cZ^{\psi,h}_{\gL,\gb}}\sum_{\phi\in \gO^+_{\gL}}e^{-\beta\cH^{\psi}_{\gL}(\phi)+h|\phi^{-1}(0)|},
&\text{if } \psi(y)<\infty \text{ for all } y \in \partial \gL,\\
g(\infty) &\text{if } \psi(y)=\infty \text{ for some } y \in \partial \gL.
\end{cases}
\end{multline*}
The statement is valid for every measure in the sequence for $N$ sufficiently large and passes to the limit by continuity
(see \cite[Eq.\,(5.3)]{cf:GL} for a similar argument).

Similarly to \eqref{evian} we have
\begin{equation*}
\lim_{N\to \infty}\frac{1}{N^2}\hat \nu^{h}\bigl(\log \cZ^{h',n\vee \hat \phi}_{N,\gb}\bigr)\le \tf(\gb,h').
\end{equation*}
and thus \eqref{difencial} implies readily that
$\nu'(\phi({\bf 0})=0)\ge \partial^-_h\tf(\gb,h)$
and hence as from stochastic comparison $\nu'(\phi({\bf 0})=0)\le \bP^{n,h}_{\gb}(\phi({\bf 0})=0)$, we conclude from \eqref{variousgibbs} that
\[
\nu'(\phi({\bf 0})=0)\ge \partial^-_h\tf(\gb,h).
\]

To conclude, we need to prove that $\nu'=\bP^{n,h}_{\gb}$.
By stochastic domination, there exists a coupling $(\phi_1,\phi_2)$ of the measures $\nu'$ and $\bP^{n,h}_{\gb}$,
such that almost surely $\phi_1(x)\ge \phi_2(x)$ for all $x$. As we have $\{\phi_1(x)=0\}\subset \{\phi_2(x)=0\}$,
the fact that both events have equal probability implies that almost surely
\[
\phi^{-1}_1\{0\}=\phi^{-1}_2\{0\}.
\]

In particular, we have
\begin{equation*}
\nu'(\exists x\in \bbZ^2, \phi(x)=0)= \bP^{n,h}_{\gb}(\exists x\in \bbZ^2,\; \phi(x)=0)=1.
\end{equation*}
Then, replicating the argument in \cite[\S 5]{cf:GL}, we deduce that
\begin{equation*}
\nu'(\forall x\in \bbZ^2,\; \phi(x)<\infty)=1.
\end{equation*}
To conclude, we use the DLR relation in a neighborhood of $x$ for $\nu'$ and $\bP^{n,h}_{\gb}$ like in Equation \eqref{superdlr} to prove that the two measures coincide.

\appendix

\section*{Appendix. Cluster expansion estimates}
\refstepcounter{section}\label{apclus}

\subsection{Proof of Lemma \ref{finalfrontier}}\label{pfinal}

From the expression \eqref{fraconv} and translation invariance we have
\begin{equation*}
|\bar \gamma|\tf(w)= \sum_{x\in \bar \gamma }\sum_{\{\bC\in \cQ \sep x-(1/2,1/2) \in \bC\}} \frac{1}{|\bC|}w^T(|\bC|).
\end{equation*}
In this sum, the coefficient of $w^T(|\bC|)$ for a cluster in $\cQ(\cC_{\gga})$ is $1$ while the other clusters have a coefficient between zero and one.
Thus we have
\begin{equation*}
\left|Z[\gamma,w]- |\bar \gamma|\tf(w)| \right|\le \sum_{\{\bC\in \cQ\setminus \cQ(\cC_{\gamma}) \sep \bC\cap [\bar \gamma-(1/2,1/2)]\ne \emptyset \}} w^T(|\bC|).
\end{equation*}
Note now that, for all clusters in the right hand side, there is at least one site in $\tilde \gamma$ that belongs to $\bC$. Thus, using translation invariance,
given a fixed $x^*_0\in (\bbZ^2)^*$ we have
\begin{equation*}
\left|Z[\gamma,w]- |\bar \gamma^*|\tf(w)| \right|\le |\tilde \gamma|\sum_{\bC\in \cQ \sep x_0^* \in \bC} w^T(|\bC|)\le \frac{1}{4}\,|\tilde \gamma|,
\end{equation*}
where the last inequality is valid for $\gb$ sufficiently large as a consequence of \eqref{dabound2}.
\qed

\subsection{Proof of Proposition \ref{propinfi}}\label{ppropinf}

The second line of \eqref{convergenz} can be deduced from the first one by considering a sequence $\bL'$ that exhausts $\cC$.
The first result corresponds to evaluating the total variation distance between the respective distributions of $\Upsilon\cap \cC'_A$ under
$\bbP^w_{\bL}$ and $\bbP^w$, which is equal to
\begin{equation*}
\sum_{\gG \in \cK(\cC'_A)} \left(\bP^w_{\bL}\left(\Upsilon \cap \cC'_A=\gG \right)- \bP^w_{\bL'}\left(\Upsilon \cap \cC'_A=\gG\right) \right)_+.
\end{equation*}
We set for this proof $d:=d(A,\bL^{\cc})$.
We consider first the contribution to the above sum of $\gG$ which contains a contour $\gamma$ of large diameter, that is, such that
$|\Diam(\gamma)|\ge d/3$.
Simply using the fact that by Peierls argument (see \eg \cite[Lem.\,4.4]{cf:part1}) we have for every $\gamma$
\begin{equation}\label{pignouf}
\bbP^{w}_{\bL}[\gamma \in \Upsilon]\le \frac{w(\gamma)}{1+w(\gamma)}.
\end{equation}
We deduce from the assumption \eqref{usualchoice}, summing over all possible such contours that for $\gb$ sufficiently large
\begin{equation*}
\bbP^w_{\bL}\bigl(\exists \gamma \in \Upsilon, \ \bar\gamma \cap A\ne \emptyset, |\tilde \gamma|\ge \sfrac{2d}{3}\bigr)\le
|A| e^{-\gb \sfrac{d}{2}}.
\end{equation*}
Now, from the definition of $d(A,\bL^{\cc})$, if $\gG$ does not contain a contour of large diameter then all contours in it belongs to $\bL$ and $\bL'$.
In that case we have, from \eqref{godstim}
\begin{multline*}
\left(\bP^w_{\bL}\left(\Upsilon \cap \cC'_A=\gG \right)- \bP^w_{\bL'}\left(\Upsilon \cap \cC'_A=\gG\right) \right)_+
\\= \bP^w_{\bL}\left(\Upsilon \cap \cC'_A=\gG \right) \Bigl[ e^{ \left(\sum_{\bC'\in \cQ(\bL',A,\gG)} w^T(\bC')-
\sum_{\bC \in \cQ(\bL,A,\gG)} w^T(\bC)\right)_+ }-1 \Bigr].
\end{multline*}
And we can conclude provided that we can show that the difference in the exponential is small.
We notice that under the assumption that $\Diam(\gamma)\le d/3$ for $\gamma\in \gG$ we have
\begin{equation*}
\cQ(\bL,A,\gG)\,\triangle\,\cQ(\bL',A,\gG)\subset \cQ_1.
\end{equation*}
where $\triangle$ stands for the symmetric difference between sets and
\begin{equation*}
\cQ_1:=\{ \bC \sep \exists \gamma_1,\gamma_2 \in \bC, \min_{x\in A, y\in \bar \gamma_1} |x-y|\le d/3, \max_{x\in A, y^*\in \gamma_2}|x-y^*|\ge d \},
\end{equation*}
where the existence of $\gamma_1$ is justified by the fact that $\bC$ must contain a contour that either is connected to $\gG$ (and the inequality follows from
the assumption that contours in $\gG$ have diameter smaller than $d/3$), or belong to $\cC'_A$
and that of $\gamma_2$ by the fact that
it must contain a contour in $\bL\,\triangle\,\bL'$.
This implies in any case that the diameter of $\bC$ is larger than $d/2$ so that $L(\bC)\ge d$ and one can conclude using \eqref{dabound2} that
\begin{equation*}
\sum_{\bC\in \cQ_1} |w^T(\bC)|\le C d^2 |A|e^{-\gb d/2}\le C|A|e^{-\gb d/4},
\end{equation*}
where $|A|d^2$ gives a bound for the number of points at distance $d/3$ or less from $A$.

Let us now move to the proof of \eqref{decayze}.
This is simply a bound on the total variation distance between the distribution $\Upsilon \cap \cC'_{A\cup B}$ and what we obtain by considering the product distribution of $\Upsilon \cap \cC'_A$ and $\Upsilon \cap \cC'_B$. We must prove
\begin{multline*}
\sum_{\gG \in \cK(\cC'_{A\cup B})}\left(\bbP^w_{\bL}(\cC'_{A\cup B} \cap \Upsilon= \gG)- \bbP^w_{\bL}(\cC'_A\cap \Upsilon= \cC'_A\cap \gG)
\bbP^w_{\bL}(\cC'_B\cap \Upsilon=\cC'_B\cap \gG) \right)_+\\[-5pt]
\le |A| e^{-c\gb d(A,B)}.
\end{multline*}
A first step is to discard the possibility of having a contour that reaches the neighborhood of both $A$ and $B$. More precisely, using \eqref{pignouf}
one can check that if \eqref{usualchoice}
holds we have
\begin{equation}\label{badseeds}\begin{split}
\bbP^w_{\bL}\left(\exists \gamma \in \Upsilon \cap \cC'_{A}, \ \Diam(\gamma)\ge d/4 \right)&\le |A| e^{\gb d(A,B)/8},\\
\bbP^w_{\bL}\Bigl(\exists \gamma \in \Upsilon \cap \cC'_{B}, \ \min_{x\in A, y^*\in \gamma}|x-y^*|\le d/4 \Bigr) &\le |A| e^{-c\gb d(A,B)/8}.
\end{split}
\end{equation}
Now, if $\gG$ is such that $ \gG_1:= \cC'_A\cap \gG$ and $\gG_2:=\cC'_B\cap \gG$ are disjoint,
using \eqref{godstim}, and observing that $\cQ(\bL,A,\gG_1)\cup \cQ(\bL,B,\gG_2)= \cQ(\bL,A\cup B,\gG)$
\begin{multline*}
\left(\bbP^w_{\bL}(\cC'_{A\cup B} \cap \Upsilon= \gG)- \bbP^w_{\bL}(\cC'_A\cap \Upsilon= \gG_1)
\bbP^w_{\bL}(\cC'_B\cap \Upsilon=\gG_2) \right)_+
\\= \bbP^w_{\bL}(\cC'_{A\cup B} \cap \Upsilon= \gG)\Bigl[e^{\left(\sum_{\bC\in \cQ(\bL,A,\gG_1)\cap\cQ(\bL,B,\gG_2) } w^T(\bC)\right)_+}-1\Bigr].
\end{multline*}
Now, if $\gG$ does not include any contour of the type considered in \eqref{badseeds},
\begin{equation*}
\cQ(\bL,A,\gG_1)\cap\cQ(\bL,B,\gG_1) \subset \cQ_2,
\end{equation*}
where, using the notation $d:=d(A,B)$, we define
\begin{equation*}
\cQ_2:=\{ \bC\in \cQ \sep \exists \gamma_1, \gamma_2 \in \bC, \min_{x\in A, y\in \bar \gamma_1} |x-y|\le d/4, \ \max_{x\in A, y\in \bar \gamma_2}\ge 3d/4\}.
\end{equation*}
The existence of $\gamma_1$ is justified by the fact that some contour in $\bC$ must either be connected with $\gG_1$ or intersect $A$ and that of $\gamma_2$
by the fact that some contour in $\bC$ must either be connected with $\gG_2$ or intersect $B$.
This implies in particular that the diameter of $\bC$ is larger then $d/2$, so that one must have $L(\bC)\ge d$ and we can also conclude using \eqref{dabound2} that
\[
\sum_{\bC\in \cQ_2} |w^T(\bC)|\le C d^2 |A|e^{-\gb d/2}\le C|A|e^{-\gb d}.\eqno\qed
\]

\subsection{Proof of \eqref{infin} and \eqref{infde}}\label{lazt}

The proof is very similar to the one of the lower bound for Proposition \ref{convirj} displayed in Section \ref{lapriv}.
Let us only treat the case of \eqref{infde} as \eqref{infin} is very similar.

We only need to consider the contribution of configuration for which the only contour in $\cC'_{\{x,y\}}$ is given by the positive contour $\gamma^+_{x,y}$
of length $6$ which encloses $x$ and $y$.
Using the definition of $ \bP_{\bL,\gL,\gb}$ we obtain that
\begin{equation*}
\bP_{\bL,\gL,\gb}[\min(\phi(x),\phi(y))\ge n]\ge \bbP^{w_{\gb}}_{\bL}\left[\Upsilon\cap \cC'_{\{x,y\}}=\gamma^+_{x,y}\right]e^{-6\gb(n-1)}.
\end{equation*}
Applying \eqref{godstim} for the contour weight \eqref{specialweight}
\begin{equation*}
\bbP^{w_{\gb}}_{\bL}\left[\Upsilon\cap \cC'_{\{x,y\}}=\gamma^+_{x,y}\right]= \frac{1}{e^{6\gb}-1}
\exp\biggl(-\sum_{\bC\in \cQ(\bL,\{x,y\},\{\gamma^+_{x,y}\})} w^T(\bC)\biggr).
\end{equation*}
The sum in the exponential can be seen to be small as a consequence of \eqref{dabound2}, and this leads to the conclusion by choosing $\gb$ sufficiently large.


\bibliographystyle{jepplain+eid}
\bibliography{lacoin}
\end{document}
